% Worked solutions for Listas/lista-topic3-asymmetric-information.md (19 exercises: 15 from the books, 16-19 original).
% Compile: pdflatex twice; check: pdffonts (Type 1 only, uni = yes).
% Every number is asserted in lista-topic3-asymmetric-information-solutions_check.py.
\documentclass[11pt,a4paper]{article}
\usepackage[utf8]{inputenc}
\usepackage[T1]{fontenc}
\usepackage{lmodern}   % Type 1 vector fonts (replaces PK bitmaps)
\usepackage{cmap}      % ToUnicode CMap -> copyable text and accents
\usepackage[brazilian,english]{babel}
\usepackage{amsmath,amssymb,amsfonts}
\usepackage[a4paper,margin=2.2cm]{geometry}
\usepackage{booktabs}
\usepackage{enumitem}
\usepackage{fancyhdr}
\usepackage{xcolor}
\usepackage{microtype}

\pagestyle{fancy}\fancyhf{}
\lhead{\footnotesize Finan\c{c}as II --- Corporativas (MPE Insper)}
\rhead{\footnotesize Topic 3 list --- Worked solutions}
\cfoot{\footnotesize\thepage}
\renewcommand{\headrulewidth}{0.3pt}
\setlength{\headheight}{14pt}
\setlist{topsep=2pt,itemsep=1pt}

\newcommand{\exsec}[1]{\section*{\color{blue!55!black}#1}\addcontentsline{toc}{section}{#1}}
\newcommand{\notation}{\par\noindent\textbf{Notation.}}
\newcommand{\step}[2]{\par\medskip\noindent\textbf{Step #1 --- #2.}~}
\newcommand{\intuition}{\par\medskip\noindent\textbf{Intuition.}~}
\newcommand{\bookans}{\par\medskip\noindent\textbf{Book answer.}~}
\newcommand{\refbib}{\par\smallskip\noindent\textbf{Ref:}~}
\newcommand{\assume}{\par\smallskip\noindent\textbf{Assumption.}~}
\newcommand{\pt}{\foreignlanguage{brazilian}}
\newcommand{\asked}{\par\medskip\noindent\textbf{What is asked.}~}
\newcommand{\approach}{\par\smallskip\noindent\textbf{Approach.}~}
\newcommand{\answer}[1]{\par\smallskip\noindent\fbox{\textbf{Answer #1.}}~}
\newcommand{\pitfalls}{\par\medskip\noindent\textbf{Common mistakes.}}

\usepackage[hyphens]{url}
%% STATEMENT-PREAMBLE BEGIN (tools/statements_into_solutions.py)
\usepackage{longtable,array,calc}
\providecommand{\tightlist}{\setlength{\itemsep}{0pt}\setlength{\parskip}{0pt}}
\providecommand{\href}[2]{#2}
\providecommand{\real}[1]{#1}
\makeatletter\@ifundefined{c@none}{\newcounter{none}}{}\makeatother  % pandoc's \LTcaptype{none}
\newcommand{\stmtbegin}{\par\medskip\noindent\textbf{Statement.}\par\nopagebreak\begingroup\small\setlength{\tabcolsep}{3pt}\sloppy}
\newcommand{\stmtend}{\par\endgroup\nopagebreak\noindent\rule{\linewidth}{0.4pt}\par\medskip}
\DeclareUnicodeCharacter{2192}{\ensuremath{\rightarrow}}
\DeclareUnicodeCharacter{2212}{\ensuremath{-}}
\DeclareUnicodeCharacter{2022}{\textbullet{}}
\DeclareUnicodeCharacter{2248}{\ensuremath{\approx}}
\DeclareUnicodeCharacter{2032}{\ensuremath{^{\prime}}}
\DeclareUnicodeCharacter{2081}{\ensuremath{_{1}}}
\DeclareUnicodeCharacter{03B8}{\ensuremath{\theta}}
\DeclareUnicodeCharacter{03BB}{\ensuremath{\lambda}}
\DeclareUnicodeCharacter{03C8}{\ensuremath{\psi}}
\DeclareUnicodeCharacter{03B1}{\ensuremath{\alpha}}
\DeclareUnicodeCharacter{25A0}{\rule{0.7em}{0.7em}}
%% STATEMENT-PREAMBLE END
\begin{document}
\begin{center}{\Large\bfseries Exercise list --- Topic 3: Asymmetric information}\\[4pt]{\large Worked solutions}\end{center}

\noindent Statements are in \url{Listas/lista-topic3-asymmetric-information.md} (and its PDF); each exercise is cited by number and printed page and its full text is reproduced right before its solution. Exercises 16--19 are original exercises (no course book has one on the Stiglitz--Weiss model or on collateral as a signal), written with the user's agreement on 2026-11-12 from the lecture slides and Tirole \S6.3; they have no book answer. Short answers: \url{lista-topic3-asymmetric-information-solutions-short.pdf}. Editions are the local copies: Tirole, \emph{The Theory of Corporate Finance} (2005); Amaro de Matos, \emph{Theoretical Foundations of Corporate Finance} (2001); Berk = Berk, DeMarzo \& Harford, \emph{Fundamentals of Corporate Finance}, 2nd ed. (2011); Grinblatt--Titman, \emph{Financial Markets and Corporate Strategy}, 2nd ed. (2002); Ross = Ross, Westerfield, Jordan \& Lamb, \emph{Fundamentos de Administra\c{c}\~ao Financeira}, 9th ed. (2013).

\tableofcontents

% =============================================================================
% Topic 3 list (asymmetric information), exercises 1--8. Body fragment only.
% =============================================================================
\exsec{1. Winner's curse in IPO allocations --- Ross Cap.~15, Quest\~ao 3, p.~537}
%% BEGIN-STATEMENT 1
\stmtbegin
\textbf{Source:} Ross, Cap. 15, Questão 3 (Questões e problemas, Básico), p.~537 (PDF 561) · Answer
in book: no

\begin{enumerate}
\def\labelenumi{\arabic{enumi}.}
\setcounter{enumi}{2}
\tightlist
\item
  Subpreço de IPO {[}OA3{]} A Madeira S/A e a Metal S/A anunciaram IPOs a \$40 por ação. Uma delas
  está subavaliada em \$7 e a outra está superavaliada em \$5, mas você não tem como saber qual é
  qual. Você pretende comprar 1.000 ações de cada emissão. Se uma emissão estiver subprecificada,
  ela será racionada e você receberá apenas metade do seu pedido. Se você pudesse obter 1.000 ações
  da Madeira e 1.000 ações da Metal, qual seria seu ganho? Qual ganho você realmente espera? Qual
  princípio você ilustrou?
\end{enumerate}
\stmtend
%% END-STATEMENT

\asked (a) Your profit if both orders of 1{,}000 shares were filled in full. (b) The profit you should actually expect, given that only the underpriced issue is rationed (you get half of it) and you cannot tell the two issues apart. (c) The name of the principle the gap between (a) and (b) illustrates.

\approach This is the uninformed investor's problem in Rock's IPO model (lecture block 3.10). The allocation you receive is correlated with the quality of the issue: good issues are oversubscribed and rationed, bad ones are filled in full. So the expected gain must weight each issue by the fraction you are actually allocated, not by the fraction you ordered. Setting that expected gain to zero gives the underpricing the issuer must offer to keep you in the market.

\notation
\begin{itemize}
  \item $P_0=40$: offer price per share of both IPOs (\$).
  \item $u=7$: underpricing of the undervalued issue (true value $P_0+u=47$); $o=5$: overpricing of the overvalued one (true value $P_0-o=35$).
  \item $Q=1{,}000$: shares ordered in each issue; $\pi$: fraction of the order filled in the underpriced issue ($\pi=1/2$); $\pi'$: fraction filled in the overpriced issue ($\pi'=1$).
  \item $G^{\text{full}}$: gain with both orders filled; $E[G]$: gain you actually expect.
\end{itemize}

\step{1}{Gain if both orders were filled} You make $u$ per share on one issue and lose $o$ per share on the other:
\begin{align}
  G^{\text{full}} &= Q\,u - Q\,o \label{t3e1-full}\\
                  &= (1{,}000)(7)-(1{,}000)(5) \nonumber\\
                  &= 7{,}000-5{,}000 \nonumber\\
                  &= 2{,}000. \nonumber
\end{align}

\answer{(a)} A gain of \$2{,}000 (\$7{,}000 on the underpriced issue minus \$5{,}000 on the overpriced one).

\step{2}{Gain you actually expect} Whichever issue is the underpriced one is oversubscribed (informed buyers also want it) and fills only $\pi Q$; the overpriced one fills $\pi' Q$. Since you cannot tell which is which, this holds whatever the labels are. The shares you receive in each issue are
\begin{align}
  \pi Q &= (0.5)(1{,}000) = 500 \text{ shares of the underpriced issue}, \nonumber\\
  \pi' Q &= (1)(1{,}000) = 1{,}000 \text{ shares of the overpriced issue}, \nonumber
\end{align}
so the gain is
\begin{align}
  E[G] &= \pi Q\,u - \pi' Q\,o \label{t3e1-exp}\\
       &= (0.5)(1{,}000)(7)-(1)(1{,}000)(5) \nonumber\\
       &= 3{,}500-5{,}000 \nonumber\\
       &= -1{,}500. \nonumber
\end{align}
The full-allocation gain of \$2{,}000 becomes an expected loss of \$1{,}500. (The outcome is in fact certain here: whatever the labels, you get 500 shares of the good issue and 1{,}000 of the bad one.)

\answer{(b)} An expected (and certain) loss of \$1{,}500.

\step{3}{Principle, and the underpricing that would make you participate} This is the \emph{winner's curse} (\pt{``maldi\c{c}\~ao do vencedor''}): an uninformed buyer receives a full allocation exactly when informed buyers stay away, that is, when the issue is overpriced. To break even, \eqref{t3e1-exp} must be zero, which requires an underpricing $u^{*}$ with
\begin{align}
  \pi Q\,u^{*} - \pi' Q\,o &= 0 \nonumber\\
  \pi Q\,u^{*} &= \pi' Q\,o \nonumber\\
  u^{*} &= \frac{\pi'}{\pi}\,o \label{t3e1-ustar}\\
        &= \frac{(1)}{(0.5)}\,(5) \nonumber\\
        &= 10. \nonumber
\end{align}
The average initial return per share across the two issues would then be
\begin{align}
  \frac{u^{*}+(-o)}{2} = \frac{(10)-(5)}{2} = \frac{5}{2} = 2.5, \nonumber
\end{align}
positive on average, yet zero for the uninformed buyer.

\answer{(c)} The winner's curse: an uninformed investor is allocated more of the bad (overpriced) issues than of the good ones, so average underpricing is needed just for her to break even.

\intuition Informed investors skim the underpriced issues, so what an uninformed investor actually receives is tilted towards the bad issues. Issuers must underprice on average just to keep uninformed investors in the market (Rock's model, Amaro de Matos \S6.2.2; lecture block 3.10: $P_0<E[P_1]$ even though the uninformed only break even).

\pitfalls
\begin{itemize}
  \item Computing the expected gain with the ordered quantities ($1{,}000\times7-1{,}000\times5$), i.e.\ ignoring that the allocation depends on the quality of the issue.
  \item Using a 50\% probability that \emph{each} issue is rationed: only the underpriced issue is rationed, and that is the whole point.
  \item Concluding that IPOs are ``overpriced on average'': the average issue is underpriced (+2.5 per share in Step~3); it is the uninformed investor's \emph{allocation-weighted} return that is zero.
\end{itemize}
\refbib{Ross Cap.~15, \S15.5 ``IPOs e subprecifica\c{c}\~ao'', subsection ``Por que existe subprecifica\c{c}\~ao?'', pp.~507--508; lecture Topic~3, ``IPO underpricing''.}

% =============================================================================
\exsec{2. Investment and available cash --- Grinblatt \& Titman 19.4, p.~687}
%% BEGIN-STATEMENT 2
\stmtbegin
\textbf{Source:} Grinblatt, Ch. 19, Exercise 19.4, p.~687 (PDF 700) · Answer in book: no

Classical finance theory suggests that firms take projects with positive NPVs regardless of the
amount of cash the firm has available. However, empirical evidence suggests that the amount that
firms invest is heavily dependent on their available cash flows. Why might this be?
\stmtend
%% END-STATEMENT

\asked (a) The NPV rule says a firm should take every positive-NPV project whatever its cash, yet in the data investment moves with cash flow. Explain why: name the friction that makes internal cash cheaper than outside money, and show how it can make investment depend on cash.

\approach The friction is adverse selection in the issue of new securities (the lemons/pooling model of lecture block 3.2 and the Myers--Majluf pecking order of block 3.5). Outside investors cannot tell a good firm from a bad one, so they price any new claim at the pooled value; a good firm therefore sells its claims at a discount. That discount is proportional to the amount raised outside, so the more of the investment the firm pays for with its own cash, the cheaper its funding. A short model with the lecture's numbers makes the point quantitative; other frictions (moral hazard, issue costs) give the same prediction and are listed at the end.

\notation
\begin{itemize}
  \item Lecture model (Topic~3): project costs $I$, pays $R$ on success and 0 on failure; good type succeeds with probability $p$, bad type with $q<p$; $\alpha\in(0,1)$: share of good types; $m\equiv\alpha p+(1-\alpha)q$: investors' prior success probability.
  \item $A\in[0,I)$: internal cash (retained earnings) of the firm, the same for both types (the bad type copies the good type); $R_u$: repayment to outside investors on success.
  \item $U_G(A)$: good type's net payoff (gross payoff minus the $A$ it puts in).
\end{itemize}

\assume Both types hold the same cash $A$ and pool on the same external claim (as in the lecture's pecking-order slide, the bad type ``follows suit''); outside investors break even.

\step{1}{Why ``invest in every positive-NPV project'' fails} The classical rule assumes external funds are priced fairly. When managers know more than investors, new securities sell at the price of the average firm, so a good firm sells its claims too cheaply (Myers--Majluf, GT Result~19.10). Internal cash carries no such discount, so the cost of funds depends on how much of the investment the firm can pay for itself.

\step{2}{Good type's payoff with internal cash $A$} The firm borrows $I-A$ and promises $R_u$ on success. Because both types issue the same claim, investors cannot tell them apart and expect success with the pooled probability $m$; they break even when
\begin{align}
  m R_u &= I-A \label{t3e2-be}\\
  R_u &= \frac{I-A}{m}. \nonumber
\end{align}
The good type keeps $R-R_u$ on success (probability $p$) and has put in $A$. Her net payoff, substituting \eqref{t3e2-be} and then adding and subtracting $I$:
\begin{align}
  U_G(A) &= p(R-R_u) - A \nonumber\\
         &= pR - p\left(\frac{I-A}{m}\right) - A \nonumber\\
         &= pR - I + (I-A) - \frac{p}{m}(I-A) \nonumber\\
         &= (pR-I) - (I-A)\left(\frac{p}{m}-1\right). \label{t3e2-UG}
\end{align}
The first term is the NPV; the second, the adverse-selection discount, is proportional to the amount raised outside, $I-A$, and vanishes at $A=I$. It is positive because $m<p$: the good type pays for the bad type's higher failure rate.

\step{3}{Numbers from the lecture example} Take $R=50$, $I=30$, $p=0.8$, $q=0.7$, $\alpha=0.5$. The pooled probability and the NPV are
\begin{align}
  m &= (0.5)(0.8)+(0.5)(0.7) = 0.40+0.35 = 0.75, \nonumber\\
  pR-I &= (0.8)(50)-(30) = 40-30 = 10, \nonumber\\
  \frac{p}{m}-1 &= \frac{(0.8)}{(0.75)}-1 = 1.066667-1 = 0.066667. \nonumber
\end{align}
Substituting into \eqref{t3e2-UG}:
\begin{align}
  U_G(0) &= (10)-(30)(0.066667) = 10-2 = 8, \nonumber\\
  U_G(15) &= (10)-(15)(0.066667) = 10-1 = 9. \nonumber
\end{align}
Half the investment paid internally halves the discount.

\step{4}{Investment that depends on cash} Take the lecture's breakdown case $q=0.5$, $\alpha=0.2$:
\begin{align}
  m &= (0.2)(0.8)+(0.8)(0.5) = 0.16+0.40 = 0.56, \nonumber\\
  mR &= (0.56)(50) = 28 < 30 = I. \nonumber
\end{align}
With $A=0$ even handing investors the whole income ($R_u=R$) returns only 28 for a loan of 30, so nobody lends. Limited liability caps the repayment at $R_u\le R$, so by \eqref{t3e2-be} lending requires
\begin{align}
  \frac{I-A}{m} &\le R \nonumber\\
  I-A &\le mR \nonumber\\
  A &\ge I-mR = (30)-(28) = 2. \label{t3e2-Amin}
\end{align}
Two firms with the same positive-NPV project ($pR-I=10$) behave differently: one with $A\ge 2$ invests, one with $A<2$ cannot. Investment is a function of cash.

\step{5}{Other frictions with the same prediction} (i) Moral hazard: pledgeable income is limited, so borrowing capacity is a multiple of net worth (Tirole \S3.4, $I\le kA$). (ii) Issue (flotation) costs make external money dearer than internal money. (iii) Agency in the other direction: managers with free cash flow overinvest. All three make investment rise with internal cash; GT emphasise the information channel.

\answer{(a)} Because outside finance is dearer than internal cash. Under asymmetric information new claims are priced at the pooled value, so a good firm pays a discount $(I-A)(p/m-1)$ that grows with the amount raised outside (Step~2: 2 with no cash, 1 with half the investment in cash); in bad pools it cannot raise funds at all unless $A\ge I-mR$ (Step~4: $A\ge2$). Moral hazard and issue costs reinforce the same link.

\intuition External finance is costly because outsiders cannot tell good firms from bad ones and price every new claim at the pooled value. Internal cash avoids the discount, so firms invest out of cash first (the pecking order), and investment tracks cash flow even when the NPV rule says it should not.

\pitfalls
\begin{itemize}
  \item Answering ``because cash is cheaper'' without saying \emph{why}: in a perfect market (MM) a dollar of cash and a dollar raised outside cost the same.
  \item Claiming the good firm rejects a positive-NPV project because the project is bad: the project's NPV is unchanged; what changes is the transfer to new investors.
  \item Pricing the claim at the good type's probability $p$: investors only know the pool, so the break-even uses $m$.
\end{itemize}
\refbib{Grinblatt \& Titman Ch.~19, \S19.5 ``The Information Content of the Debt-Equity Choice'', pp.~671--678 (Example~19.8 and Result~19.10, pp.~676--677); Tirole \S6.2.2 pp.~244--249 and \S3.4 pp.~127--129; lecture Topic~3, ``Pecking-order hypothesis''.}

% =============================================================================
\exsec{3. Leverage, agency costs and information (IST) --- Berk 16.29, p.~496}
%% BEGIN-STATEMENT 3
\stmtbegin
\textbf{Source:} Berk, Ch. 16, Exercise 16.29 (starred), p.~496 (PDF 527) · Answer in book: no

Info Systems Technology (IST) manufactures microprocessor chips for use in appliances and other
applications. IST has no debt and 100 million shares outstanding. The correct price for these shares
is either \$14.50 or \$12.50 per share. Investors view both possibilities as equally likely, so the
shares currently trade for \$13.50.

IST must raise \$500 million to build a new production facility. Because the firm would suffer a
large loss of both customers and engineering talent in the event of financial distress, managers
believe that if IST borrows the \$500 million, the present value of financial distress costs will
exceed any tax benefits by \$20 million. At the same time, because investors believe that managers
know the correct share price, IST faces a lemons problem if it attempts to raise the \$500 million
by issuing equity.

\begin{enumerate}
\def\labelenumi{\alph{enumi}.}
\item
  Suppose that if IST issues equity, the share price will remain \$13.50. To maximize the long-term
  share price of the firm once its true value is known, would managers choose to issue equity or
  borrow the \$500 million if:

  \begin{enumerate}
  \def\labelenumii{\roman{enumii}.}
  \tightlist
  \item
    They know the correct value of the shares is \$12.50?
  \item
    They know the correct value of the shares is \$14.50?
  \end{enumerate}
\item
  Given your answer to part (a), what should investors conclude if IST issues equity? What will
  happen to the share price?
\item
  Given your answer to part (a), what should investors conclude if IST issues debt? What will happen
  to the share price in that case?
\item
  How would your answers change if there were no distress costs, but only the tax benefits of
  leverage?
\end{enumerate}
\stmtend
%% END-STATEMENT

\asked (a) Taking the equity issue price as fixed at \$13.50, compare the long-run share price (once the truth is out) under an equity issue and under a \$500 million loan, for a manager who knows the shares are worth \$12.50 (i) and \$14.50 (ii), and say which she picks. (b) What investors should infer from an equity issue and where the price goes. (c) The same for a debt issue. (d) Redo the reasoning when debt brings only tax benefits and no distress costs.

\approach This is the lemons problem in an equity issue (Berk \S16.6, lecture blocks 3.2 and 3.5). Equity is the claim whose value depends most on what the manager knows, so selling it at a pooled price is a gain for an overvalued firm and a loss for an undervalued one. Debt avoids that transfer but here costs \$20 million of net distress costs. Part (a) compares the two payoffs type by type; parts (b)--(c) use the fact that rational investors read the choice of security as a signal (a separating outcome); part (d) removes the cost that made the separation possible.

\notation
\begin{itemize}
  \item $N_0=100$ million shares; $v\in\{12.50,\,14.50\}$: true price per share before the financing (equally likely); $P_m=13.50$: market price.
  \item $K=500$ million: amount to raise; $C=20$ million: PV of distress costs net of tax benefits if IST borrows $K$.
  \item $n$: new shares issued at $P_m$; $P^E(v)$, $P^D(v)$: long-run share price (after the truth is known) under equity and debt financing; $X>0$: net tax benefit of debt in part (d).
\end{itemize}

\assume The true prices $v$ already include the facility's NPV, so the \$500 million raised is invested in assets worth what is paid for them. Under debt, the \$500 million debt claim offsets the \$500 million of cash raised, and equity loses $C$.

\step{1}{(a) Price of each policy once the truth is known} Equity: IST sells $n$ shares at $P_m$; once the truth is known the firm's equity is worth the old value $N_0v$ plus the cash raised $K$ (invested at zero NPV beyond what $v$ already includes), spread over $N_0+n$ shares:
\begin{align}
  n &= \frac{K}{P_m} = \frac{(500)}{(13.50)} = 37.037 \text{ million}, \label{t3e3-n}\\
  P^E(v) &= \frac{N_0 v + K}{N_0+n}. \label{t3e3-PE}
\end{align}
Debt: the \$500 million of assets bought is offset by the \$500 million debt claim, so equity keeps all the old shares and bears only the net distress cost $C$:
\begin{align}
  P^D(v) &= \frac{N_0 v + K - K - C}{N_0} = \frac{N_0 v - C}{N_0}. \label{t3e3-PD}
\end{align}

\textit{(i) $v=12.50$.} Substituting into \eqref{t3e3-PE} and \eqref{t3e3-PD}:
\begin{align}
  P^E(12.50) &= \frac{(100)(12.50)+(500)}{(100)+(37.037)} = \frac{1{,}750}{137.037} = 12.77, \nonumber\\
  P^D(12.50) &= \frac{(100)(12.50)-(20)}{(100)} = 12.30. \nonumber
\end{align}
Since $12.77>12.30$, managers \textbf{issue equity}: selling overvalued shares transfers $N_0[P^E(12.50)-12.50]=(100)(12.7703-12.50)=27.03$ million from new to old shareholders, while debt would cost them $C=20$ million.

\textit{(ii) $v=14.50$.}
\begin{align}
  P^E(14.50) &= \frac{(100)(14.50)+(500)}{(100)+(37.037)} = \frac{1{,}950}{137.037} = 14.23, \nonumber\\
  P^D(14.50) &= \frac{(100)(14.50)-(20)}{(100)} = \frac{1{,}430}{100} = 14.30. \nonumber
\end{align}
Since $14.30>14.23$, managers \textbf{borrow}: the net distress cost, $C/N_0=20/100=0.20$ per share, is smaller than the loss from selling undervalued shares, $14.50-14.23=0.27$ per share.

\answer{(a)} (i) Value \$12.50: issue equity (\$12.77 against \$12.30 with debt). (ii) Value \$14.50: borrow (\$14.30 against \$14.23 with equity).

\step{2}{(b) An equity issue} By Step~1 only the low-value firm issues equity, so investors who know the managers' incentives infer $v=12.50$ and the price falls from 13.50 to \textbf{12.50}. Check that this price is consistent: at 12.50 IST would sell
\begin{align}
  \frac{K}{12.50} = \frac{(500)}{(12.50)} = 40 \text{ million shares}, \qquad
  \frac{N_0(12.50)+K}{N_0+40} = \frac{(1{,}250)+(500)}{(100)+(40)} = \frac{1{,}750}{140} = 12.50, \nonumber
\end{align}
so the new shareholders pay exactly what their shares are worth.

\answer{(b)} An equity issue reveals that the shares are worth \$12.50; the price falls from \$13.50 to \$12.50.

\step{3}{(c) A debt issue} Only the high-value firm borrows, so investors infer $v=14.50$ and price in the distress cost: the price rises to $P^D(14.50)=$ \textbf{14.30}.

\answer{(c)} A debt issue reveals that the shares are worth \$14.50; net of the \$0.20 of distress costs per share the price rises to \$14.30.

\step{4}{(d) Tax benefits only} Now debt adds a net tax benefit $X>0$ instead of costing $C$, so \eqref{t3e3-PD} becomes
\begin{align}
  P^D(v) = \frac{N_0v+X}{N_0} = v+\frac{X}{N_0} = v+\frac{X}{100}. \nonumber
\end{align}
\textit{Part (a) redone.} The high-value firm borrows, since $14.50+X/100>14.50>14.23=P^E(14.50)$. With the issue price fixed at 13.50, the low-value firm issues equity only if
\begin{align}
  P^E(12.50) &> 12.50 + \frac{X}{100} \nonumber\\
  12.7703 - 12.50 &> \frac{X}{100} \nonumber\\
  X &< 100\,(0.2703) = 27.03 \text{ million}. \label{t3e3-X}
\end{align}
\textit{Parts (b)--(c) redone.} By the logic of (b), an equity issue would reveal the low type and be priced at 12.50, which leaves the low type with $12.50<12.50+X/100$ (the value under debt). So in equilibrium \textbf{both types borrow}: no equity is issued, a debt issue carries no information, and the price moves only by the tax benefit, to
\begin{align}
  \tfrac12(12.50)+\tfrac12(14.50)+\frac{X}{100} = 13.50+\frac{X}{100}. \nonumber
\end{align}
The answers to (b)--(c) change: equity would be read as bad news but is never chosen, and debt is no longer a signal of a high value.

\answer{(d)} Debt now dominates for both types. (Against a fixed \$13.50 equity price the low type would still issue equity if $X<\$27.03$ million, but investors price any equity issue at \$12.50, so it borrows too.) Both types issue debt, the issue reveals nothing, and the price rises only by the tax benefit, to $13.50+X/100$ per share: the pecking order.

\intuition Equity is the security whose value is most sensitive to what managers know, so issuing it says ``our shares are overpriced''. With distress costs, only the overvalued firm accepts that stigma, and the choice of security separates the types. Without distress costs, debt dominates for both, and the market learns nothing from the choice: this is the pecking order.

\pitfalls
\begin{itemize}
  \item Forgetting the cash raised in the equity case: the post-issue value is $N_0v+K$, not $N_0v$, divided by $N_0+n$ shares.
  \item Subtracting the \$500 million debt from equity but not adding the \$500 million of new assets (or vice versa) in the debt case; the two cancel and only $C$ is left.
  \item In (b)--(c), keeping the price at \$13.50: once the security choice reveals the type, the price must move to the revealed value.
\end{itemize}
\refbib{Berk Ch.~16, \S16.6 ``Additional Consequences of Leverage: Agency Costs and Information'', subsection ``Debt and Information'', pp.~485--487; Ch.~14 lemons principle, p.~431.}

% =============================================================================
\exsec{4. Issuing undervalued equity --- Grinblatt \& Titman 19.6, p.~687}
%% BEGIN-STATEMENT 4
\stmtbegin
\textbf{Source:} Grinblatt, Ch. 19, Exercise 19.6, p.~687 (PDF 700) · Answer in book: no

ABC Industries is considering an investment that requires the firm to issue new equity. The project
will cost \$100, but will add \$120 to the firm's value. Although management believes the firm's
value is \$1,000 without the new project, outside investors value the firm at \$600 without the
project. If the firm currently has 100 shares outstanding, how many new shares must it issue to
finance the project? Now assume that the true value of the firm will become known to the market
shortly after the new equity has been issued. What will the firm's stock price be at this time if it
chooses to finance this new investment? What will the stock price be if it chooses to pass up the
investment?
\stmtend
%% END-STATEMENT

\asked (a) How many new shares ABC must sell to raise the \$100 the project costs, given that outsiders value the firm at \$600. (b) The share price once the true value is known if ABC issues and invests. (c) The share price once the true value is known if ABC passes up the project.

\approach This is the Myers--Majluf underinvestment problem (lecture block 3.5; GT \S19.5). New shares are sold at the outsiders' (too low) valuation, so the new shareholders buy a claim on the undervalued assets in place as well as on the project. Comparing the post-revelation price per old share with and without the project shows whether the transfer to new shareholders exceeds the project's NPV; if it does, managers acting for old shareholders pass up a positive-NPV project.

\notation
\begin{itemize}
  \item $N_0=100$ shares; $V_{\text{in}}=1{,}000$: management's value of the firm without the project; $V_{\text{out}}=600$: outside investors' value without the project.
  \item $K=100$: cost of the project (raised with new equity); $V_P=120$: PV the project adds to the firm, so $\text{NPV}=V_P-K=20$.
  \item $P_{\text{issue}}$: issue price; $n$: new shares; $P_{\text{take}}$, $P_{\text{pass}}$: share price once the truth is known, with and without the project.
\end{itemize}

\assume The shares are sold at the outsiders' current value per share, $P_{\text{issue}}=V_{\text{out}}/N_0$, as in GT's own illustration (p.~676), where new shares are sold at the current market price. (The alternative is in Step~5.)

\step{1}{New shares} The issue price and the number of shares needed to raise $K$:
\begin{align}
  P_{\text{issue}} &= \frac{V_{\text{out}}}{N_0} = \frac{(600)}{(100)} = 6.00, \label{t3e4-Pi}\\
  n &= \frac{K}{P_{\text{issue}}} = \frac{(100)}{(6.00)} = 16.67. \label{t3e4-n}
\end{align}

\answer{(a)} 16.67 new shares at \$6.00 each.

\step{2}{Price if the project is taken} Once the truth is known the firm is worth $V_{\text{in}}+V_P$, split over $N_0+n$ shares (GT's ``share value taking the project''):
\begin{align}
  P_{\text{take}} &= \frac{V_{\text{in}}+V_P}{N_0+n} \label{t3e4-take}\\
                  &= \frac{(1{,}000)+(120)}{(100)+(16.67)} \nonumber\\
                  &= \frac{1{,}120}{116.67} = 9.60. \nonumber
\end{align}

\answer{(b)} \$9.60 per share if ABC issues and invests.

\step{3}{Price if the project is passed up} No shares are issued and the firm is worth $V_{\text{in}}$ once the truth is known:
\begin{align}
  P_{\text{pass}} &= \frac{V_{\text{in}}}{N_0} = \frac{(1{,}000)}{(100)} = 10.00. \label{t3e4-pass}
\end{align}
Since $9.60<10.00$, old shareholders are better off if the firm \textbf{passes up} the positive-NPV project.

\answer{(c)} \$10.00 per share if ABC passes up the project, so the old shareholders prefer to pass it up.

\step{4}{Where the \$0.40 goes} New shareholders pay 100 for a fraction of the firm
\begin{align}
  \frac{n}{N_0+n} = \frac{(16.67)}{(116.67)} = 0.142857 \quad(=1/7), \nonumber
\end{align}
of a firm worth 1{,}120, that is, for a stake worth $(0.142857)(1{,}120)=160$. The transfer to them is $160-100=60$, larger than the NPV of $120-100=20$, so old shareholders lose $60-20=40$ in total, or $40/100=0.40$ per share ($=10.00-9.60$). The firm takes the project only if \eqref{t3e4-take} with $n=K/P_{\text{issue}}$ is at least \eqref{t3e4-pass}. Cross-multiplying (both denominators are positive) and cancelling $N_0V_{\text{in}}$:
\begin{align}
  \frac{V_{\text{in}}+V_P}{N_0+K/P_{\text{issue}}} &\ge \frac{V_{\text{in}}}{N_0} \nonumber\\
  N_0(V_{\text{in}}+V_P) &\ge V_{\text{in}}\Bigl(N_0+\frac{K}{P_{\text{issue}}}\Bigr) \nonumber\\
  N_0V_{\text{in}}+N_0 V_P &\ge N_0V_{\text{in}}+\frac{V_{\text{in}}K}{P_{\text{issue}}} \nonumber\\
  N_0 V_P &\ge \frac{V_{\text{in}}K}{P_{\text{issue}}} \nonumber\\
  P_{\text{issue}} &\ge \frac{V_{\text{in}}K}{N_0 V_P}= \frac{(1{,}000)(100)}{(100)(120)} = \frac{100{,}000}{12{,}000} = 8.33, \label{t3e4-thr}
\end{align}
so the market's undervaluation (6 against 10) is far beyond what the project can absorb: the project is taken only if outsiders' price is at least \$8.33.

\step{5}{If the market prices the project in} If outsiders value the firm with the project at $V_{\text{out}}+V_P=600+120=720$ after the issue, the issue price must satisfy $N_0P_{\text{issue}}+K=720$:
\begin{align}
  P_{\text{issue}} &= \frac{720-K}{N_0} = \frac{(720)-(100)}{(100)} = 6.20, \nonumber\\
  n &= \frac{(100)}{(6.20)} = 16.13, \nonumber\\
  P_{\text{take}} &= \frac{(1{,}120)}{(100)+(16.13)} = \frac{1{,}120}{116.13} = 9.64 < 10. \nonumber
\end{align}
The conclusion does not change.

\intuition Equity is a claim on the assets in place too. When outsiders undervalue those assets, selling equity hands part of their hidden value to the new shareholders; if that transfer exceeds the project's NPV, the managers acting for the old shareholders forgo the project (Myers--Majluf underinvestment). Internal cash or safe debt would avoid the transfer.

\pitfalls
\begin{itemize}
  \item Using management's value (\$10 per share) as the issue price: the shares are sold at what outsiders will pay, \$6.
  \item Forgetting to add the project's \$120 (or adding only its NPV of \$20 while also counting the \$100 of cash, which is spent on the project) when computing the firm's value after investing.
  \item Concluding that the project should be rejected on its merits: its NPV is $+20$; it is rejected only because of the \$60 transfer to new shareholders.
\end{itemize}
\refbib{Grinblatt \& Titman Ch.~19, \S19.5, ``Adverse Selection and the Pecking Order'' discussion and Example~19.8, pp.~676--677, Result~19.10, p.~677; Amaro de Matos \S3.2.3 pp.~86--88 (rule (3.20)).}

% =============================================================================
\exsec{5. Dividend suspension and the stock price --- Ross CQ 17.6, p.~614}
%% BEGIN-STATEMENT 5
\stmtbegin
\textbf{Source:} Ross, Cap. 17, Exercise CQ 17.6 (Revisão de conceitos e questões instigantes),
p.~614 (PDF 638) · Answer in book: no

Mês passado, a Companhia de Energia Elétrica de Angra dos Reis, que vinha tendo problemas com custos
acima dos previstos em uma usina nuclear em construção, anunciou que estava ``suspendendo
temporariamente os pagamentos devido à diminuição dos fluxos de caixa associados a esse programa de
investimentos''. O preço da ação da empresa caiu de \$28,50 para \$25 quando esse anúncio foi feito.
Como você interpretaria essa mudança no preço da ação (ou seja, em sua opinião, o que teria causado
isso)?
\stmtend
%% END-STATEMENT

\asked (a) The utility suspended its dividend, citing lower cash flows from a nuclear plant over budget, and its share price fell from \$28.50 to \$25. Interpret the fall: what caused it?

\approach Two benchmarks frame the answer. MM dividend irrelevance says the payout alone cannot change value in a perfect market, so the price drop must come from something else. The information content of dividends (dividends as a costly signal, lecture block 3.7) supplies it: managers smooth dividends and cut them only when they expect lower cash flows, so a suspension is bad news about future cash flows. Size the reaction first, then explain it.

\notation
\begin{itemize}
  \item $P_b=28.50$, $P_a=25$: share price before and after the announcement (\$); $r_{\text{ann}}$: announcement return.
\end{itemize}

\step{1}{Size of the reaction} The change in price and the announcement return:
\begin{align}
  P_a-P_b &= (25)-(28.50) = -3.50, \nonumber\\
  r_{\text{ann}} &= \frac{P_a-P_b}{P_b} = \frac{(-3.50)}{(28.50)} = -12.28\%. \label{t3e5-r}
\end{align}

\step{2}{What caused it} In a perfect market (MM dividend irrelevance and the homemade-dividend argument) changing the payout alone does not change value: the cash stays in the firm. The drop therefore reflects \emph{information}. Firms smooth dividends and cut them only with great reluctance, so a suspension tells the market that management no longer expects the current dividend to be sustainable. Here the firm says why: cost overruns on the nuclear plant have reduced cash flows (\pt{``diminui\c{c}\~ao dos fluxos de caixa''}). Investors revise expected future cash flows, and so expected future dividends, downward, and the price falls with their present value. The price fell because future dividends are expected to be lower, not because the payout ratio changed.

\step{3}{Secondary effects} (i) Clientele: investors who held the stock for current income may sell, adding temporary pressure. (ii) Keeping the cash avoids a costly external issue to finish the plant (pecking order), which by itself would support the price; the net fall of 12.28\% shows the bad news dominated.

\answer{(a)} The \$3.50 (12.28\%) fall is the market's reaction to the information in the suspension: lower expected future cash flows (hence lower future dividends) from the cost overruns, not the change in payout itself.

\intuition Dividends work as a signal because they are costly to maintain when cash flows are weak: only a firm that expects strong cash flows can keep paying. Cutting the dividend reveals the opposite, so the stock reprices on the news about cash flows, not on the dividend itself (lecture: payout policy as a signalling device).

\pitfalls
\begin{itemize}
  \item Saying the price fell ``because shareholders lost the dividend'': under MM the retained cash is still theirs; the payout change alone does not destroy value.
  \item Computing the return on the after-price, $3.50/25=14\%$, instead of on the before-price, $3.50/28.50=12.28\%$.
  \item Ignoring that the cut was partly anticipated (the overruns were known): the price reacts only to the \emph{new} information in the announcement.
\end{itemize}
\refbib{Ross Cap.~17, \S17.5, ``Conte\'udo informacional dos dividendos'', pp.~593--594; lecture Topic~3, ``Other signaling devices: payout policy''.}

% =============================================================================
\exsec{6. Privately known private benefit and market breakdown --- Tirole 6.1, p.~273}
%% BEGIN-STATEMENT 6
\stmtbegin
\textbf{Source:} Tirole, Ch. 6, Exercise 6.1, p.~273 (PDF 286) · Answer in book: yes, p.~595

Section 6.2 illustrated the possibility of market breakdown without the possibility of signaling.
This exercise supplies another illustration. Let us consider the fixed-investment model of Section
3.2 and assume that only the borrower knows the private benefit associated with misbehavior. When
the borrower has private information about this parameter, lenders are concerned that this private
benefit might be high and induce the borrower to misbehave. In the parlance of information
economics, the ``bad types'' are the types of borrower with high private benefit. We study the case
of two possible levels of private benefit (see Exercise 6.2 for an extension to a continuum of
possible types). The borrower wants to finance a fixed-size project costing \(I\), and, for
simplicity, has no equity (\(A = 0\)). The project yields \(R\) (success) or 0 (failure). The
probability of success is \(p_H\) or \(p_L\), depending on whether the borrower works or shirks,
with \(\Delta p \equiv p_H - p_L > 0\). There is no private benefit when working. The private
benefit \(B\) enjoyed by the borrower when shirking is either \(B_L > 0\) or \(B_H > B_L\). The
borrower will be labeled a ``good borrower'' when \(B = B_L\) and a ``bad borrower'' when
\(B = B_H\). At the date of contracting, the borrower knows the level of her private benefit, while
the capital market puts (common knowledge) probabilities \(\alpha\) that the borrower is a good
borrower and \(1 - \alpha\) that she is a bad borrower. All other parameters are common knowledge
between the borrower and the lenders.

To make things interesting, let us assume that under asymmetric information, the lenders are
uncertain about whether the project should be funded:

\[p_H\left(R - \frac{B_H}{\Delta p}\right) < I < p_H\left(R - \frac{B_L}{\Delta p}\right). \qquad (1)\]

Assume that investors cannot break even if the borrower shirks:

\[p_L R < I. \qquad (2)\]

\begin{enumerate}
\def\labelenumi{(\roman{enumi})}
\item
  Note that the investor cannot finance only good borrowers. Assume that the entrepreneur receives
  no reward in the case of failure (this is indeed optimal); consider the effect of rewards \(R_b\)
  in the case of success that are (a) smaller than \(B_L/\Delta p\), (b) larger than
  \(B_H/\Delta p\), (c) between these two values.
\item
  Show that there exists \(\alpha^*\), \(0 < \alpha^* < 1\), such that
\end{enumerate}

\begin{itemize}
\tightlist
\item
  no financing occurs if \(\alpha < \alpha^*\),
\item
  financing is an equilibrium if \(\alpha \geqslant \alpha^*\).
\end{itemize}

\begin{enumerate}
\def\labelenumi{(\roman{enumi})}
\setcounter{enumi}{2}
\tightlist
\item
  Describe the ``cross-subsidies'' between types that occur when borrowing is feasible.
\end{enumerate}

\paragraph{Context needed to solve it}

\emph{From Tirole §3.2.1--3.2.2 (the fixed-investment model), pp.~115--117 (PDF 128--130):}

\textbf{3.2.1 The Fixed-Investment Model.} Variants of the following entrepreneurial model will be
used in the following: an entrepreneur (also called the ``insider'' or the ``borrower,'' ``she'')
has a project. This project requires a fixed investment \(I\). The entrepreneur initially has
``assets'' (``cash on hand'' or ``net worth'') \(A < I\). For the moment we interpret these assets
as being cash or liquid securities that can be used toward covering the cost of investment.
{[}\ldots{]} The entrepreneur's cash can either be invested in the project or used for consumption.
To implement the project the entrepreneur must borrow \(I - A\) from lenders. {[}\ldots{]}

\emph{Project.} If undertaken, the project either succeeds, that is, yields verifiable income
\(R > 0\), or fails, that is, yields no income. The probability of success is denoted by \(p\). The
project is subject to moral hazard. The entrepreneur can ``behave'' (``work,'' ``exert effort,''
``take no private benefit'') or ``misbehave'' (``shirk,'' ``take a private benefit''); or,
equivalently, the entrepreneur chooses between a project with a high probability of success and
another project which ceteris paribus she prefers (is easier to implement, is more fun, has greater
spinoffs in the future for the entrepreneur, benefits a friend, delivers perks, is more
``glamorous,'' etc.) but has a lower probability of success.

Behaving yields probability \(p = p_H\) of success and no private benefit to the entrepreneur, and
misbehaving results in probability \(p = p_L < p_H\) of success and private benefit \(B > 0\)
(measured in units of account) to the entrepreneur. In the ``effort interpretation,'' \(B\) can also
be interpreted as a disutility of effort saved by the entrepreneur when shirking. Let
\(\Delta p \equiv p_H - p_L\).

\emph{Preferences and the loan agreement.} Both the borrower and the potential lenders (or
``investors'') are risk neutral. For notational simplicity, there is no time preference; the rate of
return expected by investors (which is also the riskless rate, due to risk neutrality) is taken to
be 0. The borrower is protected by limited liability, and so her income cannot take negative values.

Lenders behave competitively in the sense that the loan, if any, makes zero profit. That is, we have
in mind that several prospective lenders compete for issuing a loan to the borrower, and that, if
the most attractive loan offer made a positive profit, the borrower could turn to an alternative
lender and offer to switch for a slightly lower interest rate. {[}\ldots{]}

Let us turn to the loan contract. A contract first stipulates whether the project is financed. If
so, it further specifies how the profit is shared between the lenders and the borrower. The
borrower's limited liability will imply that both sides receive 0 in the case of failure (the gross
payoffs are the \emph{ex post} monetary payoffs and take no account of past investments and private
benefit). Intuitively, there is no point in specifying a positive transfer from the lenders to the
borrower, as such a transfer can only weaken incentives, while it has no insurance benefit under
risk neutrality. This property will be proved more rigorously and is here taken for granted. In the
case of success, the two parties share the profit \(R\); \(R_b\) goes to the borrower and \(R_l\) to
the lenders. To sum up, we posit an incentive scheme for the entrepreneur of the following form:
\(R_b\) in the case of success, 0 in the case of failure.

The zero-profit constraint for the lenders can be written as

\[p_H R_l = I - A,\]

assuming that the loan agreement induces the borrower to behave (which under our assumptions will be
the case). {[}\ldots{]} We summarize the timing in Figure 3.1.

\emph{Figure 3.1 (timeline, left to right):} Loan agreement (sharing rule) → Investment → Moral
hazard → Outcome.

We assume that the project is viable only in the absence of moral hazard. That is, the project has
positive NPV if the entrepreneur behaves,

\[p_H R - I > 0, \qquad (3.1)\]

but negative NPV, even if one includes the borrower's private benefit, if she does not,

\[p_L R - I + B < 0. \qquad (3.2)\]

{[}\ldots{]}

\textbf{3.2.2 The Lenders' Credit Analysis.} Because the project has negative NPV in the case of
misbehavior, the loan agreement must be careful to preserve enough of a stake for the borrower in
the enterprise. The borrower faces the following tradeoff once the financing has been secured: by
misbehaving, she obtains private benefit \(B\), but she reduces the probability of success from
\(p_H\) to \(p_L\). Because she has stake \(R_b\) in the firm's income (she receives \(R_b\) in the
case of success and 0 in the case of failure), the borrower will therefore behave if the following
``incentive compatibility constraint'' is satisfied:

\[p_H R_b \geqslant p_L R_b + B \quad \text{or} \quad (\Delta p) R_b \geqslant B. \qquad (\text{IC}_b)\]

From this incentive compatibility constraint we infer that the highest income in the case of success
that can be pledged to the lenders without jeopardizing the borrower's incentives is
\(R - \dfrac{B}{\Delta p}\). The \emph{(expected) pledgeable income} is then
\(\mathcal{P} = p_H\left(R - \dfrac{B}{\Delta p}\right)\).

Because the lenders must break even in order to be willing to finance the project, a necessary
condition for the borrower to receive a loan is that the expected pledgeable income exceed the
lenders' initial outlay:

\[\mathcal{P} \equiv p_H\left(R - \frac{B}{\Delta p}\right) \geqslant I - A, \qquad (\text{IR}_l)\]

where ``IR\(_l\)'' stands for the lenders' individual rationality constraint (which we will also
often call the ``breakeven constraint'' or the ``participation constraint''). Thus a necessary
condition for financing to be arranged is

\[A \geqslant \bar{A} = p_H \frac{B}{\Delta p} - (p_H R - I). \qquad (3.3)\]

{[}\ldots{]} Conversely, if \(A \geqslant \bar{A}\), the entrepreneur is able to secure financing,
and so condition (3.3) is both a necessary and a sufficient condition for financing.

\emph{From Tirole §6.2 (privately-known-prospects model) and §6.2.1.2 (Asymmetric Information),
pp.~241--243 (PDF 254--256) --- the market-breakdown illustration the statement refers to:}

\emph{Privately-known-prospects model.} A borrower/entrepreneur has no funds (\(A = 0\)) to finance
a project costing \(I\). The project yields \(R\) in the case of success and 0 in the case of
failure. The borrower and the lenders are risk neutral, and the borrower is protected by limited
liability. The interest rate in the economy is normalized at 0.

The borrower can be one of two types. A good borrower has a probability of success equal to \(p\). A
bad borrower has a probability of success \(q\). Assume that \(p > q\) and that \(pR > I\) (at least
the good type is creditworthy). There are two subcases, which we will treat separately: either
\(pR > I > qR\) (only the good type is creditworthy), or \(pR > qR > I\) (both types are
creditworthy).

The borrower has private information about her type. The capital market, which is competitive and
demands an expected rate of return equal to 0, puts probabilities \(\alpha\) and \(1 - \alpha\) on
the borrower being a good or a bad type, respectively. Under asymmetric information, the capital
market does not know whether it faces a ``\(p\)-borrower'' (a good borrower) or a ``\(q\)-borrower''
(a bad borrower). Let

\[m \equiv \alpha p + (1 - \alpha) q\]

denote the investors' prior probability of success. {[}\ldots{]}

\textbf{6.2.1.2 Asymmetric Information.} {[}\ldots{]} Let us assume that the only feasible financial
contracts are contracts that give the borrower a compensation \(R_b \geqslant 0\) in the case of
success and 0 in the case of failure. (The validity of this assumption will be discussed in the
remark below on the optimality of contracts.) Such contracts necessarily pool the two types of
borrower as each prefers receiving financing to not being funded, and conditional on being funded,
prefers contracts with a higher compensation. The investors' profit for such a contract is therefore
on average:

\[[\alpha p + (1 - \alpha) q](R - R_b) - I = m(R - R_b) - I.\]

\emph{No lending: \(mR < I\).} This case can arise only if the bad borrower is not creditworthy. It
then arises whenever the probability that the borrower is a bad borrower is large enough, or
\(\alpha < \alpha^*\), where

\[\alpha^*(pR - I) + (1 - \alpha^*)(qR - I) = 0.\]

Because the borrower cannot receive a negative compensation (\(R_b \geqslant 0\)), investors lose
money if they choose to finance the project. Accordingly they do not and the market breaks down.

The good borrower is therefore hurt by the suspicion that she might be a bad one. There is
underinvestment.

\emph{Lending: \(mR \geqslant I\).} This case corresponds either to the situation in which both
types are creditworthy or to that in which the bad borrower is not creditworthy but
\(\alpha \geqslant \alpha^*\). The borrower's compensation \(R_b\) is then set so that investors
break even on average:

\[m(R - R_b) = I.\]

This implies that, ex post, investors make money on the good type (\(p(R - R_b) > I\)) and lose
money on the bad type (\(q(R - R_b) < I\)): there is cross-subsidization.
\stmtend
%% END-STATEMENT

\asked (i) Explain why lenders cannot fund only the good (low-$B$) borrower, and work out what happens to lenders' profit for a success reward $R_b$ below $B_L/\Delta p$, above $B_H/\Delta p$, and in between. (ii) Show that there is a cutoff $\alpha^*\in(0,1)$ for the share of good borrowers: no funding below it, funding is an equilibrium at or above it. (iii) Describe who subsidises whom when funding occurs.

\approach This is the lemons/pooling model of lecture block 3.2, with the type defined by the private benefit (moral hazard parameter) instead of the success probability. The borrower's incentive constraint (IC$_b$), $(\Delta p)R_b\ge B$, decides which types work at each $R_b$; investors' break-even then depends on the pool's success probability. Because the bad type can always copy the good one, any funded contract pools both types, so the analysis reduces to one function, lenders' pooled profit, and its sign as $\alpha$ varies (a market-breakdown threshold, as in Tirole \S6.2.1).

\notation
\begin{itemize}
  \item $I$: investment; $A=0$; $R>0$: income on success (0 on failure); $p_H>p_L$: success probability when working/shirking; $\Delta p\equiv p_H-p_L$.
  \item $B\in\{B_L,B_H\}$, $0<B_L<B_H$: private benefit from shirking (good/bad borrower); $\alpha\in(0,1)$: prior that $B=B_L$.
  \item $R_b\in[0,R]$: borrower's reward on success (0 on failure); $R-R_b$: lenders' claim.
  \item $m\equiv\alpha p_H+(1-\alpha)p_L$: success probability of the pool when the good type works and the bad type shirks.
  \item $\alpha^*$: the funding threshold; $R_b^*$: reward under competitive pooled lending.
\end{itemize}

\step{1}{Benchmark and why only good types cannot be financed} With reward $R_b$ in success and 0 in failure, type $B$ gets $p_HR_b$ by working and $p_LR_b+B$ by shirking, so she works iff
\begin{align}
  p_HR_b \ge p_LR_b+B \iff (p_H-p_L)R_b\ge B \iff (\Delta p)R_b\ge B. \tag{IC$_b$}
\end{align}
Under symmetric information lenders can be repaid at most $R-B/\Delta p$ in success without breaking (IC$_b$), so type $B$ is funded iff its pledgeable income covers $I$, $p_H(R-B/\Delta p)\ge I$; by (1) the good type is and the bad type is not. Under asymmetric information a bad borrower can offer any contract the good one offers. If she offered a different contract she would reveal her type and, by (1), not be funded; any funded contract gives her at least the good type's payoff (same if she works, more if she shirks, since $\max\{p_HR_b,p_LR_b+B_H\}\ge\max\{p_HR_b,p_LR_b+B_L\}$). So any contract that funds the good type also attracts the bad type: lenders cannot fund good types alone.

\step{2}{(i) The three ranges of $R_b$} For each range, (IC$_b$) tells us which types work; lenders' profit follows.

\textit{(a) $R_b<B_L/\Delta p$.} Then $(\Delta p)R_b<B_L<B_H$, so both types shirk and succeed with $p_L$. Since $R_b\ge0$, and by (2),
\begin{align}
  p_L(R-R_b)-I \le p_L R - I < 0. \label{t3e6-a}
\end{align}
\textit{(b) $R_b\ge B_H/\Delta p$.} Then $(\Delta p)R_b\ge B_H>B_L$, so both types work and succeed with $p_H$: the claim is information-insensitive. Since $R-R_b\le R-B_H/\Delta p$, and by the left inequality of (1),
\begin{align}
  p_H(R-R_b)-I \le p_H\left(R-\frac{B_H}{\Delta p}\right)-I < 0. \label{t3e6-b}
\end{align}
\textit{(c) $B_L/\Delta p\le R_b<B_H/\Delta p$.} Then $(\Delta p)R_b\ge B_L$ but $<B_H$: the good type works and the bad type shirks. The pool succeeds with probability $\alpha p_H+(1-\alpha)p_L=m$, so lenders' expected profit is
\begin{align}
  \Pi(R_b,\alpha) = \alpha p_H(R-R_b)+(1-\alpha)p_L(R-R_b) - I = m(R-R_b) - I. \label{t3e6-c}
\end{align}
By \eqref{t3e6-a}--\eqref{t3e6-b} only range (c) can attract funds, and by Step~1 any equilibrium is \emph{pooling}: both types or neither are funded, and when funded the good type works and the bad type shirks.

\answer{(i)} Lenders cannot separate the types (the bad type copies any funded contract). Rewards below $B_L/\Delta p$ make both types shirk and rewards from $B_H/\Delta p$ up make both work, and in both ranges lenders lose money, by (2) and (1) respectively. Only $R_b\in[B_L/\Delta p,\,B_H/\Delta p)$, where the good type works and the bad type shirks, can break even, with pooled profit $m(R-R_b)-I$.

\step{3}{(ii) The threshold $\alpha^*$} In range (c), \eqref{t3e6-c} falls with $R_b$, so it is highest at the lowest reward in the range, $R_b=B_L/\Delta p$. Define
\begin{align}
  \Pi(\alpha) \equiv [\alpha p_H+(1-\alpha)p_L]\left(R-\frac{B_L}{\Delta p}\right) - I
   = [p_L+\alpha\,\Delta p]\left(R-\frac{B_L}{\Delta p}\right) - I. \label{t3e6-Pi}
\end{align}
Differentiating the second form, $\Pi'(\alpha)=\Delta p\,(R-B_L/\Delta p)$, which is positive because $p_H(R-B_L/\Delta p)>I>0$ by (1). At the end points,
\begin{align}
  \Pi(0) &= p_L\left(R-\frac{B_L}{\Delta p}\right)-I < p_LR-I < 0 \quad\text{by (2)}, \nonumber\\
  \Pi(1) &= p_H\left(R-\frac{B_L}{\Delta p}\right)-I > 0 \quad\text{by (1)}. \nonumber
\end{align}
A continuous, strictly increasing function that is negative at 0 and positive at 1 has a unique root $\alpha^*\in(0,1)$. Setting \eqref{t3e6-Pi} to zero and solving for $\alpha$:
\begin{align}
  [p_L+\alpha^*\Delta p]\left(R-\frac{B_L}{\Delta p}\right) &= I \nonumber\\
  p_L+\alpha^*\Delta p &= \frac{I}{R-B_L/\Delta p} \nonumber\\
  \alpha^* &= \frac{1}{\Delta p}\left[\frac{I}{R-B_L/\Delta p}-p_L\right]. \label{t3e6-astar}
\end{align}
\textit{No financing if $\alpha<\alpha^*$:} then $\Pi(\alpha)<0$, and every $R_b$ loses money (ranges (a), (b) by \eqref{t3e6-a}--\eqref{t3e6-b}, range (c) because $\Pi(R_b,\alpha)\le\Pi(\alpha)$). The market breaks down.

\textit{Financing if $\alpha\ge\alpha^*$:} competitive lenders break even at
\begin{align}
  m(R-R_b^*) &= I \nonumber\\
  R_b^* &= R-\frac{I}{m}. \label{t3e6-Rbstar}
\end{align}
It lies in range (c). Lower end: $R_b^*\ge B_L/\Delta p$ iff $R-I/m\ge B_L/\Delta p$ iff $m(R-B_L/\Delta p)\ge I$, i.e.\ iff $\Pi(\alpha)\ge0$, i.e.\ iff $\alpha\ge\alpha^*$. Upper end: since $m<p_H$, $I/m>I/p_H$, so $R_b^*=R-I/m<R-I/p_H$; and the left inequality of (1), $p_H(R-B_H/\Delta p)<I$, divided by $p_H$ and rearranged, reads $R-I/p_H<B_H/\Delta p$. Together, $R_b^*<B_H/\Delta p$. So the good type works, lenders break even on average, and both types offering $\{R_b^*,0\}$ is an equilibrium.

\answer{(ii)} $\alpha^*=\dfrac{1}{\Delta p}\Bigl[\dfrac{I}{R-B_L/\Delta p}-p_L\Bigr]\in(0,1)$. For $\alpha<\alpha^*$ every contract loses money and the market breaks down; for $\alpha\ge\alpha^*$ the pooling contract $R_b^*=R-I/[\alpha p_H+(1-\alpha)p_L]$ (good type works, bad type shirks) breaks even and is an equilibrium.

\step{4}{(iii) Cross-subsidies} Substituting $R-R_b^*=I/m$ from \eqref{t3e6-Rbstar}, the lenders' profit on each type is
\begin{align}
  \text{good: } p_H\left(\frac{I}{m}\right)-I &= \frac{p_HI-mI}{m} = \frac{I(p_H-m)}{m} > 0, \label{t3e6-good}\\
  \text{bad: } p_L\left(\frac{I}{m}\right)-I &= \frac{p_LI-mI}{m} = -\frac{I(m-p_L)}{m} < 0, \label{t3e6-bad}
\end{align}
and $\alpha\cdot\eqref{t3e6-good}+(1-\alpha)\cdot\eqref{t3e6-bad}=I[\alpha p_H+(1-\alpha)p_L-m]/m=0$. Lenders make money on the good type and lose on the bad type. The good type's symmetric-information reward is $R-I/p_H$ (lenders break even on her alone), so she now pays
\begin{align}
  p_H\Bigl[\Bigl(R-\frac{I}{p_H}\Bigr)-R_b^*\Bigr] = p_H\Bigl[\frac{I}{m}-\frac{I}{p_H}\Bigr] = \frac{I(p_H-m)}{m} \nonumber
\end{align}
more in expectation, exactly the lenders' profit on her. The bad type, who would not be funded at all under symmetric information, gets $p_LR_b^*+B_H>0$. There is \emph{more} lending than under symmetric information (the bad project is funded) and the quality of lending falls: NPV conditional on funding is $mR-I$ instead of $p_HR-I$.

\answer{(iii)} Good borrowers cross-subsidise bad ones: lenders earn $I(p_H-m)/m>0$ on each good borrower and lose $I(m-p_L)/m$ on each bad one, netting to zero. The bad (negative-NPV) project gets funded and the good borrower's reward falls below its symmetric-information level.

\step{5}{Numerical illustration} Take $p_H=0.8$, $p_L=0.2$ ($\Delta p=0.6$), $R=100$, $I=50$, $B_L=9$, $B_H=27$. Then
\begin{align}
  \frac{B_L}{\Delta p} = \frac{(9)}{(0.6)} = 15, \qquad \frac{B_H}{\Delta p} = \frac{(27)}{(0.6)} = 45; \nonumber
\end{align}
(1) reads $0.8(100-45)=44<50<0.8(100-15)=68$ and (2) reads $p_LR=(0.2)(100)=20<50$. From \eqref{t3e6-astar},
\begin{align}
  \alpha^* = \frac{1}{(0.6)}\left[\frac{(50)}{(100)-(15)}-(0.2)\right] = \frac{1}{0.6}\Bigl[\frac{10}{17}-\frac{1}{5}\Bigr] = \frac{1}{0.6}\cdot\frac{33}{85} = \frac{11}{17} = 0.6471. \nonumber
\end{align}
At $\alpha=0.8$:
\begin{align}
  m &= (0.8)(0.8)+(0.2)(0.2) = 0.64+0.04 = 0.68, \nonumber\\
  R_b^* &= (100)-\frac{(50)}{(0.68)} = 100-73.53 = 26.47\in[15,45), \nonumber\\
  \text{good: } & (0.8)(73.53)-50 = 58.82-50 = 8.82, \nonumber\\
  \text{bad: } & (0.2)(73.53)-50 = 14.71-50 = -35.29, \nonumber
\end{align}
which average to zero ($0.8\times8.82=0.2\times35.29=7.06$). The good type gets $0.8\times26.47=21.18$ instead of $p_HR-I=30$ under symmetric information; the bad type gets $0.2\times26.47+27=32.29$ instead of $0$.

\intuition The lender's claim is information-sensitive only in the range where it can be profitable, and there the bad type always pools. Contracts cannot screen, so the pool's quality $m$ decides everything: with few good borrowers the market breaks down (underinvestment); with many, everyone is funded at a rate that taxes good borrowers to subsidise bad ones (overinvestment). This is the lecture's $m$, $\alpha^*$, cross-subsidy logic, with types defined by the private benefit instead of the success probability.

\pitfalls
\begin{itemize}
  \item Pricing the pooled loan with $p_H$ for both types: at $R_b\in[B_L/\Delta p,B_H/\Delta p)$ the bad type \emph{shirks}, so the pool succeeds with $\alpha p_H+(1-\alpha)p_L$.
  \item Trying to screen with a menu: whatever reward funds the good type is weakly preferred by the bad type, so there is no separating contract.
  \item Deriving $\alpha^*$ at an arbitrary $R_b$: the threshold uses the most favourable contract for lenders in range (c), $R_b=B_L/\Delta p$.
\end{itemize}
\bookans Matches the book's answer (pp.~595--596): the same three cases, the same $\alpha^*$ defined by $[\alpha^*p_H+(1-\alpha^*)p_L](R-B_L/\Delta p)=I$, the pooling contract $R_b^*=R-I/[\alpha p_H+(1-\alpha)p_L]$, and the cross-subsidy and fall in lending quality.
\refbib{Tirole \S3.2.1--3.2.2 pp.~115--117 (fixed-investment model); \S6.2.1 ``Market Breakdown and Cross-Subsidization'', pp.~241--244; answer pp.~595--596.}

% =============================================================================
\exsec{7. Rock's winner's curse model --- Amaro de Matos 6.2.7 \#4 and \#6, p.~175}
%% BEGIN-STATEMENT 7
\stmtbegin
\textbf{Source:} Amaro de Matos, Ch. 6, Exercises 6.2.7 \#4 and \#6, p.~175 (PDF 193) · Answer in
book: yes, p.~269 (\#6 solved; \#4 only refers to Rock 1986, p.~192)

\textbf{\#4.} Show that the probability of an uninformed investor receiving an allocation of an
underpriced issue (\(\pi\)) is less than or equal to the probability of receiving an allocation of
an overpriced issue (\(\pi'\)).

\textbf{\#6.} When there are informed and uninformed investors, it is said that as the proportion of
informed investors (\(\alpha\)) rises, the offer price (\(p_0\)) must fall to maintain the
equilibrium. Explain why.

\paragraph{Context needed to solve it}

\emph{From Amaro de Matos §6.2.2 ``Asymmetry between Investors'', p.~163 (PDF 181):}

\begin{quote}
The uninformed investors do not have a criterion to decide if an IPO is underpriced or not, so they
will go either for all of them or for none. The informed investors, on the other hand, will go only
for the underpriced issues, since they can distinguish between the two types.
\end{quote}

\begin{quote}
Suppose that the cost of information is \(c\) per investor, the price per share in the public offer
is \(p^\circ\), and the true value per share is a random variable \(v\), to which is associated a
density of probability \(f(v)\). Let \(n\) be the total number of shares offered, \(N\) be the total
number of informed investors, and \(\alpha\) be the fraction of shares allocated to them.
\end{quote}

(The book writes \(p^\circ\) in §6.2.2 and \(p_0\) in the exercises; they are the same offer price.)

\emph{From Amaro de Matos §6.2.2, p.~164 (PDF 182), equations typeset from the book:}

\begin{quote}
Clearly, if \(p^\circ > v\), there will be only uninformed investors interested, whereas if
\(p^\circ < v\), both kinds of investors will appear.
\end{quote}

Break-even of the uninformed, eq. (6.7):
\[\int_0^{p^\circ} n(p^\circ - v)f(v)\,dv = (1-\alpha)\int_{p^\circ}^{\infty} n(v - p^\circ)f(v)\,dv\]

Break-even of the informed, eq. (6.8): \[Nc = \alpha\int_{p^\circ}^{\infty} n(v - p^\circ)f(v)\,dv\]

which together give (6.9): \(E(v) - p^\circ = \dfrac{Nc}{n} > 0\).

\begin{quote}
If the demand is great enough---or for low enough prices---there may be an oversubscription, in
which case the underwriter rations the available shares.
\end{quote}

\emph{From Amaro de Matos §6.2.3, p.~166 (PDF 184):}

\begin{quote}
{[}\ldots{]} and it follows as before that as the proportion of informed investors (\(\alpha\))
rises, the offer price \(p^\circ\) must fall in order to maintain the equilibrium.
\end{quote}

\emph{Note on \(\pi\), \(\pi'\):} the book uses these symbols only in exercise \#4 and does not
define them anywhere in §6.2. They are Rock's (1986) notation, which is not in \texttt{Livros/}, and
mean what the statement says: the probability that an uninformed investor's order is filled for an
underpriced issue (\(\pi\)) and for an overpriced issue (\(\pi'\)).
\stmtend
%% END-STATEMENT

\asked (\#4) Prove that an uninformed investor's order is filled with a probability $\pi$ on an underpriced issue that is no larger than the probability $\pi'$ on an overpriced one. (\#6) Explain why, in equilibrium, the offer price $p_0$ must fall when the fraction $\alpha$ of the shares that goes to informed investors rises.

\approach Both parts are Rock's (1986) IPO model as set out in Amaro de Matos \S6.2.2 (lecture block 3.10). For \#4, write the fill probability under pro-rata rationing as a function of total demand and note that informed investors add demand only to underpriced issues. For \#6, the equilibrium condition is the uninformed investors' break-even (6.7); treat it as an implicit equation $G(p_0,\alpha)=0$ and sign $dp_0/d\alpha$ with the implicit function theorem.

\notation
\begin{itemize}
  \item $n$: shares offered; $v\ge0$: true value per share, density $f$, c.d.f.\ $F$; $p_0$ ($=p^\circ$ in \S6.2.2): offer price.
  \item $D_U$, $D_I$: uninformed and informed demand (shares); informed demand only underpriced issues ($v>p_0$).
  \item $\pi$, $\pi'$: probability that an uninformed order is filled in an underpriced / overpriced issue.
  \item $\alpha\in[0,1)$: fraction of the shares of an underpriced issue that goes to informed investors; $N$, $c$: number of informed investors and cost of information each.
\end{itemize}

\assume Rationing as in Rock (Amaro's answer to \#5): when an issue is oversubscribed every order has the same chance of being filled, so the fill probability of an order is $\min\{1,\,n/D\}$, where $D$ is total demand.

\step{1}{(\#4) $\pi\le\pi'$} An overpriced issue attracts only uninformed demand; an underpriced one attracts both:
\begin{align}
  \pi' &= \min\left\{1,\frac{n}{D_U}\right\}, \label{t3e7-pip}\\
  \pi  &= \min\left\{1,\frac{n}{D_U+D_I}\right\}. \label{t3e7-pi}
\end{align}
The map $D\mapsto\min\{1,n/D\}$ is non-increasing and $D_U+D_I\ge D_U$, so
\begin{align}
  \pi \le \pi', \label{t3e7-ineq}
\end{align}
with equality only when even the underpriced issue is not oversubscribed ($D_U+D_I\le n$), and strict inequality whenever informed demand causes rationing. $\blacksquare$

\answer{(\#4)} Under pro-rata rationing the fill probability is $\min\{1,n/D\}$, non-increasing in total demand $D$; underpriced issues carry the extra informed demand, so $\pi=\min\{1,n/(D_U+D_I)\}\le\min\{1,n/D_U\}=\pi'$.

\step{2}{Link to (6.7)} The uninformed hold the whole overpriced issue but only a fraction $1-\alpha$ of an underpriced one. With uninformed orders of exactly $D_U=n$, \eqref{t3e7-pip} gives $\pi'=\min\{1,n/n\}=1$, and the uninformed get $\pi n=(1-\alpha)n$ shares of an underpriced issue, so $\pi=1-\alpha$: the weight $1-\alpha$ on the right of (6.7) is $\pi/\pi'$. In the Ross IPO question (Cap.~15, Quest\~ao~3), $\pi=1/2$ and $\pi'=1$.

\step{3}{(\#6) The uninformed break-even condition} Let $L(p_0)\equiv\int_0^{p_0}(p_0-v)f(v)\,dv$ (expected loss per share on overpriced issues) and $U(p_0)\equiv\int_{p_0}^{\infty}(v-p_0)f(v)\,dv$ (expected gain on underpriced issues). Dividing both sides of (6.7) by $n$ gives $L(p_0)=(1-\alpha)U(p_0)$, so equilibrium requires
\begin{align}
  G(p_0,\alpha) \equiv (1-\alpha)U(p_0) - L(p_0) = 0. \label{t3e7-G}
\end{align}
By Leibniz's rule (the integrand vanishes at the moving limit $v=p_0$, so only the derivative inside the integral is left),
\begin{align}
  U'(p_0) &= -(p_0-p_0)f(p_0)+\int_{p_0}^{\infty}(-1)f(v)\,dv = -[1-F(p_0)], \nonumber\\
  L'(p_0) &= (p_0-p_0)f(p_0)+\int_0^{p_0}(1)f(v)\,dv = F(p_0). \nonumber
\end{align}
Differentiating \eqref{t3e7-G}:
\begin{align}
  \frac{\partial G}{\partial\alpha} &= -U(p_0) < 0, \label{t3e7-Ga}\\
  \frac{\partial G}{\partial p_0} &= -(1-\alpha)[1-F(p_0)] - F(p_0) < 0. \label{t3e7-Gp}
\end{align}

\step{4}{Comparative statics} By the implicit function theorem applied to \eqref{t3e7-G}, using \eqref{t3e7-Ga} and \eqref{t3e7-Gp},
\begin{align}
  \frac{dp_0}{d\alpha} &= -\frac{\partial G/\partial\alpha}{\partial G/\partial p_0} \nonumber\\
   &= -\frac{(-U(p_0))}{(-(1-\alpha)[1-F(p_0)]-F(p_0))} \nonumber\\
   &= -\frac{U(p_0)}{F(p_0)+(1-\alpha)[1-F(p_0)]} < 0. \label{t3e7-dp}
\end{align}
Equivalently $\partial\alpha/\partial p_0<0$, the form of the book's answer. A higher $\alpha$ shrinks the uninformed share of the good issues, so at the old price \eqref{t3e7-G} is negative (they lose on average); only a lower $p_0$, which raises $U$ and lowers $L$, restores zero expected profit.

\step{5}{Numerical illustration} Let $v$ be uniform on $[0,2]$ ($f=1/2$, $E(v)=1$). Then
\begin{align}
  L(p_0) &= \int_0^{p_0}(p_0-v)\tfrac12\,dv = \tfrac12\Bigl[p_0v-\tfrac{v^2}{2}\Bigr]_0^{p_0} = \tfrac12\cdot\frac{p_0^2}{2} = \frac{p_0^2}{4}, \nonumber\\
  U(p_0) &= \int_{p_0}^{2}(v-p_0)\tfrac12\,dv = \tfrac12\cdot\frac{(2-p_0)^2}{2} = \frac{(2-p_0)^2}{4}. \nonumber
\end{align}
Substituting into \eqref{t3e7-G}, multiplying by 4 and taking positive square roots (both sides are nonnegative for $p_0\in[0,2]$):
\begin{align}
  (1-\alpha)(2-p_0)^2 &= p_0^2 \nonumber\\
  \sqrt{1-\alpha}\,(2-p_0) &= p_0 \nonumber\\
  2\sqrt{1-\alpha} &= p_0\bigl(1+\sqrt{1-\alpha}\bigr) \nonumber\\
  p_0 &= \frac{2\sqrt{1-\alpha}}{1+\sqrt{1-\alpha}}. \nonumber
\end{align}
$\alpha=0$: $p_0=2/2=1=E(v)$; $\alpha=0.36$: $\sqrt{0.64}=0.8$, $p_0=1.6/1.8=0.8889$; $\alpha=0.64$: $\sqrt{0.36}=0.6$, $p_0=1.2/1.6=0.75$. At $\alpha=0.36$: $U=(2-0.8889)^2/4=0.308642$, $F(p_0)=p_0/2=0.444444$ and $1-F=0.555556$, so \eqref{t3e7-dp} gives
\begin{align}
  \frac{dp_0}{d\alpha} = -\frac{0.308642}{0.444444+(0.64)(0.555556)} = -\frac{0.308642}{0.8} = -0.3858, \nonumber
\end{align}
and the underpricing $E(v)-p_0=1-0.8889=0.1111$ equals $\alpha U(p_0)=(0.36)(0.308642)=0.1111$, which by (6.8) is $Nc/n$, as (6.9) requires.

\answer{(\#6)} The uninformed get only $1-\alpha$ of each good issue but all of each bad one; a higher $\alpha$ makes their break-even (6.7) negative at the old price, and only a lower $p_0$ restores it: $dp_0/d\alpha=-U(p_0)/\{F(p_0)+(1-\alpha)[1-F(p_0)]\}<0$.

\intuition The informed crowd the uninformed out of exactly the issues worth having, so the more of them there are, the worse the uninformed's average allocation (the winner's curse). The uninformed are needed to absorb the issue, and they stay only if they break even, which needs a lower offer price. In the lecture's version, $P_0=(\hat V-\lambda\alpha V_G)/(1-\lambda\alpha)$ falls as the informed capital share $\lambda$ rises.

\pitfalls
\begin{itemize}
  \item Claiming $\pi<\pi'$ always: equality holds when even the underpriced issue is not oversubscribed.
  \item Arguing that a higher $\alpha$ lowers $p_0$ ``because informed investors want cheap shares'': informed investors take the price as given; the price falls to keep the \emph{uninformed} willing to participate.
  \item Forgetting that $p_0$ moves both integrals: a lower price raises the gain on good issues \emph{and} cuts the loss on bad ones.
\end{itemize}
\bookans \#6 matches the book's answer (p.~269), which argues verbally from (6.7) that $\partial\alpha/\partial p_0<0$; \eqref{t3e7-dp} makes the sign explicit. For \#4 the book only refers to Rock (1986), p.~192, and gives no derivation.
\refbib{Amaro de Matos Ch.~6, \S6.2.2 ``Asymmetry between Investors'', pp.~163--164, eqs.~(6.7)--(6.9); \S6.2.3, p.~166; Appendix C.6.2, p.~269; lecture Topic~3, ``IPO underpricing''.}

% =============================================================================
\exsec{8. Pooling in credit markets with a continuum of types --- Tirole 6.2, p.~274}
%% BEGIN-STATEMENT 8
\stmtbegin
\textbf{Source:} Tirole, Ch. 6, Exercise 6.2, p.~274 (PDF 287) · Answer in book: yes, p.~596

\textbf{Exercise 6.2 (more on pooling in credit markets).} Consider the model of Exercise 6.1, in
which the borrower has private information about her benefit of misbehaving, except that the
borrower's type is drawn from a continuous distribution instead of a binary one. We will also assume
that there is a monopoly lender, who makes a credit offer to the borrower. The borrower has no
equity (\(A = 0\)).

Only the borrower knows the private benefit \(B\) of misbehaving. The lender only knows that this
private benefit is drawn from an \emph{ex ante} cumulative distribution \(H(B)\) on an interval
\([0, \bar B]\) (so, \(H(0) = 0\), \(H(\bar B) = 1\)). (Alternatively, one can imagine that lenders
face a population of borrowers with characteristic \(B\) distributed according to distribution
\(H\), and are unable to tell different types of borrower apart in their credit analysis.) The
lender knows all other parameters. For a loan agreement specifying share \(R_b\) for the borrower in
the case of success, and 0 in the case of failure, show that the lender's expected profit is

\[U_l = H((\Delta p)R_b)\,p_H(R - R_b) + [1 - H((\Delta p)R_b)]\,p_L(R - R_b) - I.\]

Show that

\begin{itemize}
\tightlist
\item
  the proportion of ``high-quality borrowers'' (that is, of borrowers who behave) is endogenous and
  increases with \(R_b\);\textsuperscript{78}
\item
  adverse selection reduces the quality of lending (if lending occurs, which as we will see cannot
  be taken for granted);
\item
  there is an externality among different types of borrower, in that the low-quality types (\(B\)
  large) force the lender to charge an interest rate that generates strictly positive profit on
  high-quality types (those with small \(B\));
\item
  the credit market may ``break down,'' that is, it may be the case that no credit is extended at
  all even though the borrower may be creditworthy (that is, have a low private benefit). To
  illustrate this, suppose that \(p_L = 0\) and \(H\) is uniform (\(H(B) = B/\bar B\)). Show that if
\end{itemize}

\[\frac{p_H^2}{\bar B}\,\frac{R^2}{4} < I\]

(which is the case for \(\bar B\) large enough), no loan agreement can enable the lender to recoup
on average his investment.

\textsuperscript{78} In this model, the loan agreement attracts all types of borrowers if it
attracts any type willing to behave. It is easy to find variants of the model in which this is not
the case and an increase in \(R_b\) attracts higher-quality borrowers, where ``higher-quality''
refers to an \emph{ex ante} selection effect and not only to an \emph{ex post} behavior like here.

\paragraph{Context needed to solve it}

\textbf{Part 0.} \emph{From Tirole, Ch. 6, Exercise 6.1, p.~273 (PDF 286):}

\textbf{Exercise 6.1 (privately known private benefit and market breakdown).} Section 6.2
illustrated the possibility of market breakdown without the possibility of signaling. This exercise
supplies another illustration. Let us consider the fixed-investment model of Section 3.2 and assume
that \emph{only the borrower knows the private benefit associated with misbehavior}. When the
borrower has private information about this parameter, lenders are concerned that this private
benefit might be high and induce the borrower to misbehave. In the parlance of information
economics, the ``bad types'' are the types of borrower with high private benefit. We study the case
of two possible levels of private benefit (see Exercise 6.2 for an extension to a continuum of
possible types). The borrower wants to finance a fixed-size project costing \(I\), and, for
simplicity, has no equity (\(A = 0\)). The project yields \(R\) (success) or 0 (failure). The
probability of success is \(p_H\) or \(p_L\), depending on whether the borrower works or shirks,
with \(\Delta p \equiv p_H - p_L > 0\). There is no private benefit when working. The private
benefit \(B\) enjoyed by the borrower when shirking is either \(B_L > 0\) or \(B_H > B_L\). The
borrower will be labeled a ``good borrower'' when \(B = B_L\) and a ``bad borrower'' when
\(B = B_H\). At the date of contracting, the borrower knows the level of her private benefit, while
the capital market puts (common knowledge) probabilities \(\alpha\) that the borrower is a good
borrower and \(1 - \alpha\) that she is a bad borrower. All other parameters are common knowledge
between the borrower and the lenders.

To make things interesting, let us assume that under asymmetric information, the lenders are
uncertain about whether the project should be funded:

\[p_H\left(R - \frac{B_H}{\Delta p}\right) < I < p_H\left(R - \frac{B_L}{\Delta p}\right). \qquad (1)\]

Assume that investors cannot break even if the borrower shirks:

\[p_L R < I. \qquad (2)\]

\begin{enumerate}
\def\labelenumi{(\roman{enumi})}
\item
  Note that the investor cannot finance only good borrowers. Assume that the entrepreneur receives
  no reward in the case of failure (this is indeed optimal); consider the effect of rewards \(R_b\)
  in the case of success that are (a) smaller than \(B_L/\Delta p\), (b) larger than
  \(B_H/\Delta p\), (c) between these two values.
\item
  Show that there exists \(\alpha^*\), \(0 < \alpha^* < 1\), such that
\end{enumerate}

\begin{itemize}
\tightlist
\item
  no financing occurs if \(\alpha < \alpha^*\),
\item
  financing is an equilibrium if \(\alpha \geqslant \alpha^*\).
\end{itemize}

\begin{enumerate}
\def\labelenumi{(\roman{enumi})}
\setcounter{enumi}{2}
\tightlist
\item
  Describe the ``cross-subsidies'' between types that occur when borrowing is feasible.
\end{enumerate}

\emph{From Tirole §3.2.1 ``The Fixed-Investment Model'', pp.~115--116 (PDF 128--129):}

\emph{Preferences and the loan agreement.} Both the borrower and the potential lenders (or
``investors'') are risk neutral. For notational simplicity, there is no time preference; the rate of
return expected by investors (which is also the riskless rate, due to risk neutrality) is taken to
be 0. The borrower is protected by limited liability, and so her income cannot take negative values.

{[}\ldots{]} Let us turn to the loan contract. A contract first stipulates whether the project is
financed. If so, it further specifies how the profit is shared between the lenders and the borrower.
The borrower's limited liability will imply that both sides receive 0 in the case of failure
{[}\ldots{]}. In the case of success, the two parties share the profit \(R\); \(R_b\) goes to the
borrower and \(R_l\) to the lenders. To sum up, we posit an incentive scheme for the entrepreneur of
the following form: \(R_b\) in the case of success, 0 in the case of failure.

\emph{From Tirole §3.2.2 ``The Lenders' Credit Analysis'', p.~116 (PDF 129):}

{[}\ldots{]} Because she has stake \(R_b\) in the firm's income (she receives \(R_b\) in the case of
success and 0 in the case of failure), the borrower will therefore behave if the following
``incentive compatibility constraint'' is satisfied:

\[p_H R_b \geqslant p_L R_b + B \quad \text{or} \quad (\Delta p) R_b \geqslant B. \qquad (\text{IC}_b)\]
\stmtend
%% END-STATEMENT

\asked (0) Derive the lender's expected profit $U_l$ for a loan that pays the borrower $R_b$ in success. Then show: (1) the share of borrowers who behave is endogenous and rises with $R_b$; (2) adverse selection lowers the quality of lending; (3) high-$B$ types impose an externality, forcing a strictly positive profit on low-$B$ types; (4) with $p_L=0$ and $B$ uniform on $[0,\bar B]$, if $p_H^2R^2/(4\bar B)<I$ no loan lets the lender break even, although some borrowers are creditworthy.

\approach The model is Tirole's fixed-investment model (\S3.2) with the private benefit privately known (Part~0, Tirole Exercise~6.1), now drawn from a continuous distribution $H$, and a monopoly lender. A single contract $R_b$ applies to all types; (IC$_b$), $(\Delta p)R_b\ge B$, sorts them into behaving (low $B$) and shirking (high $B$). Everything follows from writing the lender's profit as a mixture of the two groups. The breakdown result is a bound: maximise expected repayment over $R_b$ and compare it with $I$ (lecture block 3.2, market breakdown).

\notation
\begin{itemize}
  \item $I$, $A=0$, $R$, $p_H>p_L$, $\Delta p$: as in Part~0 (Tirole Exercise~6.1): investment, borrower's cash, income in success, success probabilities when working/shirking, $\Delta p=p_H-p_L$.
  \item $B\in[0,\bar B]$: borrower's private benefit, c.d.f.\ $H$ ($H(0)=0$, $H(\bar B)=1$), density $h$ where it exists.
  \item $R_b\in[0,R]$: borrower's reward on success; the lender is a monopolist and receives $R-R_b$ on success.
  \item $U_l(R_b)$: lender's expected profit; $\mu(R_b)\equiv p_L+\Delta p\,H((\Delta p)R_b)$: success probability of the funded pool.
\end{itemize}

\assume As in Part~0, $p_LR<I$ (condition (2) of Exercise~6.1) and $p_HR>I$ (implied by its condition (1)): low-$B$ types are creditworthy. Ties ($B=(\Delta p)R_b$) behave; with a continuous $H$ they have zero mass.

\step{1}{The lender's profit} With $A=0$ every type accepts any $R_b\ge0$ (she gets $\max\{p_HR_b,\;p_LR_b+B\}\ge0$, her outside option). Type $B$ works iff $p_HR_b\ge p_LR_b+B$, i.e.\ by (IC$_b$) iff $B\le(\Delta p)R_b$. Those types form the mass $\Pr[B\le(\Delta p)R_b]=H((\Delta p)R_b)$, and their projects succeed with probability $p_H$; the rest, mass $1-H((\Delta p)R_b)$, shirk and succeed with $p_L$. The lender collects $R-R_b$ on success and lends $I$:
\begin{align}
  U_l = H((\Delta p)R_b)\,p_H(R-R_b) + [1-H((\Delta p)R_b)]\,p_L(R-R_b) - I. \label{t3e8-Ul}
\end{align}

\answer{(0)} $U_l=H((\Delta p)R_b)\,p_H(R-R_b)+[1-H((\Delta p)R_b)]\,p_L(R-R_b)-I$, as stated.

\step{2}{Endogenous quality rises with $R_b$} The share of borrowers who behave, $H((\Delta p)R_b)$, is a c.d.f.\ evaluated at an increasing function of $R_b$, so it is non-decreasing in $R_b$; where the density exists, by the chain rule
\begin{align}
  \frac{d}{dR_b}H((\Delta p)R_b) = (\Delta p)\,h((\Delta p)R_b) \ge 0, \nonumber
\end{align}
strictly where $h>0$. Note that every type is funded whatever $R_b$: the reward changes \emph{behaviour}, not who applies (Tirole's footnote 78).

\answer{(1)} The behaving share is $H((\Delta p)R_b)$, chosen implicitly through $R_b$ and increasing in it (slope $(\Delta p)h\ge0$).

\step{3}{Adverse selection lowers the quality of lending} The success probability of funded loans is $\mu(R_b)=p_L+\Delta p\,H((\Delta p)R_b)\le p_H$, strictly below $p_H$ unless $R_b\ge\bar B/\Delta p$. Under symmetric information the lender would offer each type $R_b=B/\Delta p$, fund only types with $p_H(R-B/\Delta p)\ge I$, and every funded borrower would behave, so funded loans succeed with $p_H$ and no loan is made to a shirker, whose loan loses money by (2). Under asymmetric information a single $R_b$ applies to all, and any $R_b<\bar B/\Delta p$ funds a mass $1-H>0$ of shirkers; avoiding them requires $R_b\ge\bar B/\Delta p$, whose profit $p_H(R-\bar B/\Delta p)-I$ is negative when $\bar B$ is large.

\answer{(2)} Funded loans succeed with $\mu(R_b)=p_L+\Delta p\,H((\Delta p)R_b)<p_H$ whenever some types shirk, against $p_H$ for every funded loan under symmetric information: shirkers (whose loans lose money) are financed alongside behaving types.

\step{4}{Externality across types} Splitting \eqref{t3e8-Ul} by type (write $I=HI+(1-H)I$):
\begin{align}
  U_l = H\,[p_H(R-R_b)-I] + (1-H)\,[p_L(R-R_b)-I]. \label{t3e8-split}
\end{align}
If the lender lends, $U_l\ge0$; moving the second term to the right of the inequality,
\begin{align}
  H\,[p_H(R-R_b)-I] \ge (1-H)\,[I-p_L(R-R_b)] > 0 \quad\text{when } H<1, \label{t3e8-ext}
\end{align}
because $p_L(R-R_b)\le p_LR<I$ by (2). So the lender earns a strictly positive profit on each high-quality (behaving) borrower. Equivalently, $U_l\ge0$ means $\mu(R_b)(R-R_b)\ge I$, so the repayment $R-R_b\ge I/\mu(R_b)>I/p_H$ exceeds the one that would break even on that borrower alone. The losses on the high-$B$ types are recovered from the low-$B$ types: this is the externality.

\answer{(3)} Whenever the lender lends and some types shirk, $H\,[p_H(R-R_b)-I]\ge(1-H)\,[I-p_L(R-R_b)]>0$: the low-$B$ types pay a repayment above their own break-even level ($R-R_b>I/p_H$) to cover the losses on high-$B$ types.

\step{5}{Market breakdown with $p_L=0$ and uniform $H$} Now $\Delta p=p_H-0=p_H$ and, for the uniform distribution, $H((\Delta p)R_b)=\min\{p_HR_b/\bar B,\,1\}\le p_HR_b/\bar B$. With $p_L=0$ the second term of \eqref{t3e8-Ul} vanishes. Substituting, for every $R_b\in[0,R]$:
\begin{align}
  U_l + I &= H(p_HR_b)\,p_H(R-R_b) \nonumber\\
          &\le \left(\frac{p_HR_b}{\bar B}\right)p_H(R-R_b) \nonumber\\
          &= \frac{p_H^2}{\bar B}\,R_b(R-R_b) \nonumber\\
          &\le \frac{p_H^2}{\bar B}\,\frac{R^2}{4}, \label{t3e8-bound}
\end{align}
the last line because $g(R_b)\equiv R_b(R-R_b)=RR_b-R_b^2$ has $g'(R_b)=R-2R_b=0$ at $R_b=R/2$ and $g''=-2<0$, so its maximum is $g(R/2)=(R/2)(R/2)=R^2/4$. Hence if $\dfrac{p_H^2}{\bar B}\dfrac{R^2}{4}<I$, then $U_l<0$ for every $R_b$: no loan agreement lets the lender recoup $I$ on average. The left side falls with $\bar B$, so the condition holds for $\bar B>p_H^2R^2/(4I)$. Yet the types with $B\le p_HR-I$ have pledgeable income $p_H(R-B/p_H)=p_HR-B\ge I$ and would be funded under symmetric information; they are a mass $(p_HR-I)/\bar B>0$. When $R/2\le\bar B/p_H$ the bound in \eqref{t3e8-bound} is attained at $R_b=R/2$, the book's optimum.

\answer{(4)} With $p_L=0$ and uniform $H$, expected repayment is at most $p_H^2R^2/(4\bar B)$ (at $R_b=R/2$), so if this is below $I$ every loan loses money and credit is not extended, even though the types $B\le p_HR-I$ (a positive mass) would be funded under symmetric information.

\step{6}{Numerical illustration} $p_H=0.8$, $p_L=0$, $R=100$, $I=50$, $\bar B=50$: the maximum expected repayment is
\begin{align}
  \frac{p_H^2}{\bar B}\,\frac{R^2}{4} = \frac{(0.64)(10{,}000)}{(4)(50)} = \frac{6{,}400}{200} = 32 < 50, \nonumber
\end{align}
reached at $R_b=R/2=50\le\bar B/p_H=50/0.8=62.5$. Breakdown occurs for any $\bar B>p_H^2R^2/(4I)=6{,}400/200=32$, although the types $B\le p_HR-I=0.8(100)-50=30$, a mass $30/50=0.6$ of borrowers, are creditworthy. At $R_b=50$, $H=\min\{0.8\cdot50/50,1\}=0.8$ and the split \eqref{t3e8-split} gives $U_l=(0.8)[(0.8)(50)-50]+(0.2)[(0)(50)-50]=(0.8)(-10)+(0.2)(-50)=-8-10=-18$, that is, $32-50$.

\intuition A single price for all types cannot keep everyone honest without giving away most of the income: raising $R_b$ makes more types behave but leaves less for the lender. With many high-$B$ types, the best trade-off still does not cover $I$, and the market shuts even for borrowers who would be good risks under symmetric information, the lecture's market breakdown with a continuum of types. When it does lend, the honest types pay for the shirkers.

\pitfalls
\begin{itemize}
  \item Thinking a higher $R_b$ attracts better applicants: here every type always applies; $R_b$ only changes who \emph{behaves} (footnote 78).
  \item Maximising $p_H(R-R_b)H$ without the cap $H\le1$: for $R_b>\bar B/p_H$ the share of behaving types stops rising, so the interior first-order condition $R_b=R/2$ is valid only when $R/2\le\bar B/p_H$.
  \item Reading ``breakdown'' as ``no creditworthy borrower'': the types with $B\le p_HR-I$ are creditworthy; they are shut out by the pool.
\end{itemize}
\bookans Matches the book's answer (pp.~596--597): the proportion $H(R_b\Delta p)$ is endogenous and rising in $R_b$, the externality, and, with $p_L=0$ and uniform $H$, $R_b=R/2$ and maximal pledgeable income $p_H^2R^2/(4\bar B)$, below $I$ for $\bar B$ large. Two small differences: the book's displayed $U_l$ omits the ``$-I$'' (it is the expected repayment), and its first-order condition covers only the interior case $R/2\le\bar B/p_H$; the bound \eqref{t3e8-bound} covers every $R_b$. The book gives no argument for the ``quality of lending'' bullet; Step~3 supplies one.
\refbib{Tirole \S3.2.2 p.~116 (IC$_b$); \S6.2.1 pp.~241--244; Exercise~6.1; answer pp.~596--597.}

% =============================================================================
% Topic 3 list, exercises 9--15 (detailed). Label prefix t3e<n>-.
% =============================================================================
\exsec{9. Adverse selection and ratings --- Tirole 6.6, p.~276}
%% BEGIN-STATEMENT 9
\stmtbegin
\textbf{Source:} Tirole, Ch. 6, Exercise 6.6, p.~276 (PDF 289) · Answer in book: yes, p.~598

A borrower has assets \emph{A} and must find financing for a fixed investment \emph{I}
\textgreater{} \emph{A}. As usual, the project yields \(R\) (success) or 0 (failure). The borrower
is protected by limited liability. The probability of success is \(p_H\) or \(p_L\), depending on
whether the borrower works or shirks, with \(\Delta p \equiv p_H - p_L > 0\). There is no private
benefit when working. The private benefit enjoyed by the borrower when shirking is either \(b\)
(with probability \(\alpha\)) or \(B\) (with probability \(1 - \alpha\)). At the date of
contracting, the borrower knows her private benefit, but the market (which is risk neutral and
charges a 0 average rate of interest) does not know it. Assume that \(p_L R + B < I\) (the project
is always inefficient if the borrower shirks) and that

\[p_H\left(R - \frac{B}{\Delta p}\right) < I - A < p_H\left(R - \frac{b}{\Delta p}\right) \qquad (1)\]

and

\[[\alpha p_H + (1 - \alpha)p_L]\left[R - \frac{b}{\Delta p}\right] < I - A. \qquad (2)\]

\begin{enumerate}
\def\labelenumi{(\roman{enumi})}
\item
  Interpret conditions (1) and (2) and show that there is no lending in equilibrium.
\item
  Suppose now that the borrower can at cost \(r(x) = rx\) (which is paid from the cash endowment
  \(A\)) purchase a signal with quality \(x \in [0, 1]\). (This quality can be interpreted as the
  reputation or the number of rating agencies that the borrower contracts with.) With probability
  \(x\), the signal reveals the borrower's type (\(b\) or \(B\)) perfectly; with probability
  \(1 - x\), the signal reveals nothing. The financial market observes both the quality \(x\) of the
  signal chosen by the borrower and the outcome of the signal (full or no information). The borrower
  then offers a contract that gives the borrower \(R_b\) and the lenders \(R - R_b\) in the case of
  success (so, a contract is the choice of an \(R_b \in [0, R]\)). The timing is summarized in
  Figure 6.5.
\end{enumerate}

\emph{Figure 6.5}

{\def\LTcaptype{none} % do not increment counter
\begin{longtable}[]{@{}
  >{\raggedright\arraybackslash}p{(\linewidth - 4\tabcolsep) * \real{0.3333}}
  >{\raggedright\arraybackslash}p{(\linewidth - 4\tabcolsep) * \real{0.3333}}
  >{\raggedright\arraybackslash}p{(\linewidth - 4\tabcolsep) * \real{0.3333}}@{}}
\toprule\noalign{}
\begin{minipage}[b]{\linewidth}\raggedright
Step 1
\end{minipage} & \begin{minipage}[b]{\linewidth}\raggedright
Step 2
\end{minipage} & \begin{minipage}[b]{\linewidth}\raggedright
Step 3
\end{minipage} \\
\midrule\noalign{}
\endhead
\bottomrule\noalign{}
\endlastfoot
Borrower chooses quality of signal \(x\) (this quality is observed by the capital market). &
Borrower's type revealed with probability \(x\). Nothing revealed with probability \(1 - x\). &
Borrower goes to the capital market. \\
\end{longtable}
}

Look for a pure strategy, \emph{separating} equilibrium, that is, an equilibrium in which the two
types pick different signal qualities.\textsuperscript{79}

\begin{itemize}
\tightlist
\item
  Argue that the bad borrower (borrower \(B\)) does not purchase a signal in a separating
  equilibrium.
\item
  Argue that the good borrower (borrower \(b\)) borrows under the same conditions regardless of the
  signal's realization, in a separating equilibrium.
\item
  Show that the good borrower chooses signal quality \(x \in (0, 1)\) given by
  \[A = x(A - rx) + (1 - x)\left[p_L\left(R - \frac{I - A + rx}{p_H}\right) + B\right].\]
\item
  Show that this separating equilibrium exists only if \(r\) is ``not too large.''
\end{itemize}

\begin{quote}
\begin{enumerate}
\def\labelenumi{\arabic{enumi}.}
\setcounter{enumi}{78}
\tightlist
\item
  One will assume that if the signal reveals the borrower's type, the investors put probability 1 on
  this type, even when they put weight 0 on the corresponding type after observing the quality of
  the signal.
\end{enumerate}
\end{quote}
\stmtend
%% END-STATEMENT

\asked (i) Say what conditions (1) and (2) mean economically, and prove that without ratings nobody lends. (ii) With a rating of quality $x$ (cost $rx$, reveals the type with probability $x$), look for a separating equilibrium and show: (ii-1) the bad type buys no rating; (ii-2) the good type borrows on the same terms whether or not the rating reveals her; (ii-3) her rating quality $\hat x\in(0,1)$ solves the displayed no-mimicking equation; (ii-4) such an equilibrium exists only if $r$ is not too large.

\approach Part (i) is the pooling/market-breakdown logic of lecture block 3.2 with types defined by the private benefit: (IC$_b$) sorts types into workers and shirkers at each $R_b$, the bad type can copy any funded contract, and condition (2) makes the pool unprofitable. Part (ii) is costly signalling through certification (Tirole \S6.3, Application 4; lecture block 3.7). In a separating equilibrium the good type buys the cheapest rating quality at which the bad type prefers her sure outside option to gambling that the rating reveals nothing; that no-mimicking constraint binds and pins down $\hat x$. Existence then requires checking that the good type is still funded, still works, can pay the fee and does not prefer to deviate (with pessimistic off-path beliefs).

\notation
\begin{itemize}
  \item $A>0$: borrower's cash; $I>A$: fixed investment; $R>0$: income in success, $0$ in failure; market rate $0$.
  \item $p_H>p_L\ge0$: success probabilities when working / shirking; $\Delta p\equiv p_H-p_L>0$.
  \item $b<B$: private benefits from shirking of the good type (probability $\alpha\in(0,1)$) and of the bad type (probability $1-\alpha$).
  \item $R_b\in[0,R]$: borrower's reward in success (0 in failure); $R-R_b$: investors' claim.
  \item $x\in[0,1]$: signal quality, cost $rx$ paid out of $A$, $r>0$; $\hat x$: the good type's equilibrium choice.
  \item $R_b^G(x)$: the good type's reward when recognised as good with net cash $A-rx$; $U_G$: her payoff.
  \item $c(x)$: bad type's payoff when treated as good with signal quality $x$; $c_0\equiv c(0)$.
  \item $\kappa\equiv p_L/p_H\in[0,1)$; $y\equiv r\hat x$: total signalling cost.
\end{itemize}

\assume (a) The good borrower's project has positive NPV when she works, $p_HR>I$ (otherwise there is
nothing to signal). (b) Following footnote 79, a revealing signal fixes the investors' belief; at a
signal quality not chosen by anyone in equilibrium, and with no revelation, investors believe they face
type $B$ (pessimistic off-path beliefs). (c) In part (i), contracts also pay $0$ in failure, as in part (ii);
a payment in failure would only weaken incentives and reduce what investors can be paid.

\step{1}{Interpreting (1) and (2)} A borrower with private benefit $\beta\in\{b,B\}$ works if and only if
$(\Delta p)R_b\ge\beta$, so the most she can promise investors while working is the pledgeable income
\begin{align}
  \mathcal{P}(\beta) = p_H\Bigl(R-\frac{\beta}{\Delta p}\Bigr)
     = p_HR - p_H\frac{\beta}{\Delta p}. \label{t3e9-pledge}
\end{align}
Condition (1) of the statement reads $\mathcal{P}(B)<I-A<\mathcal{P}(b)$: under symmetric information the
good type could raise $I-A$ and the bad type could not. Condition (2) of the statement says that a
contract that gives the minimum incentive stake of the good type, $R_b=b/\Delta p$, and is taken by both
types (the good one works, the bad one shirks) returns to investors, at the prior, less than $I-A$.
The assumption $p_LR+B<I$ says shirking destroys value even counting the private benefit.

\step{2}{The bad type always wants what the good type wants} Under a contract $R_b$, type $\beta$
chooses the better of working and shirking:
\begin{align}
  u_\beta(R_b) = \max\{p_HR_b,\; p_LR_b+\beta\}. \label{t3e9-u}
\end{align}
Since $B>b$, \eqref{t3e9-u} gives $u_B(R_b)\ge u_b(R_b)$ for every $R_b$, and both types' outside option is
to keep $A$. Hence any contract the good type accepts is also accepted by the bad type.

\step{3}{No lending in equilibrium (part (i))} Take any $R_b$ and split by range.

\emph{(a) $R_b\ge B/\Delta p$: both types work.} Investors receive
\begin{align}
  p_H(R-R_b) \le p_H\Bigl(R-\frac{B}{\Delta p}\Bigr)=\mathcal{P}(B) < I-A, \nonumber
\end{align}
by \eqref{t3e9-pledge} and condition (1): a loss whatever the investors' beliefs.

\emph{(b) $R_b<b/\Delta p$: both types shirk.} Investors receive $p_L(R-R_b)$ from either type. If they
break even, $p_L(R-R_b)\ge I-A$, then type $\beta$ obtains
\begin{align}
  p_LR_b+\beta &\le p_LR-(I-A)+\beta \nonumber\\
  &\le (p_LR+B-I)+A \nonumber\\
  &< A && \text{(using } p_LR+B<I), \nonumber
\end{align}
so neither type accepts: there is no lending in this range.

\emph{(c) $b/\Delta p\le R_b<B/\Delta p$: the good type works, the bad type shirks.} Because
$(\Delta p)R_b<B$, \eqref{t3e9-u} gives $u_B(R_b)=p_LR_b+B>p_HR_b=u_b(R_b)$. If the good type is funded at
$R_b$ she needs $p_HR_b\ge A$, so the bad type gets $u_B(R_b)>A$ there. By (a) and (b) (and the same
computation as (b) for a recognised bad type in this range) she cannot get more than $A$ from any contract
on which she is recognised. So she offers $R_b$ too, the investors' posterior on $R_b$ equals the prior,
and their revenue is
\begin{align}
  [\alpha p_H+(1-\alpha)p_L](R-R_b)
  \le [\alpha p_H+(1-\alpha)p_L]\Bigl(R-\frac{b}{\Delta p}\Bigr) < I-A, \nonumber
\end{align}
by condition (2). Every range yields a loss or a refusal: no lending occurs in equilibrium.
$\blacksquare$

\answer{(i)} (1): with symmetric information the good type could raise $I-A$ and the bad type could not. (2): the cheapest contract that makes the good type work, taken by both types at the prior (the bad type then shirks), does not repay $I-A$. Since the bad type copies any contract the good type is funded on, every contract either loses money or is refused: no lending.

\step{4}{The bad type buys no signal (part (ii), first bullet)} In a separating equilibrium the bad type's
quality $x_B\ne\hat x$ identifies her whatever the signal shows. With net cash $A-rx_B$ she needs
$I-A+rx_B\ge I-A$, so cases (a) and (b) of Step~3 apply \emph{a fortiori}: she is not funded on terms she
accepts and keeps $A-rx_B$. This is maximised at $x_B=0$:
\begin{align}
  U_B = A. \label{t3e9-UB}
\end{align}

\answer{(ii-1)} In a separating equilibrium the bad type is identified by her choice of $x$, is not funded on terms she accepts, and so any fee is wasted: she picks $x_B=0$ and gets $A$.

\step{5}{The good type borrows on the same terms in both events (second bullet)} At $\hat x$ the
investors infer type $b$ from the choice of $x$ alone; a revealing signal shows $b$, which is the same
belief. With the same belief the competitive market accepts the same contract. The good type has net
cash $A-r\hat x$, borrows $I-A+r\hat x$, and offers the smallest stake that investors accept:
\begin{align}
  p_H(R-R_b^G) &= I-A+rx, \label{t3e9-be}\\
  R-R_b^G &= \frac{I-A+rx}{p_H}, \nonumber\\
  R_b^G(x) &= R-\frac{I-A+rx}{p_H}. \label{t3e9-RbG}
\end{align}
She works if $(\Delta p)R_b^G\ge b$; substituting \eqref{t3e9-RbG}, dividing by $\Delta p$ and multiplying by $p_H$,
\begin{align}
  (\Delta p)\Bigl(R-\frac{I-A+rx}{p_H}\Bigr) \ge b
  \iff R-\frac{b}{\Delta p}\ge\frac{I-A+rx}{p_H}
  \iff I-A+rx \le p_H\Bigl(R-\frac{b}{\Delta p}\Bigr). \label{t3e9-icG}
\end{align}
Her payoff, substituting \eqref{t3e9-RbG}, is
\begin{align}
  U_G(x) = p_H\Bigl(R-\frac{I-A+rx}{p_H}\Bigr) = p_HR-(I-A+rx) = p_HR-I+A-rx, \label{t3e9-UG}
\end{align}
decreasing in $x$: she wants the cheapest signal that separates.

\answer{(ii-2)} At $\hat x$ the investors already infer type $b$ from the choice of $x$, and a revealing signal shows the same type, so in both events she borrows on $R_b^G=R-(I-A+r\hat x)/p_H$.

\step{6}{The mimicking constraint pins down $\hat x\in(0,1)$ (third bullet)} If the bad type chooses
$x$, the signal reveals her with probability $x$ (no funding, she keeps $A-rx$); with probability $1-x$
nothing is revealed and she is treated as good, receiving \eqref{t3e9-RbG}. She then shirks, because
\begin{align}
  (\Delta p)R_b^G < B \iff p_H\Bigl(R-\frac{B}{\Delta p}\Bigr) < I-A+rx, \nonumber
\end{align}
which holds by condition (1) and $rx\ge0$. Her payoff in that event is
\begin{align}
  c(x) = p_LR_b^G(x)+B = p_L\Bigl(R-\frac{I-A+rx}{p_H}\Bigr)+B
       = \underbrace{p_L\Bigl(R-\frac{I-A}{p_H}\Bigr)+B}_{c_0}-\frac{p_L}{p_H}\,rx = c_0-\kappa rx. \label{t3e9-c}
\end{align} Her mimicking payoff is $x(A-rx)+(1-x)c(x)$ (if $c(x)<A-rx$ she would
decline the loan and get $A-rx<A$, so the condition below covers that case too). She does not mimic iff
\begin{align}
  A \ge x(A-rx)+(1-x)\Bigl[p_L\Bigl(R-\frac{I-A+rx}{p_H}\Bigr)+B\Bigr]. \label{t3e9-nm}
\end{align}
Let $\Phi(x)$ be the right side of \eqref{t3e9-nm} minus $A$. Substituting \eqref{t3e9-c} and expanding,
\begin{align}
  \Phi(x) &= xA-rx^2+(1-x)(c_0-\kappa rx)-A \nonumber\\
          &= xA-rx^2+c_0-\kappa rx-xc_0+\kappa rx^2-A \nonumber\\
          &= -r(1-\kappa)x^2+(A-c_0-\kappa r)x+(c_0-A), \label{t3e9-Phi}
\end{align}
a strictly concave quadratic ($\kappa<1$). At $x=0$, the bad type shirks under $R_b^G(0)$, so
\begin{align}
  c_0 = p_LR_b^G(0)+B > p_HR_b^G(0) = p_HR-I+A > A, \nonumber
\end{align}
using \eqref{t3e9-UG} at $x=0$ and assumption (a): $\Phi(0)=c_0-A>0$. At $x=1$, \eqref{t3e9-Phi} gives
$\Phi(1)=-r(1-\kappa)+A-c_0-\kappa r+c_0-A=-r<0$. A concave function positive at 0 and negative at 1 has a
single root $\hat x\in(0,1)$, with $\Phi>0$ on $[0,\hat x)$ and $\Phi<0$ on $(\hat x,1]$. By
\eqref{t3e9-UG} the good type picks the smallest admissible $x$, i.e.\ $\hat x$, where \eqref{t3e9-nm} binds:
\begin{align}
  A = \hat x(A-r\hat x)+(1-\hat x)\Bigl[p_L\Bigl(R-\frac{I-A+r\hat x}{p_H}\Bigr)+B\Bigr]. \qquad\blacksquare \label{t3e9-xhat}
\end{align}

\answer{(ii-3)} $\hat x\in(0,1)$ is the unique root in $(0,1)$ of the bad type's indifference condition \eqref{t3e9-xhat}, the equation in the statement: the cheapest quality at which mimicking no longer pays.

\step{7}{The separating equilibrium exists only if $r$ is not too large (fourth bullet)} Write $y=r\hat x$
in \eqref{t3e9-xhat}. Using \eqref{t3e9-c},
\begin{align}
  A &= \hat xA-\hat xy+(1-\hat x)(c_0-\kappa y) \nonumber\\
  A-\hat xA &= -\hat xy+(1-\hat x)c_0-(1-\hat x)\kappa y \nonumber\\
  (1-\hat x)A-(1-\hat x)c_0 &= -y\,[\hat x+(1-\hat x)\kappa] \nonumber\\
  (1-\hat x)(c_0-A) &= y\,[\hat x+(1-\hat x)\kappa] \nonumber\\
  y &= \frac{(1-\hat x)(c_0-A)}{\hat x+(1-\hat x)\kappa}, \label{t3e9-y}\\
  r = \frac{y}{\hat x} &= \frac{(1-\hat x)(c_0-A)}{\hat x\,[\hat x+(1-\hat x)\kappa]}. \label{t3e9-r}
\end{align}
The right side of \eqref{t3e9-r} falls strictly from $+\infty$ to $0$ as $\hat x$ goes from $0$ to $1$, so
$\hat x$ is a strictly decreasing function of $r$ ($\hat x\to1$ as $r\to0$). The right side of
\eqref{t3e9-y} is strictly decreasing in $\hat x$, so the signalling cost $y(r)$ is strictly increasing in
$r$, from $0$ (as $r\to0$) to $(c_0-A)/\kappa$ (as $r\to\infty$, if $p_L>0$).

The candidate is an equilibrium only if three inequalities hold, each an upper bound on $y$:
\begin{itemize}
  \item[(E1)] the good type is funded and works, \eqref{t3e9-icG} at $\hat x$: $y\le p_H(R-b/\Delta p)-(I-A)$, positive by (1);
  \item[(E2)] the good type does not deviate. A deviation to $x'\ne\hat x$ is treated as bad unless the signal
     reveals $b$, so it yields at most $A-rx'+x'(p_HR-I)$. Subtracting this from \eqref{t3e9-UG} at $\hat x$,
     \begin{align}
       D(x') = (1-x')(p_HR-I)-r(\hat x-x'), \nonumber
     \end{align}
     which is positive for $x'\ge\hat x$ and linear on $[0,\hat x]$ with $D(\hat x)=(1-\hat x)(p_HR-I)>0$; so
     the binding deviation is $x'=0$: $y\le p_HR-I$;
  \item[(E3)] the fee is paid from cash: $y\le A$.
\end{itemize}
The bad type's deviations give her at most $A-rx'\le A$, so \eqref{t3e9-UB} is optimal. Let
$\bar y\equiv\min\{p_H(R-b/\Delta p)-(I-A),\,p_HR-I,\,A\}>0$. Since $y(r)$ rises from $0$, the
separating equilibrium exists if and only if $y(r)\le\bar y$, i.e.\ $r\le\bar r$, where $\bar r$ solves
$y(\bar r)=\bar y$ ($\bar r=+\infty$ if $(c_0-A)/\kappa\le\bar y$). Inverting \eqref{t3e9-y} at $y=\bar y$ (multiply out and collect the terms in $\bar x$):
\begin{align}
  \bar y\,[\bar x+(1-\bar x)\kappa] &= (1-\bar x)(c_0-A) \nonumber\\
  \bar y\kappa+\bar x\,\bar y(1-\kappa) &= (c_0-A)-\bar x(c_0-A) \nonumber\\
  \bar x\,[(c_0-A)+(1-\kappa)\bar y] &= c_0-A-\kappa\bar y \nonumber\\
  \bar x = \frac{c_0-A-\kappa\bar y}{c_0-A+(1-\kappa)\bar y}, &\qquad \bar r = \frac{\bar y}{\bar x}. \label{t3e9-rbar}
\end{align}

\answer{(ii-4)} The signalling cost $y=r\hat x$ rises with $r$, and the equilibrium needs $y\le\bar y=\min\{p_H(R-b/\Delta p)-(I-A),\;p_HR-I,\;A\}$ (good type works, does not prefer to deviate to $x=0$, can pay the fee). Hence it exists iff $r\le\bar r=\bar y/\bar x$ in \eqref{t3e9-rbar}.

\step{8}{Numerical illustration} Take $R=10$, $I=6$, $A=1$, $p_H=0.8$, $p_L=0.4$ ($\Delta p=0.4$), $b=0.4$,
$B=1.8$, $\alpha=0.3$, so $b/\Delta p=1$, $B/\Delta p=4.5$ and $I-A=5$. Condition (1): $0.8(10-4.5)=4.4<5<0.8(10-1)=7.2$. Condition (2):
$[0.3(0.8)+0.7(0.4)](10-1)=(0.24+0.28)(9)=(0.52)(9)=4.68<5$; and $p_LR+B=4+1.8=5.8<6$. Then
\begin{align}
  c_0 &= (0.4)\Bigl(10-\frac{5}{0.8}\Bigr)+1.8 = (0.4)(3.75)+1.8 = 1.5+1.8 = 3.3, \qquad \kappa=\frac{0.4}{0.8}=0.5. \nonumber
\end{align}
With $r=1$, \eqref{t3e9-Phi} has coefficients $-r(1-\kappa)=-0.5$, $A-c_0-\kappa r=1-3.3-0.5=-2.8$ and $c_0-A=2.3$:
\begin{align}
  -0.5x^2-2.8x+2.3 &= 0 \iff x^2+5.6x-4.6=0, \nonumber\\
  \hat x &= \frac{-5.6+\sqrt{5.6^2+4(4.6)}}{2} = \frac{-5.6+\sqrt{49.76}}{2} = -2.8+\sqrt{12.44} = 0.727038. \nonumber
\end{align}
Then $y=r\hat x=0.727038$, $U_G=p_HR-I+A-y=8-6+1-0.727038=2.272962$
and $R_b^G=10-(5+0.727038)/0.8=10-7.158798=2.841202\ge b/\Delta p=1$. Here
\begin{align}
  \bar y=\min\{7.2-5,\;8-6,\;1\}=\min\{2.2,\,2,\,1\}=1, \nonumber
\end{align}
so \eqref{t3e9-rbar} gives $\bar x=(2.3-0.5)/(2.3+0.5)=1.8/2.8=0.642857$ and $\bar r=1/0.642857=1.555556$: the rating equilibrium exists for
$r\le1.555556$.

\intuition Lending collapses because whatever contract attracts the good borrower attracts the bad
one, and condition (2) says the pool is unprofitable. A rating of quality $x$ lets the good type escape the
pool, but only partly: the bad type can buy the same rating and hope that it reveals nothing. The good type
therefore buys just enough quality that the gamble ``with probability $1-x$ I pass as good'' is not worth
the bad type's sure $A$. The cost $rx$ is pure dissipation, and when it eats the good type's surplus (or
cash, or incentive margin) separation fails.

\pitfalls
\begin{itemize}
  \item Forgetting that the fee $rx$ is paid out of $A$, so the good type must borrow $I-A+rx$, not $I-A$; this enters both her terms and the bad type's mimicking payoff.
  \item Writing the mimicking payoff as if the bad type worked under $R_b^G$: since $R_b^G<B/\Delta p$ she shirks, and her payoff is $p_LR_b^G+B$.
  \item Choosing $x=1$ (``buy the best rating''): the good type wants the \emph{cheapest} separating quality, so the no-mimicking constraint binds at $\hat x<1$.
\end{itemize}

\bookans Matches the book's answer (pp.~598--599): the bad type chooses $x=0$ and is not funded, the
good type's terms solve $p_H(R-R_b^G)=I-A+rx$, and the no-mimicking condition is the displayed equation,
binding at the equilibrium. The book does not spell out the ``$r$ not too large'' conditions; (E1)--(E3)
and \eqref{t3e9-rbar} are derived here.

\refbib{Tirole Ch.~6, \S6.3 Application 4 ``Certification'', pp.~249--251; Exercise 6.6, p.~276;
answer pp.~598--599. Lecture Topic~3, signalling block.}

% =============================================================================
\exsec{10. Maturity structure --- Tirole 6.10, p.~279}
%% BEGIN-STATEMENT 10
\stmtbegin
\textbf{Source:} Tirole, Ch. 6, Exercise 6.10, p.~279 (PDF 292) · Answer in book: yes, p.~600

At date 0 the entrepreneur has cash on hand \emph{A} and needs to finance an investment of fixed
size \(I\). At date 1, a deterministic income \(r\) accrues; a liquidity shock must be met in order
for the firm to continue. Liquidation yields nothing. The probability of success in the case of
continuation depends on a date-1 effort: for a good borrower, this probability is \(p_H\) or \(p_L\)
depending on whether she behaves (no private benefit) or misbehaves (private benefit \(B\));
similarly, for a bad borrower, it is \(q_H\) or \(q_L\). We assume that

\[p_H - p_L = q_H - q_L = \Delta p,\]

and so the incentive compatibility constraint in the case of continuation is the same for both types
of borrower:

\[(p_H - p_L)R_b = (q_H - q_L)R_b = (\Delta p)R_b \geqslant B,\]

where \(R_b\) is the borrower's reward in the case of success.

The borrower knows at the date of contracting whether she is a ``\(p\)-type'' or a ``\(q\)-type.''
Let

\[\rho_0^G \equiv p_H\left(R - \frac{B}{\Delta p}\right) \quad \text{and} \quad \rho_0^B = q_H\left(R - \frac{B}{\Delta p}\right)\]

denote the date-1 pledgeable incomes for the good and bad types.

The liquidity shock is deterministic and equal to \(\rho\). Information is asymmetric at date 0, but
the capital market learns the borrower's type perfectly at date 1, before the liquidity shock has to
be met. Assume that

\[\rho_0^G > \rho > \rho_0^B.\]

Suppose further that under symmetric information only the good borrower is creditworthy (provided
that she is incentivized to behave).

Assume that \(r < I - A < r + [\rho_0^G - \rho]\). Show that the good borrower can costlessly signal
her type.

\paragraph{Context needed to solve it}

\emph{From Tirole §6.3, Application 6 ``Short-Term Maturities'', pp.~254--255 (PDF 267--268) --- the
model of which the exercise is a variant (there the liquidity shock is random; in the exercise it is
deterministic and the asymmetry is about the success probabilities). It defines \(R\) and the
preferences the statement leaves implicit:}

Let us consider the following variant of the model of debt maturities set up in Section 5.2. At date
0 the entrepreneur has a project of fixed size \(I\), has wealth \(A\), and must borrow \(I - A\).
At date 1, the investment yields a deterministic and verifiable short-term (date-1) profit
\(r > 0\). With probability \(\lambda\), continuation requires reinvesting \(\rho\) (the liquidation
value is equal to 0); with probability \(1 - \lambda\), continuation requires no
reinvestment.\textsuperscript{51}

In the case of continuation, the firm at date 2 yields \(R\) in the case of success and 0 in the
case of failure. The probability of success is \(p_H\) in the case of good behavior and
\(p_L = p_H - \Delta p\) in the case of misbehavior (which yields private benefit \(B\)). Moral
hazard is introduced in order to create an entrepreneurial rent from continuation, or, equivalently,
a cost for the entrepreneur associated with early termination. Put differently, moral hazard
introduces a friction in the date-1 refinancing market. Let

\[\rho_1 \equiv p_H R \quad \text{and} \quad \rho_0 \equiv p_H\left(R - \frac{B}{\Delta p}\right)\]

denote the continuation NPV and pledgeable income, respectively.

{[}\ldots{]} the borrower can use retentions \((r - d)\) together with what can be raised in the
capital market \((\rho_0)\) to meet shock \(\rho\).

\begin{quote}
\begin{enumerate}
\def\labelenumi{\arabic{enumi}.}
\setcounter{enumi}{50}
\tightlist
\item
  As usual, the entrepreneur is risk neutral and protected by limited liability; and the investors
  are risk neutral and demand a 0 rate of return.
\end{enumerate}
\end{quote}

\emph{Figure 6.1 (timeline, left to right):} Date 0: Entrepreneur has wealth \(A\) and fixed-size
investment project \(I - A\). → Date 1: Short-run income \(r\). Reinvestment need \(\rho\) (prob.
\(\lambda\)) or 0 (prob. \(1 - \lambda\)). → (if reinvestment) Moral hazard (\(p = p_H\) or
\(p_L\)). → Outcome: success (\(R\)) or failure (0).
\stmtend
%% END-STATEMENT

\asked (a) Two types of entrepreneur differ only in their date-2 success probabilities; the market learns the type at date 1, before a deterministic liquidity shock $\rho$ must be paid. Only the good type can refinance the shock in the market ($\rho_0^G>\rho>\rho_0^B$) and only she is creditworthy. Show that she can separate herself from the bad type at no cost, i.e.\ while obtaining her symmetric-information payoff.

\approach This is signalling through the maturity of debt (Tirole \S6.3, Application 6, built on the liquidity model of Ch.~5; lecture block 3.7). A signal separates if it is cheap for the good type and expensive for the bad one. Short-term debt that takes all the date-1 income ($d=r$) leaves no hoarded liquidity, so the firm must go back to the market at date 1, after its type is known. The good type can refinance there; the bad type cannot and is liquidated. So the good type offers that contract, checks that it gives her the symmetric-information payoff (incentive and break-even constraints), and checks that the bad type does not want to copy it (no-mimicking, including renegotiation), with pessimistic off-path beliefs.

\notation
\begin{itemize}
  \item $A$: cash at date 0; $I$: fixed investment, $I-A>0$ borrowed; $r>0$: deterministic date-1 income.
  \item $\rho$: deterministic date-1 liquidity shock; liquidation value $0$.
  \item $p_H,p_L$ (good) and $q_H,q_L$ (bad): date-2 success probabilities when behaving / misbehaving; $p_H-p_L=q_H-q_L=\Delta p$.
  \item $B$: private benefit from misbehaving; $R$: date-2 income in success, 0 in failure; $R_b$: borrower's reward in success.
  \item $\rho_0^G=p_H(R-B/\Delta p)$, $\rho_0^B=q_H(R-B/\Delta p)$: date-1 pledgeable incomes, $\rho_0^G>\rho>\rho_0^B$.
  \item $d\in[0,r]$: short-term debt repaid at date 1 (retention $r-d$); $L$: date-0 investors' long-term claim on date-2 income (expected value).
  \item $R_n$, $R_l$: shares of $R$ promised to date-1 (new) and date-0 (old) investors.
\end{itemize}

\assume (a) As in Tirole's Application 6, a contract specifies $d$, a long-term claim and $R_b$; new claims
issued at date 1 may dilute the old long-term claim. (b) A continuation in which the borrower misbehaves
cannot be financed ($p_LR<\rho$ and $q_LR<\rho$), so date-1 investors fund only continuations with
$(\Delta p)R_b\ge B$. (c) Off-path beliefs: any date-0 offer other than the good type's equilibrium
contract is attributed to the bad type.

\step{1}{Symmetric-information benchmark} With the type known, the good borrower can always be continued
(the shock is deterministic and $\rho_0^G>\rho$, so the date-1 market covers it). Investors break even, so
her gross payoff is the project's value net of all outlays:
\begin{align}
  U_G^{SI} = A + r + p_HR - \rho - I. \label{t3e10-USI}
\end{align}
Funding is feasible because total pledgeable income covers the loan: investors can receive the date-1 income $r$ and, at date 2, the pledgeable income $\rho_0^G$ minus the $\rho$ that must be paid to whoever finances the shock, so they break even iff
\begin{align}
  r+(\rho_0^G-\rho)\ge I-A, \nonumber
\end{align}
which holds by the right inequality of the statement. The bad type is not creditworthy (statement), so her
symmetric-information payoff is to keep her cash:
\begin{align}
  U_B^{SI} = A. \label{t3e10-UBSI}
\end{align}

\step{2}{The signalling contract: short-term debt $d=r$} The good type offers: short-term debt $d=r$ (no
retention, no hoarded liquidity), a dilutable long-term claim worth
\begin{align}
  L = I-A-r > 0 \qquad (\text{by } r<I-A), \label{t3e10-L}
\end{align}
and a reward $R_b$ in success. Date-0 investors receive $d+L=r+(I-A-r)=I-A$: they break even. At date 1
the market learns she is good and she must raise $\rho$ by issuing a new claim $R_n$ with
$p_HR_n=\rho$. The old claim is $p_HR_l=L$. The borrower keeps the rest, so
\begin{align}
  p_HR_b &= p_HR-p_HR_n-p_HR_l \nonumber\\
         &= p_HR-\rho-L \nonumber\\
         &= p_HR-\rho-(I-A-r) && \text{(substituting \eqref{t3e10-L})} \nonumber\\
         &= A+r+p_HR-\rho-I = U_G^{SI} && \text{(by \eqref{t3e10-USI})}. \label{t3e10-UG}
\end{align}
She behaves if $(\Delta p)R_b\ge B$, i.e.\ (multiplying by $p_H/\Delta p$) $p_HR_b\ge p_HB/\Delta p$; by the second line of
\eqref{t3e10-UG} this is
\begin{align}
  p_HR-\rho-L &\ge p_H\frac{B}{\Delta p} \nonumber\\
  p_H\Bigl(R-\frac{B}{\Delta p}\Bigr) &\ge \rho+L \iff \rho_0^G\ge\rho+(I-A-r) \nonumber\\
  I-A &\le r+(\rho_0^G-\rho), \label{t3e10-ic}
\end{align}
which is the right inequality assumed in the statement. So the contract is feasible and gives the good
type exactly her symmetric-information payoff: signalling is costless.

\step{3}{The bad type does not mimic} Suppose the bad type offers the same contract. At date 1 all her
income $r$ goes to repay $d=r$. The market learns she is bad; even with the old claim fully diluted, the
most new investors can be promised while she behaves is $\rho_0^B<\rho$, and by assumption (b) a misbehaving
continuation cannot be funded either. The project is liquidated and yields $0$, so
\begin{align}
  U_B^{\text{mimic}} = 0 \le A = U_B^{SI}. \label{t3e10-mimic}
\end{align}
Could the old investors waive $d$ to let her continue? Then the firm would keep $r$ and would need
$\rho-r$ from the market; out of the bad type's date-2 pledgeable income $\rho_0^B$, the new money takes $\rho-r$, so the old investors could get at most
\begin{align}
  \rho_0^B-(\rho-r)=r-(\rho-\rho_0^B)<r, \nonumber
\end{align}
less than collecting $d=r$ (because $\rho>\rho_0^B$). So they do not renegotiate, and \eqref{t3e10-mimic} stands.

\step{4}{The separating equilibrium} The good type offers the Step~2 contract and the bad type stays out.
Investors believe the Step~2 contract comes from the good type (correct, by Step~3) and fund it; any other
offer is attributed to the bad type and, since she is not creditworthy, is not funded. The good type
cannot gain by deviating: she already gets $U_G^{SI}$ in \eqref{t3e10-UG}, and other offers are not funded.
The bad type gets $A$, by \eqref{t3e10-mimic} at least as much as by mimicking. Hence the good borrower
signals her type at zero cost. $\blacksquare$

Signalling works because liquidity and refinancing are worth different amounts to the two types. A
contract that hoarded enough liquidity, $r-d\ge\rho-\rho_0^B$, would let a mimicking bad type meet the
shock with retentions plus new money, so she would get a rent instead of 0. The good type gives up
nothing by not hoarding, because the date-1 market will see that she is good and lend to her.

\answer{(a)} The good type offers short-term debt $d=r$ (no hoarded liquidity) plus a dilutable long-term claim worth $L=I-A-r$. Date-0 investors break even, she refinances $\rho$ at date 1 once revealed good and obtains $p_HR_b=A+r+p_HR-\rho-I=U_G^{SI}$, so the signal costs her nothing. A mimicking bad type would be liquidated at date 1 (payoff $0<A$), so she stays out: a costless separating equilibrium.

\step{5}{Numerical illustration} Take $p_H=0.8$, $q_H=0.5$, $\Delta p=0.4$, $B=1.2$ ($B/\Delta p=1.2/0.4=3$),
$R=10$, $\rho=4.5$, $r=1$, $I-A=1.6$. Then $\rho_0^G=0.8(10-3)=0.8(7)=5.6$ and $\rho_0^B=0.5(7)=3.5$, so
$\rho_0^G>\rho>\rho_0^B$; $r=1<I-A=1.6<r+\rho_0^G-\rho=1+5.6-4.5=2.1$; and the bad type fails even with full hoarding,
$r+\rho_0^B-\rho=1+3.5-4.5=0<1.6$. By \eqref{t3e10-L}, $L=1.6-1=0.6$. Then
\begin{align}
  R_n &= \frac{\rho}{p_H} = \frac{(4.5)}{(0.8)} = 5.625, \qquad R_l = \frac{L}{p_H} = \frac{(0.6)}{(0.8)} = 0.75, \nonumber\\
  p_HR_b &= r+p_HR-\rho-(I-A) = 1+8-4.5-1.6 = 2.9, \qquad R_b = \frac{2.9}{0.8} = 3.625\ge\frac{B}{\Delta p}=3, \nonumber
\end{align}
and $R_b+R_n+R_l=3.625+5.625+0.75=10=R$.

\intuition Short-term debt drains the firm's cash at date 1 and forces it back to the capital market just
after the market has learned the firm's type. The good type is not afraid of that meeting; the bad type is,
because the market will refuse to fund her shock. Short maturity is an action whose cost falls only on
the type that has something to hide, which is why it separates. Here it costs the good type nothing only
because the shock is certain and small enough to refinance. With a random shock that may exceed
$\rho_0^G$, she would want to hoard liquidity even under symmetric information, and the signal would cost
her some continuation (Application 6).

\pitfalls
\begin{itemize}
  \item Thinking the good type should hoard liquidity, as in Ch.~5: with a certain shock and a market that will know she is good, hoarding is unnecessary, and it would let the bad type survive by mimicking.
  \item Forgetting the date-1 claim structure: the old long-term claim is diluted by the new money raised for $\rho$, and the incentive constraint must hold \emph{after} both claims are paid ($\rho+L\le\rho_0^G$).
  \item Ignoring renegotiation: one must check that old investors would not waive $d$ to keep a revealed bad type alive.
\end{itemize}

\bookans Matches the book's answer (p.~600): separation by not hoarding liquidity, $d=r$, the good type
refinancing at date 1 and diluting existing claims, the bad type liquidated. The long-term claim
$L=I-A-r$, the renegotiation check and assumption (b) are added here; the book leaves them implicit.

\refbib{Tirole Ch.~6, \S6.3 Application 6 ``Short-Term Maturities'', pp.~254--257; Ch.~5, \S5.2
(liquidity hoarding and short-term debt); Exercise 6.10, p.~279; answer p.~600.}

% =============================================================================
\exsec{11. Myers--Majluf: whose interest does management serve? --- Amaro de Matos 3.2.6 \#4, p.~95}
%% BEGIN-STATEMENT 11
\stmtbegin
\textbf{Source:} Amaro de Matos, Ch. 3, Exercise 3.2.6 \#4, p.~95 (PDF 113) · Answer in book: yes,
p.~260

Instead of Assumption 3.2.6, assume that the management acts in the interest of all shareholders.
How would that affect the investment decision of the management? If management acts in the interest
of the new shareholders, what is the optimal equity issuance rule?

\paragraph{Context needed to solve it}

\emph{From Amaro de Matos §3.2, p.~78 (PDF 96):}

\textbf{Assumption 3.2.1.} \emph{Managers have better information about the value of the firm than
other agents in the economy.}

\emph{From Amaro de Matos §3.2.3 (Signals Conditioned by Investment Opportunities), pp.~86--88 (PDF
104--106):}

To see how all these conclusions are derived, the model must be formalized. Let \(I\) be the
required investment and \(S\) denote the \emph{slack} of the firm (i.e., large holdings of cash or
marketable securities) or the ability to issue default-risk-free debt.

If \(S \geq I\), it will not be necessary to issue any equity at all in order to finance the
investment. The only interesting case happens when \(S < I\), where the required amount of equity to
be issued is \(E = I - S\).

As opposed to the other models previously discussed, this is a three-date model with three distinct
times---say, time 0, time 1, and time 2. Let us introduce the following notation. At time 0, the
possible value of the assets of the firm at time 1 is described by a random variable
\(\widetilde{A}\). Similarly, let \(\widetilde{B}\) denote the possible net present value of the
investment opportunity at time 1.

At time 0, the information is symmetric. All the agents, managers and market, have the same
information---namely, the distributions of \(\widetilde{A}\) and \(\widetilde{B}\) and the value of
\(S\).

At time 1, the asymmetry of information between managers and outsider investors is specified.
Assumption 3.2.1 takes the form

\textbf{Assumption 3.2.5.} \emph{At time 1, the management knows the realizations of
\(\widetilde{A}\) and \(\widetilde{B}\), denoted respectively by \(a\) and \(b\), which are unknown
to the outsider investors. Negative values of \(a\) and \(b\) are ruled out.}

At this stage, the only thing that the outsiders know is that the amount of equity to be issued is
\(E = \max[0, I - S]\). Also, at time 1, let \(P\) be the market value of the shares of old
shareholders if the stock is not issued and \(P'\) be the price of the same shares if it is issued.
At time 2, finally, all the information is disclosed.

Let \(V^\circ\) be the value of the firm for the old shareholders. If the management, knowing the
values of \(a\) and \(b\), decides not to issue equity, then the company is still owned by the same
old shareholders with assets of value \(a\) and slack \(S\). Thus

\[V^\circ = a + S.\]

If, on the other hand, knowing the value of \(a\) and \(b\), the management decides to issue equity
\(E\), then the ownership of the company is shared between the old and new shareholders. The total
value of the company will include the issued equity \(E\) and the net present value of the
investment opportunity apart from the assets of value \(a\) and slack \(S\). But the old
shareholders will have only a proportion \(P'/(P' + E)\) of all of this. Thus, in this case

\[V^\circ = \frac{P'}{P' + E}(E + S + a + b).\]

To explain the decision of the managers, the following assumption is required:

\textbf{Assumption 3.2.6.} \emph{The management acts in the best interest of the old shareholders.
Moreover, old shareholders are assumed to be passive---that is, they do not adjust their portfolios
in response to the issue--invest decision of the firm.}

If that is so, the decision to issue new equity is taken only if the old shareholders are better off
in such case---that is, only if

\[\frac{P'}{P' + E}(E + S + a + b) \geq a + S.\]

This inequality may be rewritten as

\[b \geq -E + \frac{E}{P'}(S + a) \qquad (3.20)\]

and should naturally be restricted by Assumption 3.2.5 to the region \(a \geq 0\), \(b \geq 0\).

Hence, inequality (3.20) divides the domain of the joint probability distribution of
\(\widetilde{A}\) and \(\widetilde{B}\) into two disjoint regions identified in Figure 3.1: a region
\(M'\) in which all pairs \((a, b)\) satisfy the inequality and the firm issues new equity and
invests in the new project, and a region \(M\) in which no pair \((a, b)\) satisfies the inequality
and there is no investment. In the first region, the amount of equity issued is \(E = I - S\),
whereas in the second, it is zero.

\emph{Figure 3.1 (The Regions of Investment and No-Investment Decisions According to Myers and
Majluf's Model):} horizontal axis \(a\), vertical axis \(b\), both starting at 0. An upward-sloping
straight line labelled \(b = -E + \frac{E}{P'}(S + a)\) starts on the horizontal axis at a positive
value of \(a\) (where \(b = 0\)). The area above/left of the line is labelled ``region \(M'\)''; the
area below/right of the line is labelled ``region \(M\)''.

It is simple to see that the price \(P\) for the shares when \((a, b) \in M\) is greater than the
price \(P'\) when \((a, b) \in M'\). The equilibrium prices should be given by

\[P = S + E(\widetilde{A} \mid M) \qquad (3.21)\]

and

\[P' = S + E(\widetilde{A} + \widetilde{B} \mid M').\]

For any \((a, b) \in M\), the inequality (3.20) is not satisfied and hence

\[b < -E + \frac{E}{P'}(S + a) \quad \text{or} \quad a > P'(1 + b/E) - S.\]

Since \(b/E \geq 0\) from Assumption 3.2.5, it follows that for any \((a, b) \in M\),
\(a > P' - S\). Thus \(E(\widetilde{A} \mid M) > P' - S\) and (3.21) gives directly \(P > P'\).
\stmtend
%% END-STATEMENT

\asked (a) In the Myers--Majluf model, replace ``management serves old shareholders'' (Assumption 3.2.6) with ``management serves all shareholders'' (old and new together): how does the issue-and-invest decision change? (b) If management serves the new shareholders, what is the optimal issue rule?

\approach This is the Myers--Majluf model of Amaro de Matos \S3.2.3 (lecture block 3.5). The issue price $P'$ moves wealth between old and new shareholders, while the project's NPV $b$ adds to the total. Which side of that transfer management cares about decides the rule: write each group's wealth with and without the issue, and compare. The book's rule (3.20) is the old-shareholder comparison; the all-shareholder and new-shareholder rules come from the same algebra. A last step closes the model with the book's pricing equation to see what a rational market does with the new-shareholder rule.

\notation
\begin{itemize}
  \item $I$: required investment; $S<I$: slack (cash, riskless debt capacity); $E=I-S>0$: equity to issue.
  \item $a\ge0$, $b\ge0$: realisations (known to managers at time 1) of the value of assets in place $\widetilde A$ and of the project's NPV $\widetilde B$.
  \item $P'$: market value of the old shares if equity is issued; $P$: their value if not; the new shareholders hold the fraction $E/(P'+E)$.
  \item $M'$ / $M$: regions of $(a,b)$ with / without issue; $Z\equiv\widetilde A+\widetilde B$.
\end{itemize}

\step{1}{What each group gets} If the firm does not issue, the old shareholders own $a+S$ and the potential
new shareholders keep their cash $E$. If it issues $E$ and invests, the firm is worth $E+S+a+b$, split
in proportions $P'/(P'+E)$ and $E/(P'+E)$. So
\begin{align}
  \text{old:}&\quad a+S \;\;\text{vs.}\;\; \frac{P'}{P'+E}(E+S+a+b), \label{t3e11-old}\\
  \text{new:}&\quad E \;\;\text{vs.}\;\; \frac{E}{P'+E}(E+S+a+b), \label{t3e11-new}\\
  \text{all:}&\quad a+S+E \;\;\text{vs.}\;\; E+S+a+b. \label{t3e11-all}
\end{align}
Assumption 3.2.6 uses \eqref{t3e11-old}; multiplying it by $(P'+E)/P'$ and rearranging gives rule (3.20):
\begin{align}
  \frac{P'}{P'+E}(E+S+a+b) &\ge a+S \nonumber\\
  E+S+a+b &\ge \frac{P'+E}{P'}(a+S) = (a+S)+\frac{E}{P'}(S+a) \nonumber\\
  b &\ge -E+\frac{E}{P'}(S+a). \tag{3.20}
\end{align}

\step{2}{Management serves all shareholders} By \eqref{t3e11-all}, issue and invest iff
\begin{align}
  E+S+a+b &\ge E+S+a \nonumber\\
  b &\ge 0. \label{t3e11-allrule}
\end{align}
The issue price only moves wealth between old and new shareholders, and that transfer nets out in their
joint wealth. So management takes every positive-NPV project, the first best. By Assumption 3.2.5
($b\ge0$) every firm with $S<I$ issues: $M'$ is the whole domain, $M$ is empty, the issue carries no
information, and the book's pricing equation gives the pooled price
\begin{align}
  P' = S+E(\widetilde A+\widetilde B). \label{t3e11-Pall}
\end{align}
There is no underinvestment and no price reaction to the issue.

\answer{(a)} Management then issues and invests iff $b\ge0$: every positive-NPV project is taken (first best), every firm with $S<I$ issues, and the issue carries no information.

\step{3}{Management serves the new shareholders: the issue rule} By \eqref{t3e11-new}, issue iff the new shareholders' stake is worth at least the $E$ they pay; dividing by $E/(P'+E)>0$ and subtracting $E$ from both sides,
\begin{align}
  \frac{E}{P'+E}(E+S+a+b) &\ge E \nonumber\\
  E+S+a+b &\ge P'+E \nonumber\\
  b &\ge P'-(S+a). \label{t3e11-newrule}
\end{align}
The firm issues when its post-issue value, net of the new money, is at least the price $P'$, i.e.\ when
new shares are \emph{not overpriced}.

\answer{(b)} Issue (and invest) iff $b\ge P'-(S+a)$, equivalently $a+b\ge P'-S$: only when the new shares are not overpriced.

\step{4}{Comparison with the Myers--Majluf rule} Write $a_0\equiv P'-S$. At $a=a_0$ both thresholds are
zero: \eqref{t3e11-newrule} gives $b\ge0$, and (3.20) gives $b\ge-E+\frac{E}{P'}P'=0$. The new-shareholder
threshold $P'-S-a$ falls with $a$ (slope $-1$); the old-shareholder threshold
$-E+\frac{E}{P'}(S+a)$ rises with $a$ (slope $E/P'$). Hence, for a given $P'$:
\begin{itemize}
  \item $a>a_0$ (assets in place undervalued): the new-shareholder threshold is negative, so every such firm
        issues; the Myers--Majluf rule requires $b>0$ and passes up projects. Here the new rule is \emph{less} restrictive.
  \item $a<a_0$ (assets in place overvalued): the Myers--Majluf threshold is negative, so every such firm
        issues; the new-shareholder rule requires $b\ge a_0-a>0$. Here it is \emph{more} restrictive.
\end{itemize}
For example, with $P'=10$, $S=2$, $E=5$ ($a_0=10-2=8$): at $a=4$ the thresholds are $10-2-4=4$ (new) and $-5+\frac{5}{10}(2+4)=-2$ (old); at
$a=12$ they are $10-2-12=-4$ (new) and $-5+\frac{5}{10}(2+12)=2$ (old).

\step{5}{Equilibrium price under the new-shareholder rule} Close the model with the book's pricing
equation for the issue state, $P'=S+E(\widetilde A+\widetilde B\mid M')$, and
$M'=\{a+b\ge P'-S\}$ from \eqref{t3e11-newrule}. With $c\equiv P'-S$,
\begin{align}
  P'-S &= E(\widetilde A+\widetilde B\mid \widetilde A+\widetilde B\ge P'-S) \nonumber\\
  c &= E(Z\mid Z\ge c) \nonumber\\
  0 &= E(Z-c\mid Z\ge c). \label{t3e11-fix}
\end{align}
On the event $Z\ge c$ the integrand $Z-c$ is nonnegative, so \eqref{t3e11-fix} forces $Z=c$ almost surely
on that event: only firms at the top of the support of $Z$ can issue. With a continuous distribution the
issue region has probability zero: new investors would only be offered shares that are cheap for them,
so the price rises until nothing is cheap, and the market for new equity breaks down. Every firm with
$S<I$ then passes up its project. This is worse than under Myers--Majluf, where only the projects in $M$
are lost.

\intuition Who management works for decides which way the information problem cuts. Working for old
shareholders, it sells overpriced equity and refuses to sell underpriced equity (Myers--Majluf
underinvestment). Working for all shareholders, the issue price is a transfer inside the group and only
the NPV counts. Working for the new shareholders, it sells only underpriced equity, so the rule
favours undervalued firms. But a rational market prices that in and leaves no underpriced issues.

\pitfalls
\begin{itemize}
  \item Keeping the old-shareholder rule (3.20) in part (a): for all shareholders together the issue price is an internal transfer and drops out, leaving $b\ge0$.
  \item Sign errors when isolating $b$: the old shareholders' threshold rises with $a$, the new shareholders' falls with $a$.
  \item Calling the new-shareholder rule simply ``less restrictive'': it is less restrictive only for $a>P'-S$ and more restrictive for $a<P'-S$ (Step~4).
\end{itemize}

\bookans Both rules match the book's answer (p.~260): $b\ge0$ for all shareholders and
$b\ge P'-(S+a)$ for new shareholders. The book calls the second rule ``less restrictive'' than Myers and
Majluf's. Step~4 shows this holds only for $a>P'-S$; for $a<P'-S$ it is more restrictive. The
equilibrium in Step~5 goes beyond the book.

\refbib{Amaro de Matos Ch.~3, \S3.2.3 ``Signals Conditioned by Investment Opportunities'', pp.~86--88
(rule (3.20), prices (3.21)); Appendix C.3.2, p.~260.}

% =============================================================================
\exsec{12. Pecking order with variable investment --- Tirole 6.8, p.~278}
%% BEGIN-STATEMENT 12
\stmtbegin
\textbf{Source:} Tirole, Ch. 6, Exercise 6.8, p.~278 (PDF 291) · Answer in book: yes, p.~599

Consider the privately-known-prospects model with risk neutrality and variable investment. For
investment \(I\), the realized income is either \(R^S I\) (in the case of success) or \(R^F I\) (in
the case of failure), where \(R^S > R^F \geqslant 0\). A good borrower has probability \(p_H\) of
success when working and \(p_L\) when shirking; similarly, a bad borrower has probability \(q_H\) of
success when working and \(q_L\) when shirking, where \(p_H - p_L = \Delta p = q_H - q_L\), for
simplicity. The entrepreneur's private benefit is 0 when working and \(BI\) when shirking. The
entrepreneur is risk neutral and protected by limited liability; the investors are risk neutral and
demand a rate of return equal to 0.

\begin{enumerate}
\def\labelenumi{(\roman{enumi})}
\item
  Let \(\tilde{U}_b^{SI}\) denote the bad borrower's gross utility under symmetric
  information.\textsuperscript{80} Consider the problem of maximizing the good borrower's utility
  subject to the investors' breaking even on that borrower, to the mimicking constraint that the
  good borrower's terms not be preferred by the bad borrower to her symmetric-information terms, and
  to the no-shirking constraint. Let \(\{R_b^S, R_b^F\}\) denote the (nonnegative) rewards of the
  good borrower in the cases of success and failure. Write the separating program.
\item
  Show that \(R_b^F = 0\).
\item
  (Only if you have read the supplementary section.) Show that the separating outcome is the only
  perfect Bayesian equilibrium of the issuance game if and only if \(\alpha \leqslant \alpha^*\) for
  some threshold \(\alpha^*\).
\end{enumerate}

\begin{quote}
\begin{enumerate}
\def\labelenumi{\arabic{enumi}.}
\setcounter{enumi}{79}
\tightlist
\item
  This utility was derived in Section 3.4.2. It is equal to
  \[\left[1 + \frac{q_H R - 1}{1 - q_H(R - B/\Delta p)}\right]A\] if \(q_H R \geqslant 1\), and to
  \(A\) otherwise.
\end{enumerate}
\end{quote}

\paragraph{Context needed to solve it}

\emph{From Tirole §6.2 (privately-known-prospects model), p.~241 (PDF 254):}

\emph{Privately-known-prospects model.} A borrower/entrepreneur has no funds (\(A = 0\)) to finance
a project costing \(I\). The project yields \(R\) in the case of success and 0 in the case of
failure. The borrower and the lenders are risk neutral, and the borrower is protected by limited
liability. The interest rate in the economy is normalized at 0.

The borrower can be one of two types. A good borrower has a probability of success equal to \(p\). A
bad borrower has a probability of success \(q\). {[}\ldots{]}

The borrower has private information about her type. The capital market, which is competitive and
demands an expected rate of return equal to 0, puts probabilities \(\alpha\) and \(1 - \alpha\) on
the borrower being a good or a bad type, respectively. Under asymmetric information, the capital
market does not know whether it faces a ``\(p\)-borrower'' (a good borrower) or a ``\(q\)-borrower''
(a bad borrower). {[}\ldots{]}

\emph{From Tirole §3.4.1--3.4.2 (the continuous-investment model, where footnote 80's formula is
derived), pp.~127--128 (PDF 140--141):}

\textbf{3.4.1 The Continuous-Investment Model.} The continuous-investment model of this section is
the polar opposite of the fixed-investment one. The fixed-investment model depicts a situation in
which returns are sharply decreasing beyond some investment level. In contrast, we now assume that
there are constant returns to scale in the investment technology. An investment
\(I \in [0, \infty)\) yields income \(RI\), proportional to \(I\), in the case of success, and 0 in
the case of failure. The borrower's private benefit from misbehaving is also taken to be
proportional to investment. As before, the borrower has a choice between behaving, in which case she
derives no private benefit and the probability of success is \(p_H\), and misbehaving, that is,
enjoying private benefit \(BI\), and reducing the probability of success to
\(p_L = p_H - \Delta p < p_H\). {[}\ldots{]}

The borrower initially has cash \(A\), and must therefore borrow \(I - A\) to finance a project of
size \(I\). A loan agreement specifies that the lenders (who as before are assumed to make no
profit) and the borrower receive 0 each in the case of failure, and \(R_l\) and \(R_b\),
respectively, in the case of success, where \(R_l + R_b = RI\).

As in Section 3.2, we assume that investment has positive NPV (net present value), here per unit of
investment, if the borrower behaves,

\[p_H R > 1, \qquad (3.7)\]

but negative NPV otherwise,

\[1 > p_L R + B, \qquad (3.8)\]

so that unless one can control the agency problem the investment cannot be funded. We also make an
assumption that guarantees that the equilibrium investment is finite:

\[p_H R < 1 + \frac{p_H B}{\Delta p}. \qquad (3.9)\]

{[}\ldots{]} Finally, we keep assuming that the capital market is competitive. {[}\ldots{]}

\textbf{3.4.2 The Lenders' Credit Analysis.} Following the steps of Section 3.2, the incentive
compatibility and the breakeven conditions are

\[(\Delta p) R_b \geqslant BI \qquad (\text{IC}_b)\]

and

\[p_H(RI - R_b) \geqslant I - A. \qquad (\text{IR}_l)\]

In equilibrium, competitive lenders make no profit on the contract that is most advantageous for the
borrower; the borrower's net utility is therefore equal to the social surplus brought about by the
investment:

\[U_b = (p_H R - 1)I. \qquad (3.10)\]

From (3.10) it is optimal for the borrower to invest as much as possible. The upper bound on
investment and in turn her borrowing capacity (``outside financing capacity'' or ``debt capacity'')
are determined by constraints (IC\(_b\)) and (IR\(_l\)). Substituting (IC\(_b\)) into (IR\(_l\)), we
obtain

\[I \leqslant kA, \qquad (3.11)\]

where

\[k = \frac{1}{1 - p_H(R - B/\Delta p)} > 1. \qquad (3.12)\]

{[}\ldots{]} Another important concept {[}\ldots{]} is the shadow value \(v\) of equity (here cash).
The entrepreneur derives gross utility \(v > 1\) from one more unit of equity. Letting
\(U_b^g \equiv A + U_b\) denote the borrower's gross utility, and using (3.10) and (3.11), we have

\[U_b^g \equiv vA, \qquad (3.14)\]

where the shadow value of equity is

\[v = \frac{p_H B/\Delta p}{1 - p_H(R - B/\Delta p)} > 1. \qquad (3.15)\]

\emph{From Tirole, Supplementary Section §6.4 (Contract Design by an Informed Party: An
Introduction), pp.~264--268 (PDF 277--281) --- the ``issuance game,'' the separating program and the
uniqueness threshold \(\alpha^*\) used in part (iii):}

A borrower who attempts to raise funds from lenders has private information about some
characteristic (private benefit, value of assets in place, prospects of the firm, value of
collateral) that affects the lenders' payoff. The borrower may have type \(b\) or \(\tilde{b}\).
While the borrower knows her type, the lenders only know that this type is \(b\) with probability
\(\alpha\) and \(\tilde{b}\) with probability \(1 - \alpha\).

\emph{Figure 6.4 (timeline, left to right):} Borrower offers a contract → Investors accept/refuse →
Borrower picks action(s).

Let us, abstractly, denote the contractual terms faced by the borrower by \(c\). Let \(U_b(c)\) and
\(\tilde{U}_b(c)\) denote the two types' net utilities for arbitrary contractual terms
\(c\).\textsuperscript{68} Let \(U_l(c)\) and \(\tilde{U}_l(c)\) denote the investors' expected
profit when contractual terms are \(c\) and the borrower turns out to have type \(b\) and
\(\tilde{b}\), respectively.

\emph{Example (privately-known-prospects).} In Section 6.2, the borrower had possible types
\(b = p\) and \(\tilde{b} = q\). The contractual terms \(c\) were just the borrower's reward
\(R_b^S\) in the case of success. More generally, they also contain the probability of investment,
\(x\), her reward in the case of failure, \(R_b^F\), and in the absence of investment, \(R_b^0\),
even though the latter in equilibrium can be taken to be equal to 1, 0, and 0, respectively (see
Section 6.5). We have

\[U_b(c) = x[pR_b^S + (1 - p)R_b^F] + (1 - x)R_b^0,\]
\[\tilde{U}_b(c) = x[qR_b^S + (1 - q)R_b^F] + (1 - x)R_b^0,\]
\[U_l(c) = x[p(R - R_b^S) - (1 - p)R_b^F] - (1 - x)R_b^0,\]
\[\tilde{U}_l(c) = x[q(R - R_b^S) - (1 - q)R_b^F] - (1 - x)R_b^0.\]

{[}\ldots{]} Figure 6.4 describes the timing of the \emph{issuance game}. As earlier, we assume that
the borrower designs the issue and offers the associated claims to a competitive capital market.
Investors purchase the claims if and only if they expect a nonnegative profit. Lastly, the borrower
chooses some action(s).

A few clarifications are in order. First, we allow for post-contracting actions by the borrower in
order to accommodate situations in which the borrower can waste resources (as in Chapters 3--5).
Second, we said that investors subscribe ``if and only if they expect a nonnegative profit.'' The
expectation should be taken relative to the updated beliefs, that is, the investors' beliefs after
they observe the contract offer and thus possibly learn something about the borrower. Third, we will
analyze the perfect Bayesian equilibrium (or equilibria) of the issuance game.\textsuperscript{70}
Fourth, a ``contract'' can in principle be anything that the borrower sees fit to design. However,
for the purpose of the analysis, it can be shown that there is no loss of generality in assuming
that the borrower offers an ``option contract'' \((c, \tilde{c})\), that is, as many contractual
terms as there are possible types. {[}\ldots{]} It can further be shown that there is no loss of
generality in assuming that the option contract is ``incentive compatible,'' that is, that type
\(b\) prefers contractual terms \(c\) to \(\tilde{c}\) and type \(\tilde{b}\) prefers \(\tilde{c}\)
to \(c\) {[}\ldots{]}.

{[}\ldots{]} Let an \emph{allocation} be a pair of (possibly identical) type-contingent contractual
terms \((c, \tilde{c})\). That is, it defines contractual terms \(c\) for type \(b\) and
\(\tilde{c}\) for type \(\tilde{b}\). (Note that an option contract defines an allocation.)

\emph{Definition.} An allocation \((c, \tilde{c})\) is \emph{incentive compatible} if type \(b\)
prefers \(c\) to \(\tilde{c}\) and type \(\tilde{b}\) prefers \(\tilde{c}\) to \(c\):

\[U_b(c) \geqslant U_b(\tilde{c}) \quad \text{and} \quad \tilde{U}_b(\tilde{c}) \geqslant \tilde{U}_b(c).\]

{[}\ldots{]}

\emph{Definition.} An incentive-compatible allocation \((c, \tilde{c})\) is \emph{profitable
type-by-type} if

\[U_l(c) \geqslant 0 \quad \text{and} \quad \tilde{U}_l(\tilde{c}) \geqslant 0.\]

\emph{Definition.} An incentive-compatible allocation \((c, \tilde{c})\) is \emph{profitable in
expectation} (relative to the prior beliefs) if

\[\alpha U_l(c) + (1 - \alpha)\tilde{U}_l(\tilde{c}) \geqslant 0.\]

An incentive-compatible allocation \((c, \tilde{c})\) that is profitable in expectation is
\emph{interim efficient} if it is Pareto-optimal for the two types of borrower in the set of
incentive-compatible, profitable-in-expectation allocations.

{[}\ldots{]} \emph{Definition.} Utility \(U_b(c_0)\) for borrower type \(b\) is the
\emph{low-information-intensity optimum} for that type if \(c_0\) maximizes type \(b\)'s utility in
the set of incentive-compatible, profitable-type-by-type allocations. That is, it is (part of) the
solution to the following program.

\emph{Program I (type \(b\)):}

\[\max_{\{c, \tilde{c}\}} U_b(c) \quad \text{s.t.} \quad U_b(c) \geqslant U_b(\tilde{c}), \quad \tilde{U}_b(\tilde{c}) \geqslant \tilde{U}_b(c), \quad U_l(c) \geqslant 0, \quad \tilde{U}_l(\tilde{c}) \geqslant 0.\]

The low-information-intensity optimum \(\tilde{c}_0\) for type \(\tilde{b}\) is defined similarly
(Program I (type \(\tilde{b}\))).

The payoff pair \((U_b(c_0), U_b(\tilde{c}_0))\) is called the \emph{low-information-intensity
optimum}. (By abuse of terminology, we will sometimes call the allocation \((c_0, \tilde{c}_0)\)
itself the low-information-intensity optimum.) {[}\ldots{]}

\textbf{Lemma 6.1.} \emph{The borrower can guarantee herself her low-information-intensity optimum
(\(U_b(c_0)\) if she has type \(b\), and \(\tilde{U}_b(\tilde{c}_0)\) if she has type
\(\tilde{b}\)).}

{[}\ldots{]}

\textbf{Proposition 6.1.}

\begin{enumerate}
\def\labelenumi{(\alph{enumi})}
\item
  \emph{The issuance game has a unique perfect Bayesian equilibrium if the low-information-intensity
  optimum is interim efficient (relative to prior beliefs \((\alpha, 1 - \alpha)\)). The borrower
  then obtains her low-information-intensity optimum (\(U_b(c_0)\) for type \(b\),
  \(\tilde{U}_b(\tilde{c}_0)\) for type \(\tilde{b}\)).}
\item
  \emph{If the low-information-intensity optimum is not interim efficient, then the set of
  equilibrium payoffs for the two types of borrowers is the set of payoffs that result from an
  incentive-compatible, profitable-in-expectation allocation and (weakly) Pareto-dominate the
  low-information-intensity optimum.}
\end{enumerate}

{[}\ldots{]} Let \(\tilde{c}^{SI}\) denote the \emph{symmetric information} contractual terms for
the bad borrower. It solves

\[\max_{\{\tilde{c}\}} \tilde{U}_b(\tilde{c}) \quad \text{s.t.} \quad \tilde{U}_l(\tilde{c}) \geqslant 0.\]

\textbf{Assumption 6.4 (weak monotonic profit).} \emph{Investors make a nonnegative profit if the
contractual terms are those of the bad borrower under symmetric information and the borrower is a
good one:}

\[U_l(\tilde{c}^{SI}) \geqslant 0.\]

This assumption is in general satisfied when both types are creditworthy under symmetric information
(and it is satisfied in all of our illustrations). It is always satisfied when the bad borrower is
not creditworthy under symmetric information: in that case, \(\tilde{c}^{SI}\) is the absence of
funding and thus \(U_l(\tilde{c}^{SI}) = 0\).

\emph{Definition.} The \emph{separating allocation} is the allocation \(c^*\) for the good borrower
and the symmetric-information contractual terms \(\tilde{c}^{SI}\) for the bad borrower, where
\(c^*\) maximizes the good borrower's payoff subject to the investors breaking even for the good
borrower and to the bad borrower not preferring \(c^*\) to \(\tilde{c}^{SI}\):

\[\max_{\{c\}} U_b(c) \quad \text{s.t.} \quad U_l(c) \geqslant 0, \quad \tilde{U}_b(c) \leqslant \tilde{U}_b(\tilde{c}^{SI}).\]

Note that the separating allocation is profitable type-by-type.

\textbf{Lemma 6.2.} \emph{Under the weak monotonic-profit assumption, the separating allocation is
the low-information-intensity optimum.}

\emph{Proof.} First, note that the bad borrower's symmetric-information program has the same
objective function and fewer constraints than her low-information-intensity optimum program, Program
I (type \(\tilde{b}\)). Hence,

\[\tilde{U}_b^{SI} \equiv \tilde{U}_b(\tilde{c}^{SI}) \geqslant \tilde{U}_b(\tilde{c}_0).\]

{[}\ldots{]}

\textbf{Lemma 6.3.} \emph{Under the weak monotonic-profit assumption, there exists a threshold level
\(\alpha^*\) for prior beliefs such that the low-information-intensity optimum (that is, the
separating allocation under the weak monotonic-profit assumption) is interim efficient if and only
if \(\alpha \leqslant \alpha^*\).}

\begin{quote}
\begin{enumerate}
\def\labelenumi{\arabic{enumi}.}
\setcounter{enumi}{67}
\item
  So, for example, if the borrower has initial cash \(A\), and has quasi-linear preferences, the net
  utility is equal to the gross utility minus \(A\).
\item
  Very roughly speaking, a perfect Bayesian equilibrium of a game is a set of strategies and beliefs
  such that at any stage of the game players act optimally given their beliefs at that stage (the
  equilibrium is ``perfect'') and beliefs are updated by the players according to Bayes' rule using
  equilibrium strategies and observed actions (the updating is Bayesian). {[}\ldots{]} Here
  investors are assumed to update their beliefs about the borrower's type using the borrower's
  equilibrium type-contingent contract offer and the actual contract offer. The previous condition,
  that they subscribe if and only if they expect a nonnegative profit given their updated beliefs,
  is an optimization requirement.
\end{enumerate}
\end{quote}
\stmtend
%% END-STATEMENT

\asked (i) Write the separating program: maximise the good borrower's utility over her contract, subject to investors breaking even on her, to the bad borrower not preferring the good borrower's terms to her own symmetric-information terms, and to the good borrower working. (ii) Show that the good borrower takes no reward in failure, $R_b^F=0$. (iii) Using the supplementary section, show that the separating outcome is the unique perfect Bayesian equilibrium of the issuance game iff the prior $\alpha$ is at most some $\alpha^*$.

\approach This is the pecking-order application of the privately-known-prospects model (Tirole \S6.2.2, Application 3; lecture block 3.5) with the continuous-investment technology of \S3.4. The separating program is a constrained maximisation; part (ii) is a perturbation argument: shift the reward from failure to success keeping the good type's expected reward fixed, and check that every constraint is relaxed, which, by the information-sensitivity of claims, makes investors hold the safe (debt-like) part of the cash flow. Part (iii) applies the general results on contract design by an informed party (\S6.4: Lemmas 6.2--6.3 and Proposition 6.1) after checking their key assumption, weak monotonic profit.

\notation
\begin{itemize}
  \item $I\ge0$: investment; $R^SI$ / $R^FI$: income in success / failure, $R^S>R^F\ge0$.
  \item $p_H,p_L$ (good) and $q_H,q_L$ (bad): success probabilities working / shirking, $p_H-p_L=q_H-q_L=\Delta p>0$, $q_H<p_H$.
  \item $BI$: private benefit from shirking; $A$: entrepreneur's cash; investors demand a 0 return.
  \item $R_b^S,R_b^F\ge0$: good borrower's rewards in success / failure; $\tilde U_b^{SI}$: bad borrower's symmetric-information gross utility.
  \item $\alpha$: prior probability of the good type; $\alpha^*$: threshold of part (iii).
\end{itemize}

\assume (a) Investment has positive NPV for the good type when she works,
$N\equiv p_HR^S+(1-p_H)R^F-1>0$, and finite scale, $p_H(R^S-B/\Delta p)+(1-p_H)R^F<1$ (analogues of
(3.7) and (3.9)). (b) At the optimum, investors hold a positive claim in success, $R_b^S<R^SI$. (c) Footnote 80's
formula is for $R^F=0$ ($R=R^S$); with $R^F>0$, $\tilde U_b^{SI}$ is computed the same way with the
bad type's expected income per unit.

\step{1}{Payoffs, and why the bad type works when she mimics} The good type's gross utility is
$p_HR_b^S+(1-p_H)R_b^F$. A bad type taking the same terms works, because her gain from working is
\begin{align}
  [q_HR_b^S+(1-q_H)R_b^F]-[q_LR_b^S+(1-q_L)R_b^F+BI]
  = (\Delta p)(R_b^S-R_b^F)-BI, \nonumber
\end{align}
which is the good type's incentive constraint: the two types share the same (IC$_b$). Her utility from
mimicking is therefore $q_HR_b^S+(1-q_H)R_b^F$.

\step{2}{The separating program (part (i))}
\begin{align}
  \max_{\{I,R_b^S,R_b^F\}}\;& p_HR_b^S+(1-p_H)R_b^F \label{t3e12-obj}\\
  \text{s.t.}\;& p_H(R^SI-R_b^S)+(1-p_H)(R^FI-R_b^F)\ge I-A, \tag{IR$_l$}\\
  & q_HR_b^S+(1-q_H)R_b^F\le\tilde U_b^{SI}, \tag{M}\\
  & (\Delta p)(R_b^S-R_b^F)\ge BI, \tag{IC$_b$}\\
  & 0\le R_b^S\le R^SI,\quad 0\le R_b^F\le R^FI. \nonumber
\end{align}
The book's answer lists only the rewards as choice variables, but $I$ is chosen too. With $R^F=0$,
footnote 80's value comes from the bad type's symmetric-information problem in \S3.4.2. Her net utility is the NPV $(q_HR-1)I$, increasing in $I$ when $q_HR\ge1$, so she borrows as much as the constraints allow: (IC$_b$) binds, $R_b=BI/\Delta p$, and substituting it into the break-even constraint with equality,
\begin{align}
  q_H\Bigl(RI-\frac{BI}{\Delta p}\Bigr) &= I-A \nonumber\\
  A &= I\Bigl[1-q_H\Bigl(R-\frac{B}{\Delta p}\Bigr)\Bigr] \nonumber\\
  I &= \frac{A}{1-q_H(R-B/\Delta p)}. \nonumber
\end{align}
Her gross utility is the cash plus the NPV,
\begin{align}
  \tilde U_b^{SI} &= A+(q_HR-1)\frac{A}{1-q_H(R-B/\Delta p)} \nonumber\\
  &= \Bigl[1+\frac{q_HR-1}{1-q_H(R-B/\Delta p)}\Bigr]A \qquad (q_HR\ge1). \label{t3e12-USI}
\end{align}

\answer{(i)} Program \eqref{t3e12-obj}: maximise $p_HR_b^S+(1-p_H)R_b^F$ over $\{I,R_b^S,R_b^F\}$ subject to (IR$_l$), (M) with $\tilde U_b^{SI}$ from footnote 80 (or $A$ if $q_HR<1$), (IC$_b$) and limited liability.

\step{3}{A reward in failure can be moved to success at no cost (part (ii))} Suppose an optimum has
$R_b^F>0$. Change the rewards by $\delta R^S>0$ and $\delta R^F<0$ with
\begin{align}
  p_H\,\delta R^S+(1-p_H)\,\delta R^F = 0
  \iff \delta R^F = -\frac{p_H}{1-p_H}\,\delta R^S. \label{t3e12-pert}
\end{align}
The objective \eqref{t3e12-obj} is unchanged by \eqref{t3e12-pert}. So is (IR$_l$): investors receive
expected income minus the borrower's expected reward, and both are unchanged. (IC$_b$) is relaxed:
its left side changes by $(\Delta p)(\delta R^S-\delta R^F)>0$. (M) is relaxed too; substituting
\eqref{t3e12-pert},
\begin{align}
  q_H\,\delta R^S+(1-q_H)\,\delta R^F
  &= q_H\,\delta R^S-(1-q_H)\frac{p_H}{1-p_H}\,\delta R^S \nonumber\\
  &= \frac{q_H(1-p_H)-p_H(1-q_H)}{1-p_H}\,\delta R^S \nonumber\\
  &= -\frac{p_H-q_H}{1-p_H}\,\delta R^S < 0. \label{t3e12-dM}
\end{align}
By assumption (b), a small $\delta R^S$ keeps $R_b^S\le R^SI$.

\step{4}{Slack constraints allow a strict improvement, so $R_b^F=0$} After the change, (M) and (IC$_b$)
are slack. Now raise $I$ by $\varepsilon>0$ and $R_b^S$ by $\eta=N\varepsilon/p_H>0$. Each constraint changes as follows:
(IR$_l$), left side minus right side, changes by
\begin{align}
  p_H(R^S\varepsilon-\eta)+(1-p_H)R^F\varepsilon-\varepsilon
     &= [p_HR^S+(1-p_H)R^F-1]\varepsilon-p_H\eta \nonumber\
     &= N\varepsilon-N\varepsilon = 0; \nonumber
\end{align}
(IC$_b$), left minus right, changes by $(\Delta p)\eta-B\varepsilon$; the left side of (M) by $+q_H\eta$; and the limited-liability margin $R^SI-R_b^S$ by $R^S\varepsilon-\eta=(R^S-N/p_H)\varepsilon$.
For $\varepsilon$ small enough the slack in (IC$_b$), (M) and $R_b^S<R^SI$ absorbs the last three changes, whatever their sign,
and the objective rises by $p_H\eta=N\varepsilon>0$. That contradicts optimality, so
\begin{align}
  R_b^F = 0. \qquad\blacksquare \label{t3e12-RF0}
\end{align}
The same point in one line: per unit of the bad type's mimicking utility, a reward in success buys
$p_H/q_H>1$ units of the good type's utility and a reward in failure buys $(1-p_H)/(1-q_H)<1$. Investors
receive the whole failure income $R^FI$: they hold the safe, information-insensitive part of the cash
flow, i.e.\ debt. That is the pecking order.

\answer{(ii)} $R_b^F=0$: moving reward from failure to success at constant expected cost leaves (IR$_l$) and the objective unchanged, relaxes (IC$_b$), and relaxes (M) by $(p_H-q_H)\delta R^S/(1-p_H)$; the slack then allows a larger, more profitable investment. Investors hold all of $R^FI$, a debt-like claim.

\step{5}{Uniqueness of the separating outcome iff $\alpha\le\alpha^*$ (part (iii))} First check
Assumption 6.4 (weak monotonic profit). If the bad type is not creditworthy, $\tilde c^{SI}$ is no funding
and $U_l(\tilde c^{SI})=0$. Otherwise her symmetric-information terms are $\tilde I$,
$\tilde R_b^S=B\tilde I/\Delta p$, $\tilde R_b^F=0$, with
$q_H(R^S\tilde I-\tilde R_b^S)+(1-q_H)R^F\tilde I=\tilde I-A$. A good type on those terms works (Step~1),
and the investors' profit is
\begin{align}
  U_l(\tilde c^{SI}) &= p_H(R^S\tilde I-\tilde R_b^S)+(1-p_H)R^F\tilde I-(\tilde I-A) \nonumber\\
  &= p_H(R^S\tilde I-\tilde R_b^S)+(1-p_H)R^F\tilde I-\bigl[q_H(R^S\tilde I-\tilde R_b^S)+(1-q_H)R^F\tilde I\bigr] \nonumber\\
  &= (p_H-q_H)\Bigl[\Bigl(R^S-\frac{B}{\Delta p}\Bigr)-R^F\Bigr]\tilde I \;\ge\; 0
     \quad\text{iff } R^S-\frac{B}{\Delta p}\ge R^F, \label{t3e12-wmp}
\end{align}
which always holds when $R^F=0$ (footnote 80's case). Under Assumption 6.4:
\begin{itemize}
  \item Lemma 6.2: the separating allocation (Step~2 for the good type, $\tilde c^{SI}$ for the bad) is the low-information-intensity optimum;
  \item Lemma 6.3: there is $\alpha^*$ such that this optimum is interim efficient iff $\alpha\le\alpha^*$;
  \item if $\alpha\le\alpha^*$, Proposition 6.1(a) says the issuance game has a unique PBE, and the borrower gets the low-information-intensity optimum, i.e.\ the separating outcome;
  \item if $\alpha>\alpha^*$, the optimum is not interim efficient: some incentive-compatible allocation that is profitable in expectation Pareto-dominates it, strictly for at least one type. By Proposition 6.1(b) that allocation's payoffs are equilibrium payoffs, so the separating outcome is not the only PBE.
\end{itemize}
Hence the separating outcome is the unique PBE iff $\alpha\le\alpha^*$. $\blacksquare$ When good types are
common, a pooling contract with a small cross-subsidy to the rare bad types costs the good type less than
the distortion she needs to separate, so both types can be made better off.

\answer{(iii)} Weak monotonic profit holds (always if $R^F=0$; in general iff $R^S-B/\Delta p\ge R^F$), so the separating allocation is the low-information-intensity optimum (Lemma 6.2). It is interim efficient iff $\alpha\le\alpha^*$ (Lemma 6.3), and then it is the unique PBE outcome (Proposition 6.1(a)); for $\alpha>\alpha^*$ other, Pareto-better equilibria exist (Proposition 6.1(b)).

\step{6}{Numerical check of \eqref{t3e12-USI} and \eqref{t3e12-dM}} With $q_H=0.6$, $R=3$,
$B/\Delta p=2$, $A=1$: $q_H(R-B/\Delta p)=0.6(3-2)=0.6$, so $\tilde I=1/(1-0.6)=1/0.4=2.5$ and, with $q_HR=1.8$,
$\tilde U_b^{SI}=1+(1.8-1)(2.5)=1+2=3=[1+0.8/0.4](1)$. With $p_H=0.8$ and $\delta R^S=0.01$, \eqref{t3e12-pert} gives
$\delta R^F=-(0.8/0.2)(0.01)=-0.04$, and the change in (M) is $(0.6)(0.01)+(0.4)(-0.04)=0.006-0.016=-0.01=-(0.2/0.2)(0.01)$, as \eqref{t3e12-dM} predicts.

\intuition Under asymmetric information the good type's problem is to give investors a claim whose value
depends as little as possible on her type. Income in failure is worth relatively more to the bad type (she
fails more often), so promising herself anything there invites mimicry. The good type takes her reward
only in success and gives investors all the safe salvage value $R^FI$. Investors thus hold debt-like
claims first, and that is the pecking order.

\pitfalls
\begin{itemize}
  \item Writing the mimicking constraint with the bad type shirking: because $p_H-p_L=q_H-q_L$, she faces the same (IC$_b$) and works under the good type's terms, so her mimicking utility uses $q_H$.
  \item Treating $I$ as fixed: with variable investment $I$ is a choice variable, and it is the slack created by the perturbation that lets $I$ (and the good type's utility) rise.
  \item Using the good type's symmetric-information utility in (M): the bad type compares the good type's terms with \emph{her own} symmetric-information utility $\tilde U_b^{SI}$.
\end{itemize}

\bookans Matches the book's answer (pp.~599--600): the same program (with $I$ made explicit here), the
same perturbation argument for $R_b^F=0$, and uniqueness iff $\alpha\le\alpha^*$. The book invokes
``Proposition 6.2''; with the excerpt copied in the list, the argument is Lemmas 6.2--6.3 and Proposition
6.1. The book asserts that weak monotonic profit holds; \eqref{t3e12-wmp} shows it needs
$R^S-B/\Delta p\ge R^F$, automatic in footnote 80's case $R^F=0$.

\refbib{Tirole Ch.~6, \S6.2.2 Application 3 ``Pecking-Order Hypothesis'', pp.~246--249; \S3.4.1--3.4.2,
pp.~127--128; supplementary \S6.4, pp.~264--268; Exercise 6.8, p.~278; answer pp.~599--600. Lecture
Topic~3, pecking-order slides.}

% =============================================================================
\exsec{13. Information sharing among lenders --- Tirole 6.7, pp.~276--278}
%% BEGIN-STATEMENT 13
\stmtbegin
\textbf{Source:} Tirole, Ch. 6, Exercise 6.7, pp.~276--278 (PDF 289) · Answer in book: yes, p.~599

\textbf{Exercise 6.7 (endogenous communication among lenders).} Padilla and Pagano (1997) and others
have observed that information sharing about creditworthiness is widespread among lenders (banks,
suppliers, etc.). For example, Dun \& Bradstreet Information Services, one of the leading rating
agencies, collects information from thousands of banks. Similarly, over 600,000 suppliers
communicate information about delays and defaults by their customers; and credit bureaux centralize
information about the consumer credit markets.

Padilla and Pagano (see also Pagano and Jappelli (1993) and the references therein) argue that
information sharing has both costs and benefits for the banks. By sharing information, they reduce
their differentiation and compete more with each other. But this competition protects their
borrowers' investment and therefore enhances opportunities for lending. In a sense, the ``tax rate''
(the markup that banks can charge borrowers) decreases but the ``tax base'' (the creditworthiness of
borrowers) expands. This exercise builds on the Padilla--Pagano model.

There are two periods (\(t = 1, 2\)). The discount factor between the two periods is \(\delta\). A
risk-neutral borrower protected by limited liability has no cash on hand (\(A = 0\)). Each period,
the borrower has a project with investment cost \(I\). The project delivers at the end of the period
\(R\) or 0. There is no moral hazard. The probability of success is \(p\) if the entrepreneur is
talented (which has probability \(\alpha\)), and \(q\) if she is not (which has probability
\(1 - \alpha\)). We will assume that the market rate of interest in the economy is 0, that the
lenders are risk neutral, and that only the good type is creditworthy:

\[pR > I > qR.\]

The date-1 and date-2 projects (if financed) are correlated and yield the same profit (they both
succeed or both fail).

There are \(n\) towns. Each town has one bank and one borrower. The ``local bank'' has local
expertise and thereby learns the local borrower's type; the other banks, the ``foreign banks,''
learn nothing (and therefore have beliefs \(\alpha\) that the entrepreneur is talented) at date 1.
At date 2, the foreign banks learn

\begin{itemize}
\tightlist
\item
  only whether the borrower was financed at date 1, if there is no information sharing among banks;
\item
  whether the borrower was financed at date 1 and whether she repaid (i.e., whether she was
  successful), if there is information sharing about riskiness.
\end{itemize}

In other words, information sharing is feasible on hard data (repayments), but not on soft data
(assessment of ability).

Padilla and Pagano add two twists to the model. First, banks decide \emph{ex ante} whether they will
communicate information about default and they make this decision public. Second, the borrower's
type may be endogenous (in which it refers more to an investment in the projects or industry than in
``pure talent''): at increasing and convex cost \(C(\alpha)\) \(\{C' > 0\), \(C'' > 0\),
\(C(0) = 0\), \(C'(0) = 0\), \(C'(1) = \infty\}\), the borrower develops a \(p\) project with
probability \(\alpha\) and a \(q\) project with probability \(1 - \alpha\). \(C\) can be viewed as
an investment cost and represents a nonmonetary cost borne by the borrower.

Contracts between banks and borrowers are \emph{short-term} contracts. These contracts just specify
a payment \(R_b\) for the borrower in the case of success during the period (and 0 in the case of
failure). Furthermore, in each period, banks simultaneously make take-it-or-leave-it offers to
borrowers. And at date 2, the incumbent bank (the bank that has lent at date 1) makes its offer
after the other banks.

The timing is summarized in Figure 6.6.

Figure 6.6 (timeline, left to right; p.~277, PDF 290)

{\def\LTcaptype{none} % do not increment counter
\begin{longtable}[]{@{}
  >{\raggedright\arraybackslash}p{(\linewidth - 14\tabcolsep) * \real{0.1250}}
  >{\raggedright\arraybackslash}p{(\linewidth - 14\tabcolsep) * \real{0.1250}}
  >{\raggedright\arraybackslash}p{(\linewidth - 14\tabcolsep) * \real{0.1250}}
  >{\raggedright\arraybackslash}p{(\linewidth - 14\tabcolsep) * \real{0.1250}}
  >{\raggedright\arraybackslash}p{(\linewidth - 14\tabcolsep) * \real{0.1250}}
  >{\raggedright\arraybackslash}p{(\linewidth - 14\tabcolsep) * \real{0.1250}}
  >{\raggedright\arraybackslash}p{(\linewidth - 14\tabcolsep) * \real{0.1250}}
  >{\raggedright\arraybackslash}p{(\linewidth - 14\tabcolsep) * \real{0.1250}}@{}}
\toprule\noalign{}
\begin{minipage}[b]{\linewidth}\raggedright
Date 1
\end{minipage} & \begin{minipage}[b]{\linewidth}\raggedright
\end{minipage} & \begin{minipage}[b]{\linewidth}\raggedright
\end{minipage} & \begin{minipage}[b]{\linewidth}\raggedright
\end{minipage} & \begin{minipage}[b]{\linewidth}\raggedright
\end{minipage} & \begin{minipage}[b]{\linewidth}\raggedright
Date 2
\end{minipage} & \begin{minipage}[b]{\linewidth}\raggedright
\end{minipage} & \begin{minipage}[b]{\linewidth}\raggedright
\end{minipage} \\
\midrule\noalign{}
\endhead
\bottomrule\noalign{}
\endlastfoot
The banks decide whether to share information about default. & The borrower chooses \(\alpha\). &
The borrower's type (\(p\) or \(q\)) is realized. & Banks offer contracts to borrowers. & Outcomes
(\(R\) or 0) and repayments. & Banks share information (if they agreed to do so). & Banks offer
contracts to borrowers. & Outcomes (\(R\) or 0) and repayments. \\
\end{longtable}
}

\begin{enumerate}
\def\labelenumi{(\roman{enumi})}
\item
  Suppose first that the probability \(\alpha\) of being a \(p\)-type is exogenous (there is no
  borrower investment), that
  \([\alpha p + (1 - \alpha)q]R - I + \delta(\alpha p + (1 - \alpha)q)(R - I) < 0\), and that
  \(qR - I + \delta q(R - I) < 0\). Show that the banks prefer not to share information.
\item
  Next, suppose that the borrower chooses \(\alpha\). Assuming that the two assumptions made in (i)
  still hold in the relevant range of \(\alpha\)s (for example, \(\alpha \in [0, \bar\alpha]\),
  where \(\bar\alpha\) satisfies the conditions), show that the banks choose to share information.
\end{enumerate}
\stmtend
%% END-STATEMENT

\asked (i) With the probability of talent $\alpha$ fixed, compare banks' profits with and without sharing repayment data, and show that banks prefer not to share. (ii) When the borrower chooses $\alpha$ at a convex cost, show that banks prefer to share.

\approach This is the Padilla--Pagano credit-bureau model (Tirole \S6.3; lecture block 3.4). Solve the two-period game backwards, case by case. At date 2 the incumbent bank knows the date-1 outcome; without sharing, foreign banks do not, so the incumbent holds the borrower up (an informational rent); with sharing, a success is public and Bertrand competition hands the surplus to the borrower. At date 1, the stated inequalities keep foreign banks out. Then compare banks' profits; in (ii) add the borrower's choice of $\alpha$, which depends on what she keeps at date 2.

\notation
\begin{itemize}
  \item $t=1,2$: periods; $\delta\in(0,1]$: discount factor; $I$: investment per period; $R$ / $0$: income in success / failure; market rate 0.
  \item $p$ / $q$: success probability of a talented / untalented borrower, $pR>I>qR$ (so $R>I$); $\alpha$: probability of talent; $m\equiv\alpha p+(1-\alpha)q$.
  \item Projects of the two dates succeed or fail together. $R_b^t$: borrower's reward in success at date $t$.
  \item $\pi^{ns}$, $\pi^s$: a bank's expected profit without / with information sharing; $U_b^{ns}$, $U_b^s$: borrower's utility.
  \item $C(\alpha)$: borrower's cost of reaching talent probability $\alpha$; $C'>0$, $C''>0$, $C(0)=C'(0)=0$, $C'(1)=\infty$.
\end{itemize}

\step{1}{Foreign banks do not lend at date 1 (part (i))} A foreign bank knows only $\alpha$. Its best case
is to win the borrower at date 1 and keep her at date 2 with no competition, setting $R_b=0$ both times
and relending only after a success (which guarantees a date-2 success). From a borrower of
type $\theta\in\{p,q\}$ it would earn
\begin{align}
  g(\theta) = \theta R-I+\delta\,\theta(R-I). \label{t3e13-g}
\end{align}
By the statement's assumptions, $g(q)=qR-I+\delta q(R-I)<0$ and, averaging \eqref{t3e13-g} over the prior,
\begin{align}
  \alpha g(p)+(1-\alpha)g(q) &= [\alpha p+(1-\alpha)q]R-I+\delta[\alpha p+(1-\alpha)q](R-I) \nonumber\\
  &= mR-I+\delta m(R-I)<0. \nonumber
\end{align}
It wins
every bad type (the local bank does not lend to them, Step~2) and at most a fraction $\theta\in[0,1]$
of the good types, so its profit is at most
\begin{align}
  \theta\alpha g(p)+(1-\alpha)g(q) \le \alpha g(p)+(1-\alpha)g(q) < 0. \nonumber
\end{align}
So foreign banks make no date-1 offers.

\step{2}{The local bank lends only to the good type, at $R_b^1=0$} The local bank knows the type. Even as a
two-period monopolist it would earn $g(q)<0$ on a bad type, so it does not lend to her. With no foreign
competition at date 1 it makes a take-it-or-leave-it offer to the good type: $R_b^1=0$, date-1 profit $pR-I$.

\step{3}{Date 2 without information sharing} Foreign banks see only that the borrower was financed at
date 1, which reveals a good type; they assign probability $p$ to success but do not know the date-1
outcome. The incumbent knows it, hence knows the date-2 outcome. If a foreign bank offers $R_b<R-I$, the
incumbent, moving last, outbids it ($R_b+\varepsilon$) after a success and stays out after a failure. The
foreign bank wins only failing borrowers and loses $I$, so it does not bid. The incumbent then offers
$R_b^2=0$ after a success and does not lend after a failure. Each bank's profit (banks earn nothing in
other towns) and the borrower's utility are
\begin{align}
  \pi^{ns} &= \alpha[pR-I+\delta p(R-I)], \label{t3e13-pins}\\
  U_b^{ns} &= 0. \label{t3e13-Uns}
\end{align}

\step{4}{Date 2 with information sharing} All banks now see the date-1 repayment, so after a success
everyone knows the date-2 project succeeds. With $n\ge2$ banks this is Bertrand competition: the reward
is bid up to $R_b^2=R-I$ and date-2 profit is $0$. After a failure nobody lends. At date 1 the foreign banks
still do not lend: with no date-2 rent their best case is $mR-I$, and since $R>I$ makes $\delta m(R-I)>0$,
\begin{align}
  mR-I < mR-I+\delta m(R-I) < 0 \nonumber
\end{align}
by the first assumption. The local bank still lends
to the good type only, at $R_b^1=0$. Hence
\begin{align}
  \pi^{s} &= \alpha(pR-I), \label{t3e13-pis}\\
  U_b^{s} &= \delta\alpha p(R-I). \label{t3e13-Us}
\end{align}

\step{5}{Banks prefer not to share} Subtracting \eqref{t3e13-pis} from \eqref{t3e13-pins},
\begin{align}
  \pi^{ns}-\pi^{s} = \alpha\,\delta p(R-I) > 0. \qquad\blacksquare \label{t3e13-diff}
\end{align}
With $\alpha$ fixed, sharing only hands the incumbent's date-2 informational rent to the borrower.

\answer{(i)} Without sharing $\pi^{ns}=\alpha[pR-I+\delta p(R-I)]$ and $U_b^{ns}=0$; with sharing $\pi^s=\alpha(pR-I)$ and $U_b^s=\delta\alpha p(R-I)$. Banks lose $\alpha\delta p(R-I)>0$ by sharing, so they do not share.

\step{6}{Endogenous $\alpha$: banks share (part (ii))} The banks' decision is public and is made before
the borrower chooses $\alpha$; her net utility is the gross utility of Steps~3--4 minus $C(\alpha)$.
Without sharing, \eqref{t3e13-Uns} gives net utility $-C(\alpha)$, maximised at
\begin{align}
  \alpha^{ns} = 0 \quad\Rightarrow\quad \pi^{ns} = 0\cdot[pR-I+\delta p(R-I)] = 0. \label{t3e13-ans}
\end{align}
With sharing, by \eqref{t3e13-Us} she solves $\max_\alpha\{\delta\alpha p(R-I)-C(\alpha)\}$. The derivative of the objective is $\delta p(R-I)-C'(\alpha)$, which is strictly decreasing ($C''>0$), positive at $\alpha=0$ (since $C'(0)=0<\delta p(R-I)$) and negative near $\alpha=1$ (since $C'(1)=\infty$). So the objective is strictly concave and its maximum is the unique interior zero of the derivative:
\begin{align}
  C'(\alpha^s) = \delta p(R-I), \qquad \alpha^s\in(0,1). \label{t3e13-foc}
\end{align}
Assuming, as the statement allows, that the conditions of (i) hold at $\alpha^s$, \eqref{t3e13-pis} gives
\begin{align}
  \pi^s = \alpha^s(pR-I) > 0 = \pi^{ns}, \label{t3e13-share}
\end{align}
so the banks choose to share information. $\blacksquare$ The borrower gains too:
$U_b^s=\delta\alpha^sp(R-I)-C(\alpha^s)>0$, because $C'(0)=0$ makes small $\alpha$ strictly profitable.

\answer{(ii)} Without sharing the borrower expects no rent, invests nothing ($\alpha^{ns}=0$) and banks earn $0$; with sharing she invests $\alpha^s$ with $C'(\alpha^s)=\delta p(R-I)$, and banks earn $\alpha^s(pR-I)>0$. Banks choose to share.

\step{7}{Numerical illustration} Take $p=0.8$, $q=0.2$, $R=10$, $I=6$, $\delta=0.9$, so
\begin{align}
  g(p) &= (0.8)(10)-6+(0.9)(0.8)(10-6) = 2+2.88 = 4.88, \nonumber\\
  g(q) &= (0.2)(10)-6+(0.9)(0.2)(4) = 2-6+0.72 = -3.28<0. \nonumber
\end{align}
The first assumption reads $m[R+\delta(R-I)]<I$, i.e.\ $m(10+3.6)<6$, so $m<6/13.6=0.441176$; with $m=0.2+0.6\alpha$ this is $\alpha<(0.441176-0.2)/0.6=0.402$. With
$\alpha=0.3$: $\pi^{ns}=0.3(4.88)=1.464$, $\pi^s=0.3(8-6)=0.6$, $U_b^s=0.9(0.3)(0.8)(4)=0.864$. For (ii) take
$C(\alpha)=6[-\alpha-\ln(1-\alpha)]$, which satisfies the conditions ($C(0)=0$), with
\begin{align}
  C'(\alpha)=6\Bigl[-1+\frac{1}{1-\alpha}\Bigr]=\frac{6\alpha}{1-\alpha}, \qquad C''(\alpha)=\frac{6}{(1-\alpha)^2}>0, \nonumber
\end{align}
so $C'(0)=0$ and $C'(\alpha)\to\infty$ as $\alpha\to1$.
Then \eqref{t3e13-foc} gives
\begin{align}
  \frac{6\alpha}{1-\alpha}&=\delta p(R-I)=(0.9)(0.8)(4)=2.88 \nonumber\
  6\alpha&=2.88-2.88\alpha \nonumber\
  \alpha^s&=\frac{2.88}{8.88}=0.324324<0.402, \nonumber
\end{align}
$\pi^s=0.324324(2)=0.648649>0$ and $U_b^s=2.88(0.324324)-6[-0.324324-\ln(0.675676)]=0.934054-0.406307=0.527747$.

\intuition Sharing default data is a trade-off between the size of the banks' markup and the size of the
market. With talent given, it only erodes the incumbent's hold-up rent (the markup falls, the base is
fixed), so banks keep data private. When talent is an investment, a borrower who expects to be held up
does not invest, and the base collapses to zero. Sharing commits banks to compete after a success,
protects the borrower's return on investing, and creates borrowers worth lending to.

\pitfalls
\begin{itemize}
  \item Letting a foreign bank compete at date 2 without sharing: it cannot tell successful from failed borrowers and, with the incumbent moving last, it wins only the failures (adverse selection against the outsider).
  \item Forgetting the correlation of the two projects: after a date-1 success the date-2 project succeeds for sure, which is why the date-2 rent is $\delta p(R-I)$ per good borrower.
  \item Comparing banks' profits in (ii) at the same $\alpha$: the point of (ii) is that $\alpha$ itself responds to the sharing decision ($\alpha^{ns}=0<\alpha^s$).
\end{itemize}

\bookans Matches the book's answer (p.~599): $\pi^{ns}=\alpha[pR-I+\delta p(R-I)]$, $U_b^{ns}=0$,
$\pi^s=\alpha(pR-I)$, $U_b^s=\delta\alpha p(R-I)$; with endogenous $\alpha$, $\alpha^{ns}=0$ and
$C'(\alpha^*)=\delta p(R-I)$, so sharing is chosen.

\refbib{Tirole Ch.~6, \S6.3 Application 4 (disclosure of past repayments, Padilla--Pagano), pp.~249--250;
Exercise 6.7, pp.~276--278; answer p.~599. Lecture Topic~3, credit-bureau slides.}

% =============================================================================
\exsec{14. Leland--Pyle with mean-variance preferences --- Amaro de Matos 3.2.6 \#2, p.~95}
%% BEGIN-STATEMENT 14
\stmtbegin
\textbf{Source:} Amaro de Matos, Ch. 3, Exercise 3.2.6 \#2, p.~95 (PDF 113) · Answer in book: no
(p.~260 only cites page 377 of Leland and Pyle 1977)

In the context of Leland and Pyle (1977), suppose that the expected utility of the manager is given
by \(E[U(W_1)] = G[E(W_1) - \frac{b}{2}\sigma^2_{W_1}]\), where \(G\) is an increasing function,
\(\sigma^2_{W_1}\) is the variance of the wealth At time 1, and the risk adjustment parameter
\(\lambda\) is given by
\(\lambda = \frac{E[M] - (1 + r)V_M}{\sigma^2_M}\operatorname{cov}(\varepsilon, M)\). In this case,
under what conditions does the level of debt increase with the proportion \(\alpha\) of the equity
owned by the managers?

\paragraph{Context needed to solve it}

\emph{From Amaro de Matos §3.2, p.~78 (PDF 96):}

\textbf{Assumption 3.2.1.} \emph{Managers have better information about the value of the firm than
other agents in the economy.}

\emph{From Amaro de Matos §3.2.2 ``When Managers Are Also Investors'', pp.~82--85 (PDF 100--103):}

Leland and Pyle (1977) developed a model in the same signaling context as Ross, in which risk averse
managers are also investors in the firm. Consider entrepreneurs starting a project with a random
return \(X = \overline{X} + \epsilon\), with \(E(\epsilon) = 0\), needing a capital investment
\(K\). Here, these entrepreneurs, who are going to manage the project, are the only individuals who
know the expected return \(\overline{X} = E(X)\). At time 0 they sell a fraction \((1 - \alpha)\) of
the project and retain a fraction \(\alpha\). The uninformed investors use this proportion as a
signal of the value of the firm. Underlying this idea is, clearly, Assumption 3.2.1. As the
management better knows the value of the firm, the fact that it holds a reasonable part of the
equity is a good signal, whereas no participation will naturally be considered as a bad signal. As
the managers are risk averse, they would like to diversify but cannot, since they have to hold on to
their shares. That is their cost of signaling, and that is why the signal is credible. Managers of
good projects are then easily identified with this procedure because their claim on the firm is
worth more. That is why they can afford to signal in that way and hold onto a greater portion of the
shares.

In this model, Assumption 3.2.1 that managers have superior information takes the simple form

\textbf{Assumption 3.2.4.} \emph{Only the management knows the value of the expected return \(X\).}

Uninformed investors will have to \emph{infer} in equilibrium the expected return through a function
of \(\alpha\), say, \(X(\alpha)\). Consequently, the \emph{perceived} value of the firm, say, a
function \(V(\alpha)\), will also depend on \(\alpha\). As in Ross's model, Leland and Pyle work in
a two-time framework. As in equation (3.8), the perceived value of the firm will here be given by

\[V(\alpha) = \frac{X(\alpha) - \lambda}{1 + r}, \qquad (3.15)\]

where \(\lambda\) is a risk adjustment parameter and \(r\) is the risk-free interest rate.

At time 0, managers have their initial wealth \(W_\circ\). Additionally they may issue risk-free
debt in the amount of \(B\) and sell a fraction \((1 - \alpha)\) of the perceived value of the
project. The value of such a fraction is worth \((1 - \alpha)[V(\alpha) - B]\) at time 0, since debt
must be paid from the proceeds of the project. Thus, entrepreneurs initially have

\[W_\circ + B + (1 - \alpha)[V(\alpha) - B].\]

This amount is assumed to be invested in three different assets, namely the project, the market
portfolio, and a risk-free asset. The decision variables will be the proportion \(\alpha\) of the
equity of the firm held by the entrepreneurs and the proportion \(\beta\) of the market portfolio
held. The rest of their wealth (say, \(Y\)) will be invested in a risk-free asset. The optimal way
to make this allocation is constrained by

\[W_\circ + B + (1 - \alpha)[V(\alpha) - B] = K + Y + \beta V_M,\]

where \(V_M\) is taken as the value of the market portfolio. This equation is the \emph{budget
constraint} faced by the management when deciding how to allocate its wealth. Using equation (3.15)
it follows that

\[Y = W_\circ + B - K - \beta V_M + \frac{1 - \alpha}{1 + r}[X(\alpha) - \lambda - (1 + r)B]. \qquad (3.16)\]

The objective of the managers is to maximize the expected utility of their random wealth at the end
of the process (i.e., at time 1). If the return on the market portfolio is denoted by
\(\widetilde{M}\), the wealth of the management at that time is the random variable

\[\widetilde{W}_1 = \alpha[X - B(1 + r)] + \beta\widetilde{M} + (1 + r)Y.\]

Substituting the value of \(Y\) by the expression in (3.16) (that is, incorporating in the random
wealth the budget constraint with respect to which the utility maximization must be subject), it
follows that

\[\widetilde{W}_1 = \alpha[\overline{X} + \epsilon - X(\alpha) + \lambda] + \beta[\widetilde{M} - (1 + r)V_M] + (W_\circ - K)(1 + r) + (X(\alpha) - \lambda).\]

Assuming that the management has a utility function \(U\), its problem becomes

\[\max_{(\alpha, \beta) \in [0,1] \times [0,1]} E[U(\widetilde{W}_1)], \qquad (3.17)\]

for which there will be a solution \(\alpha^* \equiv \alpha^*(\overline{X})\) and
\(\beta^* \equiv \beta^*(\overline{X})\). As in other signaling models, in equilibrium the values of
\(\alpha^*\) and \(\beta^*\) must reveal the true expected return, that is to say, the equilibrium
condition

\[\overline{X}[\alpha^*(\overline{X})] = \overline{X}\]

must be imposed on the solution of the problem. There are two necessary first-order conditions to be
satisfied in order to solve (3.17): one with respect to \(\alpha\) and the other with respect to
\(\beta\). Denoting by \(X(\alpha)_\alpha\) the derivative \(dX(\alpha)/d\alpha\), the first-order
condition reads

\[(1 - \alpha)X(\alpha)_\alpha = -\frac{E[U'(\widetilde{W}_1)(\epsilon + \lambda)]}{E[U'(\widetilde{W}_1)]}, \qquad (3.18)\]

and the second is

\[E\{U'(\widetilde{W}_1)[\widetilde{M} - (1 + r)V_M]\} = 0. \qquad (3.19)\]

Solving (3.19) for \(\beta\) and substituting this value of \(\beta\) into (3.18), one gets a
differential equation for \(X(\alpha)\) whose solution does not depend explicitly on \(\beta\) and
is uniquely defined up to a constant of integration.

The following two important results are proved by Leland and Pyle:

\textbf{Proposition 3.2.1 (Leland and Pyle, 1977).} \emph{\(X(\alpha)\) is strictly increasing in
\(\alpha\) if and only if the management demand for equity is normal.}\textsuperscript{15}

\textbf{Proposition 3.2.2 (Leland and Pyle, 1977).} \emph{In equilibrium signaling, managers with
normal demand have a value of \(\alpha\) higher than the proportion they would invest if they could
costlessly communicate the true value of \(\overline{X}\).}

This last result may be seen as a welfare cost result. In fact, the cost of signaling reflected in
the overparticipation of the management in the equity of the company may be used as a criterion to
choose the constant of integration defining \(X(\alpha)\) as the solution of the differential
equation previously described. The constant will be chosen so as to minimize this signaling cost.

The relevant issue here is how these results relate to the debt policy followed by the firms. All
other factors being fixed, including the amount invested by the management in equity, in most cases
the level of debt increases with the proportion \(\alpha\) of the equity owned by managers. From
Proposition 3.2.1, it follows that \(X(\alpha)\) increases with debt; therefore equation (3.15)
implies that the perceived value \(V(\alpha)\) increases with the debt level \(B\), just as in
Ross's model. {[}\ldots{]}

\emph{From Amaro de Matos, Notes to Chapter 3, note 15, p.~280 (PDF 298):}

\begin{enumerate}
\def\labelenumi{\arabic{enumi}.}
\setcounter{enumi}{14}
\tightlist
\item
  The demand for an asset is said to be normal if, in a portfolio choice situation without
  signaling, the demand is a decreasing function of price.
\end{enumerate}
\stmtend
%% END-STATEMENT

\asked (a) In Leland and Pyle's model, with mean-variance utility $G[E(W_1)-\frac b2\sigma^2_{W_1}]$ and the CAPM risk adjustment $\lambda$ given in the statement, find the conditions under which the firm's debt $B$ rises with the fraction $\alpha$ of equity the managers retain.

\approach This is Leland--Pyle signalling by retained ownership (Amaro de Matos \S3.2.2; lecture block 3.9). Retaining a stake is costly for a risk-averse manager because she cannot diversify, so the stake is a credible signal of $\overline X$. With mean-variance utility the two first-order conditions (3.18)--(3.19) can be solved explicitly: the CAPM term $\lambda$ cancels and only idiosyncratic risk is costly, giving a closed-form signalling schedule $X(\alpha)$. Debt does not enter the managers' wealth, so it moves with $\alpha$ only through the budget constraint with the managers' own investment held fixed (the book's ``all other factors being fixed''); differentiating $B(\alpha)$ gives the condition.

\notation
\begin{itemize}
  \item $X=\overline X+\varepsilon$: project payoff at time 1, $E(\varepsilon)=0$, $\sigma_\varepsilon^2=\operatorname{var}(\varepsilon)$; $K$: capital required.
  \item $\alpha\in[0,1)$: fraction retained by managers; $X(\alpha)$: payoff inferred by the market; $V(\alpha)=[X(\alpha)-\lambda]/(1+r)$: perceived value (3.15).
  \item $B$: riskless debt (face $(1+r)B$); $r$: riskless rate; $\widetilde M$, $V_M$, $\sigma_M^2$: market payoff, value, variance; $\beta$: fraction of the market held.
  \item $W_0$: initial wealth; $Y$: riskless holding; $\widetilde W_1$: time-1 wealth; $b>0$: risk-aversion coefficient.
  \item $\lambda=\frac{E[M]-(1+r)V_M}{\sigma_M^2}\operatorname{cov}(\varepsilon,M)$: CAPM risk adjustment (statement).
  \item $\sigma_u^2\equiv\sigma_\varepsilon^2-\operatorname{cov}(\varepsilon,M)^2/\sigma_M^2\ge0$: idiosyncratic variance of the project.
  \item $E_m\equiv K-B-(1-\alpha)[V(\alpha)-B]$: managers' own money put into the firm's equity.
\end{itemize}

\assume The book gives no worked answer (it cites Leland and Pyle 1977, p.~377), so everything below is
derived here from the \S3.2.2 excerpt. (a) $\sigma_u^2>0$ and interior $\alpha\in(0,1)$. (b) Debt is riskless,
$(1+r)B\le X$ in every state. (c) ``All other factors being fixed, including the amount invested by the
management in equity'' (\S3.2.2) is read from the budget constraint
$W_0+B+(1-\alpha)[V(\alpha)-B]=K+Y+\beta V_M$: the managers' own money in the project,
$E_m=W_0-Y-\beta V_M$, is held fixed across firms. A Remark covers the other reading.

\step{1}{Debt drops out of the managers' wealth} The book's expression for $\widetilde W_1$, obtained by
substituting (3.16), is
\begin{align}
  \widetilde W_1 = \alpha[\overline X+\varepsilon-X(\alpha)+\lambda]+\beta[\widetilde M-(1+r)V_M]+(W_0-K)(1+r)+X(\alpha)-\lambda, \label{t3e14-W1}
\end{align}
and $B$ does not appear: the $-\alpha(1+r)B$ on the firm's equity, the $(1+r)B$ of proceeds and the
$-(1-\alpha)(1+r)B$ on the shares sold cancel. Debt is irrelevant to the managers (MM), so it is pinned
down only by holding $E_m$ fixed (Step~6).

\step{2}{Mean and variance of wealth} From \eqref{t3e14-W1}, the only random terms are $\alpha\varepsilon$ and $\beta\widetilde M$; taking expectations ($E\varepsilon=0$) and using $\operatorname{var}(\alpha\varepsilon+\beta\widetilde M)=\alpha^2\sigma_\varepsilon^2+\beta^2\sigma_M^2+2\alpha\beta\operatorname{cov}(\varepsilon,M)$,
\begin{align}
  E(\widetilde W_1) &= \alpha[\overline X-X(\alpha)+\lambda]+\beta[E(M)-(1+r)V_M]+(W_0-K)(1+r)+X(\alpha)-\lambda, \label{t3e14-EW}\\
  \sigma^2_{W_1} &= \alpha^2\sigma_\varepsilon^2+\beta^2\sigma_M^2+2\alpha\beta\operatorname{cov}(\varepsilon,M). \label{t3e14-VW}
\end{align}
Since $G$ is increasing, the managers maximise $J\equiv E(\widetilde W_1)-\frac b2\sigma^2_{W_1}$.

\step{3}{First-order condition in $\beta$ (the analogue of (3.19))} Differentiating \eqref{t3e14-EW}--\eqref{t3e14-VW} in $\beta$: $\partial E(\widetilde W_1)/\partial\beta=E(M)-(1+r)V_M$ and $\partial\sigma^2_{W_1}/\partial\beta=2\beta\sigma_M^2+2\alpha\operatorname{cov}(\varepsilon,M)$, so $\partial J/\partial\beta=0$ reads
\begin{align}
  E(M)-(1+r)V_M-b[\beta\sigma_M^2+\alpha\operatorname{cov}(\varepsilon,M)] &= 0, \nonumber\\
  b\beta\operatorname{cov}(\varepsilon,M) &= \frac{E(M)-(1+r)V_M}{\sigma_M^2}\operatorname{cov}(\varepsilon,M)-b\alpha\frac{\operatorname{cov}(\varepsilon,M)^2}{\sigma_M^2} \nonumber\\
  &= \lambda-b\alpha\frac{\operatorname{cov}(\varepsilon,M)^2}{\sigma_M^2}, \label{t3e14-beta}
\end{align}
where the second line multiplies the first by $\operatorname{cov}(\varepsilon,M)/\sigma_M^2$ and rearranges, and the
last uses the statement's $\lambda$.

\step{4}{First-order condition in $\alpha$ (the analogue of (3.18))} Differentiating \eqref{t3e14-EW} in $\alpha$ (product rule on $\alpha[\cdot]$, plus the last term $X(\alpha)$),
\begin{align}
  \frac{\partial E(\widetilde W_1)}{\partial\alpha} = [\overline X-X(\alpha)+\lambda]-\alpha X'(\alpha)+X'(\alpha) = \overline X-X(\alpha)+\lambda+(1-\alpha)X'(\alpha), \nonumber
\end{align}
and $\partial\sigma^2_{W_1}/\partial\alpha=2\alpha\sigma_\varepsilon^2+2\beta\operatorname{cov}(\varepsilon,M)$, so $\partial J/\partial\alpha=0$ reads
\begin{align}
  \overline X-X(\alpha)+\lambda+(1-\alpha)X'(\alpha)-b[\alpha\sigma_\varepsilon^2+\beta\operatorname{cov}(\varepsilon,M)] = 0. \nonumber
\end{align}
Substituting \eqref{t3e14-beta} for $b\beta\operatorname{cov}(\varepsilon,M)$ and imposing the equilibrium condition $X(\alpha^*)=\overline X$ (so the first two terms cancel):
\begin{align}
  \lambda+(1-\alpha)X'(\alpha)-b\alpha\sigma_\varepsilon^2-\Bigl[\lambda-b\alpha\frac{\operatorname{cov}(\varepsilon,M)^2}{\sigma_M^2}\Bigr] &= 0 \nonumber\\
  (1-\alpha)X'(\alpha) &= b\alpha\Bigl[\sigma_\varepsilon^2-\frac{\operatorname{cov}(\varepsilon,M)^2}{\sigma_M^2}\Bigr] \nonumber\\
  (1-\alpha)X'(\alpha) &= b\alpha\sigma_u^2. \label{t3e14-ode}
\end{align}
The CAPM $\lambda$ prices the systematic risk fairly, so it cancels; only idiosyncratic risk is costly to
retain.

\step{5}{The signalling schedule} Dividing \eqref{t3e14-ode} by $1-\alpha>0$ gives $X'$. To integrate it, note
\begin{align}
  \frac{d}{d\alpha}[\alpha+\ln(1-\alpha)] = 1-\frac{1}{1-\alpha} = \frac{(1-\alpha)-1}{1-\alpha} = -\frac{\alpha}{1-\alpha}, \nonumber
\end{align}
so integrating from 0 to $\alpha$ (where $\alpha+\ln(1-\alpha)=0$ at $\alpha=0$), and using (3.15) for $V'$:
\begin{align}
  X'(\alpha) &= b\sigma_u^2\frac{\alpha}{1-\alpha} > 0, \label{t3e14-Xp}\\
  X(\alpha) &= X(0)-b\sigma_u^2[\alpha+\ln(1-\alpha)], \label{t3e14-X}\\
  V'(\alpha) &= \frac{X'(\alpha)}{1+r} = \frac{b\sigma_u^2}{1+r}\,\frac{\alpha}{1-\alpha}. \label{t3e14-Vp}
\end{align}
Normal demand holds: without signalling (price $V$ given, so the term $(1-\alpha)X'$ is absent and the market's $X(\alpha)-\lambda$ is replaced by $(1+r)V$), the same steps give
$\overline X-\lambda-(1+r)V=b\alpha\sigma_u^2$, so $\alpha=[\overline X-\lambda-(1+r)V]/(b\sigma_u^2)$ falls
with the price. Proposition 3.2.1 then agrees with \eqref{t3e14-Xp}. With costless communication
($X'\equiv0$), \eqref{t3e14-ode} gives $\alpha=0$, so every $\alpha>0$ is over-retention (Proposition 3.2.2).
The integration constant is set by the lowest type, which does not signal: $X(0)=\overline X_{\min}$.

\step{6}{Debt as a function of $\alpha$} Holding $E_m$ fixed in its definition and solving for $B$:
\begin{align}
  E_m &= K-B-(1-\alpha)[V(\alpha)-B] \nonumber\\
      &= K-B-(1-\alpha)V(\alpha)+(1-\alpha)B = K-\alpha B-(1-\alpha)V(\alpha), \nonumber\\
  \alpha B &= K-E_m-(1-\alpha)V(\alpha), \nonumber\\
  B(\alpha) &= \frac{K-E_m-(1-\alpha)V(\alpha)}{\alpha}. \label{t3e14-B}
\end{align}

\step{7}{When does debt increase with $\alpha$?} Differentiating \eqref{t3e14-B} with the quotient rule (the numerator's derivative is $V(\alpha)-(1-\alpha)V'(\alpha)$, the denominator's is 1),
\begin{align}
  B'(\alpha) &= \frac{[V(\alpha)-(1-\alpha)V'(\alpha)]\alpha-[K-E_m-(1-\alpha)V(\alpha)]}{\alpha^2} \nonumber\\
  &= \frac{\alpha V(\alpha)+(1-\alpha)V(\alpha)-(K-E_m)-\alpha(1-\alpha)V'(\alpha)}{\alpha^2} \nonumber\\
  &= \frac{V(\alpha)-(K-E_m)-\alpha(1-\alpha)V'(\alpha)}{\alpha^2} \nonumber\\
  &= \frac{1}{\alpha^2}\Bigl[V(\alpha)-K+E_m-\alpha(1-\alpha)\Bigl(\frac{b\sigma_u^2}{1+r}\frac{\alpha}{1-\alpha}\Bigr)\Bigr] && \text{(substituting \eqref{t3e14-Vp})} \nonumber\\
  &= \frac{1}{\alpha^2}\Bigl[V(\alpha)-K+E_m-\frac{b\sigma_u^2\alpha^2}{1+r}\Bigr]. \label{t3e14-Bp}
\end{align}
So debt rises with managerial ownership iff
\begin{align}
  V(\alpha)-K+E_m > \frac{b\,\sigma_u^2\,\alpha^2}{1+r}
  \iff V(\alpha)-B(\alpha) > (1-\alpha)V'(\alpha) = \frac{b\,\sigma_u^2\,\alpha}{1+r}, \label{t3e14-cond}
\end{align}
where the equivalence uses $V-B=(V-K+E_m)/\alpha$ from \eqref{t3e14-B} and \eqref{t3e14-Vp}. In words: the
NPV $V-K$ plus the managers' own equity contribution must exceed $b\sigma_u^2\alpha^2/(1+r)$, the risk
term from the idiosyncratic risk they keep. Equivalently, the market value of equity must exceed the
marginal cost to the managers of retaining one more unit of the firm, $b\sigma_u^2\alpha/(1+r)$. This holds
``in most cases'' (risk aversion, idiosyncratic risk or $\alpha$ small relative to the project's value). It
fails when managers keep a large, very risky stake in a firm with little equity value. When
\eqref{t3e14-cond} holds, $X(\alpha)$ and $V(\alpha)$ rise with $B$ along the schedule, as the book says.

\answer{(a)} With the managers' own equity contribution $E_m$ fixed, $B'(\alpha)>0$ iff $V(\alpha)-K+E_m>b\sigma_u^2\alpha^2/(1+r)$, equivalently iff the market value of equity exceeds $b\sigma_u^2\alpha/(1+r)$, where $\sigma_u^2=\sigma_\varepsilon^2-\operatorname{cov}(\varepsilon,M)^2/\sigma_M^2$ and the schedule is $X(\alpha)=X(0)-b\sigma_u^2[\alpha+\ln(1-\alpha)]$. This holds unless managers keep a large, very risky (idiosyncratically) stake in a firm with little equity value. (If instead the value of the managers' stake is held fixed, debt always rises with $\alpha$.)

\emph{Remark (other reading).} If what is held fixed is the market value of the managers' stake,
$\alpha[V(\alpha)-B]=e>0$, then $B=V(\alpha)-e/\alpha$ and $B'(\alpha)=V'(\alpha)+e/\alpha^2>0$ for every
$\alpha$ by \eqref{t3e14-Vp}: debt always rises with $\alpha$, and the only condition is
normal demand, which holds under mean-variance preferences.

\step{8}{Numerical check} Take $b=2$, $\sigma_\varepsilon^2=0.5$, $\sigma_M^2=0.04$,
$\operatorname{cov}(\varepsilon,M)=0.06$, $E(M)-(1+r)V_M=0.1$, $r=0.05$, $X(0)=10$, $K=8$, $E_m=1$. Then
$\lambda=(0.1/0.04)(0.06)=0.15$ and $\sigma_u^2=0.5-0.06^2/0.04=0.5-0.09=0.41$, so $b\sigma_u^2=0.82$. At $\alpha=0.5$: \eqref{t3e14-X} gives
$X=10-0.82(0.5+\ln0.5)=10-0.82(0.5-0.693147)=10+0.158381=10.158381$ and $V=(10.158381-0.15)/1.05=10.008381/1.05=9.531791$; \eqref{t3e14-B} gives
$B=(8-1-0.5\times9.531791)/0.5=(7-4.765896)/0.5=4.468209$; and \eqref{t3e14-Bp}, with $b\sigma_u^2\alpha^2/(1+r)=0.82(0.25)/1.05=0.195238$, gives
$B'=(9.531791-7-0.195238)/0.25=(2.531791-0.195238)/0.25=9.346212>0$. A type with $\overline X=X(0.5)$ that faces the schedule
\eqref{t3e14-X} has $\partial J/\partial\alpha=0$ at $\alpha=0.5$ (the check script verifies this by finite
differences).

\intuition Retaining shares is a credible signal because it is costly in risk, and more costly for a
manager whose firm is worth less. Mean-variance utility with CAPM pricing makes the cost explicit:
only idiosyncratic variance matters, and the signalling slope is $b\sigma_u^2\alpha/(1-\alpha)$. Debt does
not create value here (it drops out of wealth). It moves with $\alpha$ only because, with the managers'
own contribution fixed, a larger retained fraction has to be financed with less outside equity and more
debt. This is ``in most cases'', not always, because a higher-$\alpha$ firm is also worth more, and the
extra value lowers the debt needed.

\pitfalls
\begin{itemize}
  \item Keeping total variance $\sigma_\varepsilon^2$ in the schedule: the CAPM adjustment $\lambda$ prices systematic risk fairly and cancels, so only the idiosyncratic part $\sigma_u^2$ is a signalling cost.
  \item Treating debt as a value-creating signal here: $B$ drops out of the managers' wealth \eqref{t3e14-W1}; the link between $B$ and $\alpha$ comes only from what is held fixed in the budget constraint.
  \item Forgetting the equilibrium condition $X(\alpha^*)=\overline X$ when simplifying the first-order condition in $\alpha$; without it the terms $\overline X-X(\alpha)$ do not cancel and no differential equation results.
\end{itemize}

\bookans No worked answer in the book: Appendix C.3.2 (p.~260) only refers to p.~377 of Leland and Pyle
(1977). The schedule \eqref{t3e14-X} and the condition \eqref{t3e14-cond} are derived here and rest on reading
(c) of what is held fixed; the Remark gives the result under the other reading.

\refbib{Amaro de Matos Ch.~3, \S3.2.2 ``When Managers Are Also Investors'', pp.~82--85 (equations
(3.15)--(3.19), Propositions 3.2.1--3.2.2); Appendix C.3.2, p.~260. Lecture Topic~3, Leland--Pyle slides.}

% =============================================================================
\exsec{15. Assets in place and the negative stock-price reaction with a continuum of types --- Tirole 6.5, p.~275}
%% BEGIN-STATEMENT 15
\stmtbegin
\textbf{Source:} Tirole, Ch. 6, Exercise 6.5, p.~275 (PDF 288) · Answer in book: yes, p.~597

Consider the privately-known-prospects model of Application 2 in Section 6.2.2, but with a continuum
of types. The entrepreneur already owns a project, which with probability \(p\) yields profit \(R\)
and probability \(1 - p\) profit 0. The probability \(p\) is private information of the borrower.
From the point of view of the investors, \(p\) is drawn from cumulative distribution \(F(p)\) with
continuous density \(f(p) > 0\) on some interval \([\underline{p}, \bar{p}]\). Assume that the
distribution has monotone hazard rates:

\[\frac{f(p)}{F(p)} \quad \text{is decreasing in } p\]

and

\[\frac{f(p)}{1 - F(p)} \quad \text{is increasing in } p.\]

(This assumption, which is satisfied by most usual distributions, is known to imply that the
truncated means \(m^-(p)\) and \(m^+(p)\) have slope less than 1:

\[0 < (m^-(p))' \equiv \frac{d}{dp}[E(\tilde{p} \mid \tilde{p} \leqslant p)] \leqslant 1\]

and

\[0 < (m^+(p))' \equiv \frac{d}{dp}[E(\tilde{p} \mid \tilde{p} \geqslant p)] \leqslant 1\]

(see, for example, An 1998).)

The model is otherwise as in Section 6.2.2. A seasoned offering may be motivated by a profitable
deepening investment: at cost \(I\), the probability of success can be raised by an amount \(\tau\)
such that

\[\tau R > I\]

(of course, we need to assume that \(\bar{p} + \tau \leqslant 1\)). The entrepreneur has no cash on
hand, is risk neutral, and is protected by limited liability. The investors are risk neutral and
demand a rate of return equal to 0.

\begin{enumerate}
\def\labelenumi{(\roman{enumi})}
\item
  Show that in any equilibrium, only types \(p \leqslant p^*\), for some cutoff \(p^*\), raise funds
  and finance the deepening investment.
\item
  Show that \(p^* > \underline{p}\) and that if \(p^* < \bar{p}\), then
\end{enumerate}

\[\frac{\tau R}{I} = \frac{p^* + \tau}{m^-(p^*) + \tau}.\]

Show that if the benefits from investment are ``not too large,'' in that

\[\frac{\tau R}{I} < \frac{\bar{p} + \tau}{E[p] + \tau},\]

then indeed \(p^* < \bar{p}\).

Show that if there are multiple equilibria, the one with the highest cutoff \(p^*\) Pareto-dominates
(is better for all types than) the other equilibria.

\begin{enumerate}
\def\labelenumi{(\roman{enumi})}
\setcounter{enumi}{2}
\item
  Is there a negative stock price reaction upon announcement of an equity issue?
\item
  Focusing on an interior Pareto-dominant equilibrium, show that, when \(\tau\) increases, the
  volume of equity issues increases.
\end{enumerate}

\paragraph{Context needed to solve it}

\emph{From Tirole §6.2 (privately-known-prospects model), p.~241 (PDF 254):}

\emph{Privately-known-prospects model.} A borrower/entrepreneur has no funds (\(A = 0\)) to finance
a project costing \(I\). The project yields \(R\) in the case of success and 0 in the case of
failure. The borrower and the lenders are risk neutral, and the borrower is protected by limited
liability. The interest rate in the economy is normalized at 0.

The borrower can be one of two types. A good borrower has a probability of success equal to \(p\). A
bad borrower has a probability of success \(q\). Assume that \(p > q\) and that \(pR > I\) (at least
the good type is creditworthy). {[}\ldots{]}

The borrower has private information about her type. The capital market, which is competitive and
demands an expected rate of return equal to 0, puts probabilities \(\alpha\) and \(1 - \alpha\) on
the borrower being a good or a bad type, respectively. Under asymmetric information, the capital
market does not know whether it faces a ``\(p\)-borrower'' (a good borrower) or a ``\(q\)-borrower''
(a bad borrower). Let

\[m \equiv \alpha p + (1 - \alpha) q\]

denote the investors' prior probability of success.

\emph{From Tirole §6.2.2, Application 2: Assets in Place, the Negative Stock Price Reaction, and the
Going-Public Decision, pp.~244--245 (PDF 257--258):}

Let us next suppose that the entrepreneur already owns a project that, without further investment,
will succeed with probability \(p\) or \(q\), yielding profit \(R\). As before, the entrepreneur
knows the probability of success while the investors put probability \(\alpha\) on \(p\) and
\((1 - \alpha)\) on \(q\). Thus, in the absence of further information (and so the investors'
expectation of the probability of success is \(m\)), the assets in place are undervalued
(respectively, overvalued) if the true probability of success is \(p\) (respectively, \(q\)).

For computational simplicity, we will assume that the entrepreneur initially owns all shares. But
nothing is altered if she owns only a fraction of the shares. By ``stock price reaction upon the
announcement of an equity issue,'' we mean the difference between the total value of shares (whoever
owns them) before and after the announcement. This notion corresponds to the approach taken by event
studies in empirical work.

An equity offering may be motivated by a profitable ``deepening investment'' (more generally, the
key feature is that one cannot contract on the cash flow generated by this investment separately
from that generated by assets in place: the incomes generated by the two are intertwined or
fungible\textsuperscript{22}). At cost \(I\), the probability of success can be raised by an amount
\(\tau\) such that

\[\tau R > I.\]

That is, investing is efficient for both types of borrowers. Note that we assume for the moment that
the increase in profitability is uniform across types: the probability of success becomes
\(p + \tau\) for a good borrower and \(q + \tau\) for a bad one.

The entrepreneur, however, has no cash on hand. Accordingly, the full amount \(I\) must be raised
from investors. The entrepreneur must therefore issue new shares, thereby reducing the fraction of
shares she owns.

A key insight is that relinquishing shares to investors is relatively less costly to the borrower
with overvalued assets in place (the bad borrower) than to the borrower with undervalued assets in
place (the good borrower). Thus, if the good borrower conducts an equity offering, so does the bad
one.

Let us therefore investigate the possibility of an (efficient) pooling equilibrium. The entrepreneur
must offer a stake \(R_l\) in success to the investors such that

\[[\alpha(p + \tau) + (1 - \alpha)(q + \tau)]R_l = I \iff (m + \tau)R_l = I,\]

where, as earlier, \(m \equiv \alpha p + (1 - \alpha)q\) is the prior mean probability of success.
There exists a unique \(R_l\), \(0 < R_l < R\), satisfying this condition.

The good borrower, though, can guarantee herself \(pR\) by not diluting her
stake.\textsuperscript{23} Thus, she is willing to issue new shares only if

\[(p + \tau)(R - R_l) \geqslant pR \iff \tau R \geqslant \frac{p + \tau}{m + \tau} I. \qquad (6.1)\]

{[}\ldots{]} We are thus led to consider two situations:

\emph{Pooling equilibrium.} If condition (6.1) holds, then both types conduct an equity
offering.\textsuperscript{24} If the accrual of this deepening investment is
anticipated,\textsuperscript{25} the total value of shares \emph{before and after} the seasoned
equity offering is

\[(m + \tau)R - I.\]

There is no stock price reaction to the offering, which is perfectly anticipated and uninformative.

\emph{Separating equilibrium.} More interestingly, suppose that condition (6.1) is violated. The
good borrower then does not raise funds. The bad borrower still does, but under market conditions
that are not as favorable as in a pooling equilibrium. Because the investors know that the equity
offering reveals overvalued assets, they demand a higher stake \(R_l^B > R_l\) such that

\[(q + \tau)R_l^B = I.\]

The good borrower does not want to raise funds because

\[(p + \tau)(R - R_l^B) < pR \iff \tau R < \frac{p + \tau}{q + \tau} I, \qquad (6.2)\]

which holds if condition (6.1) is violated.

The announcement of a seasoned equity offering then leads to a \emph{negative stock price reaction}.
The preannouncement total value of shares is\textsuperscript{26}

\[V_0 = \alpha[pR] + (1 - \alpha)[(q + \tau)R - I].\]

After the announcement, it becomes

\[V_1 = (q + \tau)R - I.\]

Hence,

\[V_0 > V_1 \iff pR > (q + \tau)R - I.\]

But we know that

\[pR > (p + \tau)\left(R - \frac{I}{q + \tau}\right),\]

and so \emph{a fortiori} \(V_0 > V_1\).

\begin{quote}
\begin{enumerate}
\def\labelenumi{\arabic{enumi}.}
\setcounter{enumi}{21}
\item
  Otherwise, it would be optimal for the good borrower to engage in project finance so as to avoid
  having to cross-subsidize the bad one.
\item
  That this ``reservation utility'' depends on the borrower's type is the essential difference with
  the barebones model. Here, in the jargon of incentive theory, ``reservation utilities are
  type-contingent.'' {[}\ldots{]}
\item
  This pooling equilibrium is not unique whenever
  \((p + \tau)I/(m + \tau) \leqslant \tau R \leqslant (p + \tau)I/(q + \tau)\); indeed, if investors
  believe that an equity offering comes from a bad borrower and
  \(\tau R \leqslant (p + \tau)I/(q + \tau)\), then the good type indeed prefers not to raise funds.
  However, the pooling equilibrium is the \emph{Pareto-dominant} equilibrium (it is the best
  equilibrium for both the good and the bad borrower), and so we will focus on it. Note also that if
  condition (6.1) holds, then the bad borrower definitely prefers to raise funds since the analogous
  condition for her is \(\tau R \geqslant \frac{q + \tau}{m + \tau} I\), which is always satisfied.
\item
  Otherwise, the news of the existence of an investment opportunity by itself raises the value of
  shares.
\item
  As in the previous footnote, note that we assume that the investment opportunity is perfectly
  anticipated by the capital market. Otherwise, the issue of new securities could convey good news
  about the firm's opportunity set and the concomitant boost in share price might dominate the
  effect unveiled here.
\end{enumerate}
\end{quote}
\stmtend
%% END-STATEMENT

\asked (i) Show that in any equilibrium the firms that issue equity to fund the deepening investment are exactly the types below some cutoff $p^*$. (ii) Show that $p^*>\underline p$; that an interior cutoff satisfies $\tau R/I=(p^*+\tau)/(m^-(p^*)+\tau)$; that $p^*<\bar p$ when the gains from investing are not too large; and that the equilibrium with the highest cutoff is best for every type. (iii) Say whether an issue announcement lowers the stock price. (iv) At an interior Pareto-dominant equilibrium, show that the volume of equity issues rises with $\tau$.

\approach This is Tirole's Application 2 (assets in place and the negative stock-price reaction, \S6.2.2; Myers--Majluf in lecture block 3.5) with a continuum of types. An issuer pools with all other issuers at a single stake $R_l$, so issuing dilutes her assets in place at the pooled price; the dilution costs more the better her assets ($p$), which yields a cutoff rule. Investors' break-even over the issuing interval $[\underline p,p^*]$ gives $R_l$; indifference of the cutoff type gives the equilibrium condition. The price reaction compares total share value before and after the market learns $p\le p^*$; the comparative statics in $\tau$ use the implicit function theorem.

\notation
\begin{itemize}
  \item $p\in[\underline p,\bar p]$: privately known success probability of assets in place, c.d.f.\ $F$, density $f>0$; $E[p]$: prior mean.
  \item $R$: profit in success, 0 in failure; $I$: cost of the deepening investment, which raises $p$ to $p+\tau$, $\tau R>I$, $\bar p+\tau\le1$.
  \item $R_l$: investors' stake in success (they get 0 in failure); $p^*$: issue cutoff.
  \item $m^-(p)=E[\tilde p\mid\tilde p\le p]$, $m^+(p)=E[\tilde p\mid\tilde p\ge p]$; $0<(m^\pm)'\le1$.
  \item $V_0$, $V_1$: total value of shares before / after the issue announcement; $H(p^*,\tau)$: equilibrium function of (iv).
\end{itemize}

\step{1}{Issuers form a lower interval (part (i))} Not issuing gives the entrepreneur $pR$. Issuing for a
stake $R_l$ gives $(p+\tau)(R-R_l)$, so type $p$ issues iff
\begin{align}
  (p+\tau)(R-R_l) &\ge pR \nonumber\\
  pR+\tau R-(p+\tau)R_l &\ge pR \nonumber\\
  \tau R &\ge (p+\tau)R_l. \label{t3e15-issue}
\end{align}
In any equilibrium all issuers use the same stake: every type's payoff falls with $R_l$, so if two stakes
were accepted every issuer would pick the smaller. The left side of \eqref{t3e15-issue} does not depend on $p$ and the right side rises with $p$ (slope $R_l>0$),
so if type $p$ issues, every type below $p$ issues too: the issuers are the types $p\le p^*$ for some cutoff $p^*\in[\underline p,\bar p]$. $\blacksquare$

\answer{(i)} Type $p$ issues iff $\tau R\ge(p+\tau)R_l$; the cost of dilution $(p+\tau)R_l$ rises with $p$, so issuers are the types $p\le p^*$.

\step{2}{Break-even stake and the interior cutoff (part (ii))} Investors break even on the issuers, whose average success probability after investing is $E(\tilde p+\tau\mid\tilde p\le p^*)=m^-(p^*)+\tau$:
\begin{align}
  \bigl[E(\tilde p\mid\tilde p\le p^*)+\tau\bigr]R_l &= I \nonumber\\
  R_l &= \frac{I}{m^-(p^*)+\tau}. \label{t3e15-Rl}
\end{align}
If $p^*<\bar p$, the cutoff type is indifferent in \eqref{t3e15-issue} (types just below issue, types just above do not, and payoffs are continuous in $p$). Substituting \eqref{t3e15-Rl} and dividing by $I$,
\begin{align}
  \tau R &= (p^*+\tau)\frac{I}{m^-(p^*)+\tau} \nonumber\\
  \frac{\tau R}{I} &= \frac{p^*+\tau}{m^-(p^*)+\tau}. \label{t3e15-cut}
\end{align}

\step{3}{$p^*>\underline p$} Whatever the investors believe, they put mass only on $[\underline p,\bar p]$,
so the stake they require is at most $I/(\underline p+\tau)$. Under that stake, \eqref{t3e15-issue} holds
for type $\underline p$ with strict inequality:
\begin{align}
  (\underline p+\tau)\frac{I}{\underline p+\tau} = I < \tau R. \nonumber
\end{align}
By continuity it also holds strictly for a positive-measure set of types above $\underline p$. So the
cutoff cannot be $\underline p$ (equivalently, $m^-(\underline p)=\underline p$ makes the right side of
\eqref{t3e15-cut} equal 1, below $\tau R/I$).

\step{4}{$p^*<\bar p$ when benefits are not too large} If $p^*=\bar p$ (everyone issues), then
$m^-(\bar p)=E[p]$ and \eqref{t3e15-Rl} gives $R_l=I/(E[p]+\tau)$. Type $\bar p$ would then issue only if
\begin{align}
  \tau R \ge (\bar p+\tau)\frac{I}{E[p]+\tau}
  \iff \frac{\tau R}{I} \ge \frac{\bar p+\tau}{E[p]+\tau}, \nonumber
\end{align}
which the assumed condition $\tau R/I<(\bar p+\tau)/(E[p]+\tau)$ rules out. Types near $\bar p$ therefore
strictly prefer not to issue, a contradiction, so $p^*<\bar p$. An interior solution exists:
$\varphi(p)\equiv(p+\tau)/(m^-(p)+\tau)$ is continuous, $\varphi(\underline p)=1<\tau R/I$ and
$\varphi(\bar p)>\tau R/I$, so by the intermediate value theorem \eqref{t3e15-cut} has a root in
$(\underline p,\bar p)$.

\step{5}{The highest cutoff Pareto-dominates} Take equilibria with cutoffs $p_1^*<p_2^*$. Since
$(m^-)'>0$, \eqref{t3e15-Rl} gives $R_{l2}<R_{l1}$. Types $p\le p_1^*$ issue in both and pay less in the
second. Types $p\in(p_1^*,p_2^*]$ get $pR$ in the first and choose to issue in the second, so by
\eqref{t3e15-issue} they get $(p+\tau)(R-R_{l2})\ge pR$. Types above $p_2^*$ get $pR$ in both. Investors
break even in both. Hence the highest-cutoff equilibrium is weakly better for every type. $\blacksquare$
(The condition $(m^-)'\le1$ does not by itself make the root unique; for a uniform $F$, $(m^-)'=\frac12$
and the root is unique.)

\answer{(ii)} $p^*>\underline p$ (the lowest type always gains from issuing); an interior cutoff solves $\tau R/I=(p^*+\tau)/(m^-(p^*)+\tau)$ with $R_l=I/(m^-(p^*)+\tau)$; $\tau R/I<(\bar p+\tau)/(E[p]+\tau)$ rules out $p^*=\bar p$; and a higher cutoff means a lower $R_l$, so the highest-cutoff equilibrium Pareto-dominates.

\step{6}{Negative stock-price reaction (part (iii))} The opportunity is anticipated. Before the
announcement, total share value is expected value net of the funds investors will contribute (they break
even): an issuer ($p\le p^*$, probability $F(p^*)$) is worth $(p+\tau)R-I$, a non-issuer $pR$. Averaging, and using $E[p]=F(p^*)m^-(p^*)+[1-F(p^*)]m^+(p^*)$,
\begin{align}
  V_0 &= F(p^*)\bigl[(m^-(p^*)+\tau)R-I\bigr]+[1-F(p^*)]m^+(p^*)R \label{t3e15-V0}\\
      &= \bigl\{F(p^*)m^-(p^*)+[1-F(p^*)]m^+(p^*)\bigr\}R+F(p^*)(\tau R-I) \nonumber\\
      &= E[p]R+F(p^*)(\tau R-I). \nonumber
\end{align}
After the announcement the market knows $p\le p^*$:
\begin{align}
  V_1 = [m^-(p^*)+\tau]R-I. \label{t3e15-V1}
\end{align}
Subtracting \eqref{t3e15-V1} from \eqref{t3e15-V0},
\begin{align}
  V_0-V_1 = [1-F(p^*)]\Bigl\{m^+(p^*)R-\bigl[(m^-(p^*)+\tau)R-I\bigr]\Bigr\}. \label{t3e15-dV}
\end{align}
At an interior equilibrium, \eqref{t3e15-cut} gives $I=\tau R(m^-+\tau)/(p^*+\tau)$, so
\begin{align}
  (m^-+\tau)R-I &= (m^-+\tau)R\Bigl[1-\frac{\tau}{p^*+\tau}\Bigr] = (m^-+\tau)R\,\frac{p^*}{p^*+\tau}, \nonumber\\
  V_0-V_1 &= [1-F(p^*)]R\Bigl[m^+(p^*)-p^*\,\frac{m^-(p^*)+\tau}{p^*+\tau}\Bigr]. \label{t3e15-dV2}
\end{align}
Since $m^+(p^*)>p^*$ and $(m^-(p^*)+\tau)/(p^*+\tau)<1$, the bracket is positive and $1-F(p^*)>0$: the
announcement of an equity issue lowers the stock price. If everyone issues ($p^*=\bar p$, pooling), the
factor $1-F(p^*)=0$ and there is no reaction.

\answer{(iii)} Yes, at an interior equilibrium: $V_0-V_1=[1-F(p^*)]R\,[m^+(p^*)-p^*(m^-(p^*)+\tau)/(p^*+\tau)]>0$, because the announcement tells the market that $p\le p^*$. With full pooling there is no reaction.

\step{7}{Comparative statics in $\tau$ (part (iv))} Define
\begin{align}
  H(p^*,\tau) \equiv \frac{\tau R}{I}\bigl[m^-(p^*)+\tau\bigr]-(p^*+\tau), \label{t3e15-H}
\end{align}
so that \eqref{t3e15-cut}, multiplied by $m^-(p^*)+\tau$ and rearranged, is $H=0$. Then $H(\underline p,\tau)=(\tau R/I)(\underline p+\tau)-(\underline p+\tau)=(\tau R/I-1)(\underline p+\tau)>0$ (using $m^-(\underline p)=\underline p$) and
$H(\bar p,\tau)<0$ under the condition of (ii). At the highest root (the Pareto-dominant one) $H$ crosses
from positive to negative, so $H_{p^*}=(\tau R/I)(m^-)'-1<0$ (taking the crossing to be transversal).
From $H=0$, $(\tau R/I)(m^-+\tau)=p^*+\tau$, so
\begin{align}
  \frac{\tau R}{I}-1 = \frac{p^*+\tau}{m^-+\tau}-1 = \frac{p^*-m^-}{m^-+\tau}. \nonumber
\end{align}
Differentiating \eqref{t3e15-H} in $\tau$ (product rule on $(\tau R/I)(m^-+\tau)$) and substituting this,
\begin{align}
  H_\tau &= \frac RI(m^-+\tau)+\frac{\tau R}{I}-1 \nonumber\\
         &= \frac RI(m^-+\tau)+\frac{p^*-m^-(p^*)}{m^-(p^*)+\tau} > 0. \label{t3e15-Ht}
\end{align}
By the implicit function theorem,
\begin{align}
  \frac{dp^*}{d\tau} = -\frac{H_\tau}{H_{p^*}} > 0, \label{t3e15-dp}
\end{align}
and the volume of equity issues, $F(p^*)I$ (the issuers are the types below $p^*$), rises with $\tau$, since
$d[F(p^*)I]/d\tau=If(p^*)\,dp^*/d\tau>0$. $\blacksquare$ (The sign of $H_\tau$: $m^-+\tau>0$ and $p^*>m^-(p^*)$, since a truncated mean lies below the truncation point.)

\answer{(iv)} $dp^*/d\tau=-H_\tau/H_{p^*}>0$, so the issue volume $F(p^*)I$ rises with $\tau$ at rate $If(p^*)\,dp^*/d\tau>0$.

\step{8}{Numerical illustration (uniform $F$)} Take $p\sim U[0.1,0.5]$, $\tau=0.3$, $R=10$, $I=2.5$:
$\tau R=3>I$, $E[p]=0.3$, and $\tau R/I=1.2<(0.8)/(0.6)=1.333333$. With $m^-(p)=(p+0.1)/2$,
\eqref{t3e15-cut} reads $1.2[(p^*+0.1)/2+0.3]=p^*+0.3$, so $0.6p^*+0.06+0.36=p^*+0.3$, i.e.\ $0.42-0.3=0.4p^*$ and $p^*=0.12/0.4=0.3$. Then
$R_l=2.5/(0.2+0.3)=5$, and the cutoff type is indifferent: $(0.6)(5)=3=(0.3)(10)$. Also $F(p^*)=0.5$,
$m^+(p^*)=0.4$, $V_0=3+0.5(0.5)=3.25$, $V_1=(0.5)(10)-2.5=2.5$, and $V_0-V_1=0.75$, which matches
\eqref{t3e15-dV2}: $(0.5)(10)[0.4-0.3(0.5)/0.6]=0.75$. For (iv), $H_{p^*}=1.2(0.5)-1=-0.4$ and
$H_\tau=4(0.5)+0.1/0.5=2.2$, so $dp^*/d\tau=5.5$ and the issue volume $F(p^*)I=1.25$ rises at rate
$2.5\times2.5\times5.5=34.375$.

\intuition Issuing equity dilutes the entrepreneur's assets in place at the pooled price. Types whose
assets are worth more (high $p$) lose more from that, so only types below a cutoff issue, and an issue
announcement tells the market the firm is in the lower part of the distribution: the price drops.
Equilibria with higher cutoffs carry less stigma (a lower $R_l$) and are better for everyone. A more
valuable investment ($\tau$ up) makes higher types willing to accept the dilution, so the cutoff and
the volume of issues rise.

\pitfalls
\begin{itemize}
  \item Pricing the issue at the prior mean $E[p]$: investors know that only types $p\le p^*$ issue, so the stake uses $m^-(p^*)$.
  \item Measuring the price reaction on the entrepreneur's stake only: as in Tirole, the reaction is the change in the \emph{total} value of shares, net of the funds investors contribute.
  \item Writing the issue volume as $[1-F(p^*)]I$ (the book's typo): the issuers are the types \emph{below} the cutoff, so the volume is $F(p^*)I$.
\end{itemize}

\bookans Matches the book's answer (pp.~597--598): cutoff rule, $R_l=I/(m^-(p^*)+\tau)$, the interior
condition, $p^*>\underline p$, the Pareto ranking, $V_0-V_1>0$ and $dp^*/d\tau>0$. Two typos in the
book: its stock-price item is labelled ``(ii)'' instead of (iii), and it gives the volume of issues as
``$[1-F(p^*)]I$'', whereas the issuers are the types $p\le p^*$, so the volume is $F(p^*)I$. The conclusion
(volume rises with $\tau$) is unaffected.

\refbib{Tirole Ch.~6, \S6.2.2 Application 2 ``Assets in Place, the Negative Stock Price Reaction, and the
Going-Public Decision'', pp.~244--245; Exercise 6.5, p.~275; answer pp.~597--598. Lecture Topic~3,
pecking-order block.}

% =============================================================================
% Topic 3 list, original exercises 16--19 (Stiglitz--Weiss; collateral as a signal).
% No book exercise exists; written with the user's agreement (2026-11-12).
% =============================================================================
\exsec{16. Stiglitz--Weiss with two types: the lender's return is not monotone in the rate --- Original exercise}
%% BEGIN-STATEMENT 16
\stmtbegin
\textbf{Origin:} original exercise on Stiglitz--Weiss (1981) credit rationing with mean-preserving
spreads, built on Prof.~Ricca's Topic 3 slides 47--55 (Aula/Topic\_3.pdf) and Tirole §3.1
Introduction, p.~113 (PDF 126), and §6.1 Introduction, footnote 5, p.~240 (PDF 253) · Answer: in the
solutions file

A lender finances projects that cost \(I = 100\) each. Borrowers have no money of their own and
there are two types, which the lender cannot tell apart. A low-risk (safe) borrower's project yields
\(R_L = 150\) with probability \(p_L = 0.8\) and 0 otherwise. A high-risk (risky) borrower's project
yields \(R_H = 300\) with probability \(p_H = 0.4\) and 0 otherwise. A share \(\alpha = 0.75\) of
the borrowers is low-risk. Borrowers and the lender are risk neutral. The lender offers only a
standard debt contract: the borrower receives \(I\) and repays \((1 + r)I\) if the project succeeds
and nothing if it fails (limited liability; a failed project is worth 0). A borrower applies for a
loan if and only if her expected payoff is non-negative.

\begin{enumerate}
\def\labelenumi{(\alph{enumi})}
\item
  Show that both projects have the same expected return \(\bar R\) and compute it. Compute the
  variance of each project's return. Show that the high-risk return is a mean-preserving spread of
  the low-risk return, i.e.~that it can be obtained by adding to each realisation of the low-risk
  return a random variable with zero mean conditional on that realisation.
\item
  Write each type's expected payoff from borrowing at rate \(r\) as a function of \(r\), and
  evaluate both at \(r = 20\%\). Which type is hurt more by a rise in \(r\), and why?
\item
  Compute the rate \(r^*\) above which the low-risk borrowers stop applying and the rate \(r_{max}\)
  above which nobody applies. Who applies for \(r \le r^*\), for \(r^* < r \le r_{max}\) and for
  \(r > r_{max}\)?
\item
  Let \(\rho(r)\) be the lender's expected gross return per unit lent (expected repayment divided by
  \(I\)). Write \(\rho(r)\) for \(r \le r^*\) and for \(r^* < r \le r_{max}\). Evaluate it at
  \(r = 20\%\), \(r = r^*\), \(r = 60\%\), \(r = 100\%\) and \(r = r_{max}\). Compute the drop in
  \(\rho\) when \(r\) crosses \(r^*\) and show that it equals \(\alpha(p_L - p_H)(1 + r^*)\). Is
  \(\rho\) monotone in \(r\)?
\item
  Find the rate \(\tilde r > r^*\) at which the lender's return first gets back to \(\rho(r^*)\).
  Explain why a lender that faces more applicants than it can fund at \(r^*\) would never raise its
  rate to any value in \((r^*, \tilde r)\), and why this makes the supply of loans
  ``backward-bending''.
\item
  Suppose instead that the lender finances the project with equity: it takes a share \(s\) of the
  project's return, in both success and failure. Find the share \(s\) that gives the lender the same
  expected gross return per unit lent as the debt contract at \(r^*\). Which types apply, and what
  is each type's expected payoff? Explain why the adverse selection of parts (c)--(e) disappears.
\end{enumerate}

\paragraph{Context needed to solve it}

\emph{From Prof.~Ricca's Topic 3 slides 47--55 (block 3.6, Stiglitz and Weiss 1981):}

\begin{itemize}
\tightlist
\item
  Notion of risk: mean-preserving spreads. Projects have the same mean; a mean-preserving spread of
  a random variable \(X\) is obtained by adding a zero-mean random variable to each possible
  realisation of \(X\) (slide 47).
\item
  Asymmetric information and adverse selection may make the supply curve backward-bending at high
  interest rates: the pool of borrowers gets worse when rates rise, because the payoff in the
  default state does not depend on the rate, so a rise in the rate hurts the borrowers who default
  more often less (slide 48). Supply and demand may then not intersect, and demand can exceed supply
  at the prevailing rate: credit rationing.
\item
  Two-type example (slides 50--54). Low risk: \(R_L\) with probability \(p_L\), 0 otherwise; high
  risk: \(R_H\) with probability \(p_H\), 0 otherwise; \(p_L > p_H\) and
  \(p_L R_L = p_H R_H = \bar R\) (the high-risk project is a mean-preserving spread of the low-risk
  one). Borrowers and lenders are risk neutral; types are private information; the project requires
  \(I\). At rate \(r\) the expected payoffs are \(p_L[R_L - (1 + r)I] = \bar R - p_L(1 + r)I\) and
  \(p_H[R_H - (1 + r)I] = \bar R - p_H(1 + r)I\).
\item
  \(r_{max}\) solves \(\bar R - p_H(1 + r_{max})I = 0\), i.e.~\(r_{max} = \bar R/(p_H I) - 1\);
  \(r^*\) solves \(\bar R - p_L(1 + r^*)I = 0\), i.e.~\(r^* = \bar R/(p_L I) - 1\). For \(r < r^*\)
  both types apply; for \(r^* < r < r_{max}\) only high-risk borrowers apply.
\item
  With \(\alpha\) the proportion of low-risk types, the lender's expected payoff is \(m(1 + r)\)
  with \(m = \alpha p_L + (1 - \alpha)p_H\) when \(r < r^*\), and \(p_H(1 + r)\) when
  \(r^* < r < r_{max}\). The expected return drops sharply around \(r^*\) as the good types leave
  the market; it is not a monotonic function of the rate (slides 53--54).
\end{itemize}

\emph{From Tirole §6.1 Introduction, footnote 5, p.~240 (PDF 253):} ``Stiglitz and Weiss assume that
lenders can offer only debt contracts, and show that the repayment probability decreases with the
rate of interest offered by lenders, and that the loan market is characterized by credit rationing.
However, the assumptions of the model predict that investors should instead offer equity contracts,
in which case there would be no adverse selection (all firms have the same mean income) and no
credit rationing (Hart 1985).''
\stmtend
%% END-STATEMENT

\asked (a) Check that the two projects have the same mean, compute their variances, and show that the risky return is the safe return plus zero-mean noise. (b) Write each type's expected payoff from a loan at rate $r$, evaluate it at $20\%$, and say which type a rate rise hurts more. (c) Find the rate $r^*$ at which the safe types stop applying and the rate $r_{max}$ at which everybody stops, and who applies on each interval. (d) Write the lender's expected gross return per unit lent, $\rho(r)$, evaluate it at five rates, measure the drop at $r^*$, and say whether $\rho$ is monotone. (e) Find the rate $\tilde r$ at which $\rho$ recovers its value at $r^*$ and explain why rates in $(r^*,\tilde r)$ are never charged. (f) Replace debt with an equity share $s$ and show that the selection problem disappears.

\approach This is the two-type Stiglitz--Weiss example of lecture block 3.6. Under a debt contract with limited liability the borrower repays only on success, so the rate is paid with probability $p_i$: the borrower who fails more often is hurt less by a higher rate. The quoted rate therefore selects the pool of applicants, and the lender's return combines a price effect (higher $r$) with a composition effect (safe types leave at $r^*$). Equity pays a share of the return in every state, so its value depends only on the mean, which mean-preserving spreads hold fixed (Hart's point quoted by Tirole).

\notation
\begin{itemize}
  \item $I=100$: loan size; $R_L=150$, $p_L=0.8$: success payoff and probability of the low-risk (safe) project; $R_H=300$, $p_H=0.4$: same for the high-risk project; payoff 0 on failure.
  \item $\alpha=0.75$: share of low-risk borrowers; $m=\alpha p_L+(1-\alpha)p_H$: success probability of the pooled applicant.
  \item $\bar R=p_LR_L=p_HR_H$: common expected return; $X_L$, $X_H$: random project returns.
  \item $r$: quoted rate (repay $(1+r)I$ on success); $U_L(r)$, $U_H(r)$: borrowers' expected payoffs.
  \item $r^*$, $r_{max}$: exit rates of the safe and the risky type; $\rho(r)$: lender's expected gross return per unit lent; $\tilde r$: rate above $r^*$ with $\rho(\tilde r)=\rho(r^*)$.
  \item $s$: lender's equity share of the project return (part (f)).
\end{itemize}

\step{1}{Same mean, different risk (part (a))} We compute the means and the variances, then build the spread explicitly. The means are
\begin{align}
  \bar R = p_LR_L &= (0.8)(150) = 120, \nonumber\\
  p_HR_H &= (0.4)(300) = 120. \label{t3e16-mean}
\end{align}
For a return equal to $R_i$ with probability $p_i$ and 0 otherwise, $E[X_i^2]=p_iR_i^2$, so
\begin{align}
  \operatorname{Var}(X_i) = p_iR_i^2-(p_iR_i)^2 = p_i(1-p_i)R_i^2. \label{t3e16-var}
\end{align}
Substituting each type's numbers,
\begin{align}
  \operatorname{Var}(X_L) &= (0.8)(0.2)(150^2) = (0.16)(22{,}500) = 3{,}600, \nonumber\\
  \operatorname{Var}(X_H) &= (0.4)(0.6)(300^2) = (0.24)(90{,}000) = 21{,}600. \nonumber
\end{align}
To show the spread, define $Z=0$ when $X_L=0$, and $Z=+150$ or $Z=-150$ with probability $\tfrac12$ each when $X_L=150$. Then $E[Z\mid X_L]=0$ in both states, and $X_L+Z$ takes the values $150+150=300$ and $150-150=0$ after a success and $0$ after a failure, with
\begin{align}
  P(X_L+Z=300) &= P(X_L=150)\cdot\tfrac12 = (0.8)(0.5) = 0.4 = p_H, \nonumber\\
  P(X_L+Z=0) &= 0.2+(0.8)(0.5) = 0.6 = 1-p_H. \nonumber
\end{align}
So $X_L+Z$ has exactly the distribution of $X_H$. Because $E[Z\mid X_L]=0$, $Z$ is uncorrelated with $X_L$ and the variance adds up: $3{,}600+E[Z^2]=3{,}600+(0.8)(150^2)=3{,}600+18{,}000=21{,}600$, as computed above.

\answer{(a)} $\bar R=120$ for both; $\operatorname{Var}(X_L)=3{,}600$, $\operatorname{Var}(X_H)=21{,}600$; $X_H$ has the distribution of $X_L+Z$ with $Z=\pm150$ (probability $\tfrac12$ each) after a success and $Z=0$ after a failure, $E[Z\mid X_L]=0$: a mean-preserving spread.

\step{2}{Borrowers' payoffs (part (b))} A borrower keeps $R_i-(1+r)I$ on success and 0 on failure (limited liability). Multiplying out and using $p_iR_i=\bar R$,
\begin{align}
  U_i(r) = p_i\bigl[R_i-(1+r)I\bigr] = \bar R-p_i(1+r)I. \label{t3e16-U}
\end{align}
At $r=0.2$, so that $(1+r)I=120$,
\begin{align}
  U_L(0.2) &= 120-(0.8)(120) = 120-96 = 24, \nonumber\\
  U_H(0.2) &= 120-(0.4)(120) = 120-48 = 72. \nonumber
\end{align}
Differentiating \eqref{t3e16-U}, $U_i'(r)=-p_iI$: $-80$ for the safe type and $-40$ for the risky type, so one more percentage point of interest costs the safe borrower $0.8$ and the risky borrower $0.4$.

\answer{(b)} $U_L(r)=120-80(1+r)$, $U_H(r)=120-40(1+r)$; at $20\%$, $U_L=24$ and $U_H=72$. A rate rise hurts the safe type twice as much ($-80$ vs.\ $-40$ per unit of $r$): she repays with probability $0.8$, the risky type only with $0.4$.

\step{3}{Exit rates (part (c))} A type applies while $U_i(r)\ge0$. Setting \eqref{t3e16-U} to zero and solving for $r$,
\begin{align}
  \bar R-p_i(1+r)I = 0 \iff 1+r = \frac{\bar R}{p_iI}. \nonumber
\end{align}
For the safe type and for the risky type,
\begin{align}
  r^* &= \frac{\bar R}{p_LI}-1 = \frac{120}{(0.8)(100)}-1 = 1.5-1 = 0.5, \label{t3e16-rstar}\\
  r_{max} &= \frac{\bar R}{p_HI}-1 = \frac{120}{(0.4)(100)}-1 = 3-1 = 2.0. \label{t3e16-rmax}
\end{align}

\answer{(c)} $r^*=50\%$, $r_{max}=200\%$. Both types apply for $r\le50\%$, only risky borrowers for $50\%<r\le200\%$, nobody above $200\%$.

\step{4}{The lender's return (part (d))} The lender receives $(1+r)I$ when the project succeeds and 0 otherwise, so per unit lent it expects $(1+r)$ times the success probability of the average applicant. For $r\le r^*$ both types apply in proportions $\alpha$ and $1-\alpha$:
\begin{align}
  m &= \alpha p_L+(1-\alpha)p_H = (0.75)(0.8)+(0.25)(0.4) = 0.6+0.1 = 0.7, \nonumber\\
  \rho(r) &= m(1+r) = 0.7(1+r), \qquad r\le r^*. \label{t3e16-rho1}
\end{align}
For $r^*<r\le r_{max}$ only risky borrowers apply:
\begin{align}
  \rho(r) = p_H(1+r) = 0.4(1+r), \qquad r^*<r\le r_{max}. \label{t3e16-rho2}
\end{align}
Evaluating \eqref{t3e16-rho1} and \eqref{t3e16-rho2},
\begin{align}
  \rho(0.2) &= (0.7)(1.2) = 0.84, & \rho(r^*) &= (0.7)(1.5) = 1.05, \nonumber\\
  \rho(0.6) &= (0.4)(1.6) = 0.64, & \rho(1.0) &= (0.4)(2.0) = 0.80, \nonumber\\
  \rho(r_{max}) &= (0.4)(3.0) = 1.20. \nonumber
\end{align}
Just above $r^*$ the return is $p_H(1+r^*)=(0.4)(1.5)=0.60$, so the drop at $r^*$ is
\begin{align}
  \rho(r^*)-p_H(1+r^*) = (m-p_H)(1+r^*) = (0.7-0.4)(1.5) = 0.45. \label{t3e16-drop}
\end{align}
Since $m-p_H=\alpha p_L+(1-\alpha)p_H-p_H=\alpha(p_L-p_H)$, the drop is also $\alpha(p_L-p_H)(1+r^*)=(0.75)(0.4)(1.5)=0.45$. Note also $\rho(r_{max})=p_H\cdot\bar R/(p_HI)=\bar R/I=1.2$: at $r_{max}$ the lender takes the risky project's whole expected return.

\answer{(d)} $\rho(r)=0.7(1+r)$ for $r\le50\%$ and $0.4(1+r)$ for $50\%<r\le200\%$; $\rho=0.84,\ 1.05,\ 0.64,\ 0.80,\ 1.20$ at $20\%,\ 50\%,\ 60\%,\ 100\%,\ 200\%$; the drop at $r^*$ is $0.45=\alpha(p_L-p_H)(1+r^*)$. $\rho$ is not monotone: it rises, falls by $0.45$ at $r^*$, then rises again.

\step{5}{The dead zone above $r^*$ (part (e))} Set the risky-only return \eqref{t3e16-rho2} equal to $\rho(r^*)$ and solve:
\begin{align}
  p_H(1+\tilde r) &= m(1+r^*) \nonumber\\
  1+\tilde r &= \frac{1.05}{0.4} = 2.625 \nonumber\\
  \tilde r &= 1.625. \label{t3e16-rtilde}
\end{align}
For every $r\in(r^*,\tilde r)$, $\rho(r)=0.4(1+r)<0.4(2.625)=1.05=\rho(r^*)$. A lender with more applicants than funds at $r^*$ that raises its rate a little loses all its safe applicants and earns less per unit, not more. Since what lenders can pay their depositors is $\rho$, the funds offered to borrowers, as a function of the quoted rate, fall when $r$ passes $r^*$: the loan supply curve bends backward, and excess demand at $r^*$ is not removed by a higher rate.

\answer{(e)} $\tilde r=162.5\%$. On $(50\%,162.5\%)$ the return is below $1.05$: raising the rate past $r^*$ lowers the lender's return, so the loan supply bends backward at $r^*$.

\step{6}{Equity instead of debt (part (f))} With equity the lender receives $sX_i$ in every state, so its expected receipt is $sE[X_i]=s\bar R$ for both types. Matching the debt contract's return at $r^*$, $\rho(r^*)I=(1.05)(100)=105$:
\begin{align}
  s\bar R &= 105 \nonumber\\
  s &= \frac{105}{120} = 0.875. \label{t3e16-s}
\end{align}
Each type's expected payoff is $(1-s)\bar R=(0.125)(120)=15>0$, so both apply, and the lender's return $s\bar R/I$ does not depend on who applies.

\answer{(f)} $s=0.875$; both types apply and each expects $15$. Equity pays in proportion to the return, so its value depends only on the mean $\bar R$, which the two types share; the rate no longer selects the pool.

\intuition Under debt, the borrower's payoff is convex in the project return (she gets nothing below the repayment, everything above it), so riskier borrowers value a loan more and stay longer when the rate rises. The quoted rate then does two things: it raises what each success pays, and it changes who applies. At $r^*$ the second effect dominates and the return falls by $\alpha(p_L-p_H)(1+r^*)$, which is the source of the backward-bending supply. In this two-type example the return at $r_{max}$ ($1.2$) is above the return at $r^*$ ($1.05$), because at $r_{max}$ the lender takes the risky project's whole surplus; what that means for the market is the subject of exercise 17. Equity removes the selection altogether, which is why the Stiglitz--Weiss result needs the restriction to debt contracts.

\pitfalls
\begin{itemize}
  \item Using $m(1+r)$ above $r^*$: once the safe types leave, the applicant pool is all risky and the return is $p_H(1+r)$.
  \item Thinking that a riskier project means a lower mean: under a mean-preserving spread the mean is the same ($120$); only the dispersion rises.
  \item Calling $r^*$ the lender's global optimum here: in the two-type example $\rho(r_{max})=\bar R/I=1.2>1.05$; what is true is that $\rho$ falls when $r$ crosses $r^*$ and stays below $\rho(r^*)$ until $\tilde r$.
\end{itemize}

\refbib{Lecture Topic~3, block 3.6, slides 47--55 (Stiglitz and Weiss 1981, two-type example and numerical example). Tirole \S3.1 Introduction, p.~113 (adverse-selection explanation of rationing); \S6.1 Introduction, footnote~5, p.~240 (Hart's equity critique). Original exercise: no book answer.}

\exsec{17. Why raising the rate does not clear the loan market --- Original exercise}
%% BEGIN-STATEMENT 17
\stmtbegin
\textbf{Origin:} original exercise on Stiglitz--Weiss (1981) credit rationing in a competitive loan
market, built on Prof.~Ricca's Topic 3 slides 47--55 (Aula/Topic\_3.pdf) and Tirole §3.1
Introduction, p.~113 (PDF 126) · Answer: in the solutions file

There are \(N = 1{,}000\) borrowers. Each needs \(I = 100\) for a project and has no money of her
own. A share \(\alpha = 0.75\) of them is low-risk: the project yields \(R_L = 130\) with
probability \(p_L = 0.9\) and 0 otherwise. The rest are high-risk: the project yields \(R_H = 234\)
with probability \(p_H = 0.5\) and 0 otherwise. Types are private information. Banks are competitive
and risk neutral, offer only standard debt contracts (repay \((1 + r)I\) on success, nothing on
failure), make zero profit and pass to depositors the expected gross return per unit lent, \(\rho\).
Depositors supply funds

\[L^s(\rho) = \max\{0,\ 1{,}000{,}000\,(\rho - 0.96)\}.\]

A borrower applies if her expected payoff is non-negative (an indifferent borrower applies). When
applicants exceed funds, the banks choose the borrowers who get a loan at random among the
applicants.

\begin{enumerate}
\def\labelenumi{(\alph{enumi})}
\item
  Compute \(\bar R\), the rates \(r^*\) (above which the low-risk borrowers stop applying) and
  \(r_{max}\) (above which nobody applies), the banks' expected gross return \(\rho(r)\) on each
  interval, \(\rho^* \equiv \rho(r^*)\), and the rate \(\tilde r > r^*\) at which \(\rho\) is back
  to \(\rho^*\).
\item
  Compute loan demand and the supply of funds at \(r = r^*\). How many applicants are turned down,
  and how many of them are low-risk and high-risk?
\item
  One bank, facing excess demand, raises its rate to \(r = 40\%\). Who applies at that rate, what is
  \(\rho(40\%)\), and how much would depositors supply if all banks charged \(40\%\)? Explain why
  raising the rate does not clear the market.
\item
  Show that no rate \(r \le r^*\) equates supply and demand. Find the rate
  \(r_W \in (r^*, r_{max}]\) at which supply equals demand, and the return \(\rho(r_W)\). Show that
  \(r_W\) is not an equilibrium: a bank that offers \(r^*\) instead can pay its depositors more and
  still find borrowers. Conclude what the market interest rate is and whether there is credit
  rationing.
\item
  Comparative statics; change one parameter at a time and keep the others at their initial values.
\item
  The share of low-risk borrowers rises to \(\alpha = 0.9\). Compute \(\rho^*\) and \(L^s(\rho^*)\).
  Is there still rationing? If not, find the market-clearing rate.
\end{enumerate}

\begin{enumerate}
\def\labelenumi{(\roman{enumi})}
\setcounter{enumi}{1}
\tightlist
\item
  The high-risk project becomes riskier with the same mean: \(p_H = 0.4\), \(R_H = 292.5\)
  (\(\alpha = 0.75\)). Compute \(r^*\), \(r_{max}\), \(\rho^*\), \(L^s(\rho^*)\), the share of
  applicants turned down, the drop in \(\rho\) at \(r^*\), and check whether the rate that clears
  the market among high-risk borrowers only is an equilibrium.
\end{enumerate}

\begin{enumerate}
\def\labelenumi{(\alph{enumi})}
\setcounter{enumi}{5}
\tightlist
\item
  Go back to the initial parameters and the outcome of (b) and (d). A bank offers loans only to
  applicants who were turned down at \(r^*\), at a rate \(r' \in (\tilde r, r_{max}]\). Who accepts,
  and what is the bank's expected gross return? What will competition among such banks do to \(r'\)?
  Which borrowers remain without credit, and why can no rate above \(r^*\) serve them?
\end{enumerate}

\paragraph{Context needed to solve it}

\emph{From Prof.~Ricca's Topic 3 slides 47--55 (block 3.6, Stiglitz and Weiss 1981):} same model as
exercise 16, restated here.

\begin{itemize}
\tightlist
\item
  Mean-preserving spreads: \(p_L > p_H\) and \(p_L R_L = p_H R_H = \bar R\); risk-neutral borrowers
  and lenders; types are private information; investment \(I\). At rate \(r\) the expected payoffs
  are \(\bar R - p_L(1 + r)I\) (low risk) and \(\bar R - p_H(1 + r)I\) (high risk), so a rise in
  \(r\) hurts the high-risk borrowers less (slide 51).
\item
  \(r_{max} = \bar R/(p_H I) - 1\) (neither type applies when \(r > r_{max}\));
  \(r^* = \bar R/(p_L I) - 1\) (only high-risk borrowers apply when \(r^* < r < r_{max}\); both
  types apply when \(r < r^*\)) (slide 52).
\item
  Lender's expected payoff per unit: \(m(1 + r)\), \(m = \alpha p_L + (1 - \alpha)p_H\), for
  \(r < r^*\); \(p_H(1 + r)\) for \(r^* < r < r_{max}\); it drops sharply around \(r^*\) and is not
  monotone in \(r\) (slides 53--54).
\item
  The pool of borrowers gets worse when rates rise; the supply curve can be backward-bending at high
  rates; supply and demand may not intersect, so demand can exceed supply at the prevailing rate:
  credit rationing (slides 47--48).
\end{itemize}

\emph{From Tirole §3.1 Introduction, p.~113 (PDF 126):} ``Why are lenders not willing to raise
interest rates if the demand for loans exceeds their supply at the prevailing rates? {[}\ldots{]}
The adverse selection explanation is that, in a situation where lenders cannot directly tell good
and bad borrowers apart, higher interest rates tend to attract low-quality borrowers; for,
low-quality borrowers are more likely to default on their loan and therefore are less affected by a
rise in the interest rate than high-quality borrowers. Lenders may then want to keep interest rates
low in order to face a better sample of borrowers.''
\stmtend
%% END-STATEMENT

\asked (a) Compute the exit rates $r^*$ and $r_{max}$, the banks' return $\rho(r)$ on each interval, its value $\rho^*$ at $r^*$ and the rate $\tilde r$ at which it recovers. (b) Compare loan demand with deposit supply at $r^*$ and count the rejected applicants by type. (c) Show what happens if a bank raises its rate to $40\%$ and explain why that does not clear the market. (d) Show that no rate up to $r^*$ clears the market, find the clearing rate among risky borrowers only, and show that it is not an equilibrium; conclude. (e) Redo the analysis with more safe borrowers (i) and with a riskier risky project (ii). (f) Ask whether rejected applicants can borrow at a higher rate, who can and who cannot.

\approach This is the Stiglitz--Weiss model of lecture block 3.6 embedded in a market. Banks pay depositors the expected return $\rho(r)$ they earn on loans, so the supply of funds is $L^s(\rho(r))$; because $\rho$ is not monotone in $r$ (adverse selection at $r^*$), the supply of loans bends backward and may never meet demand. An equilibrium needs supply equal to demand and, in addition, no bank able to offer a different rate that pays depositors more; checking that second condition is what rules out the clearing rate. Credit rationing is excess demand that persists at the equilibrium rate.

\notation
\begin{itemize}
  \item $N=1{,}000$ borrowers, $\alpha=0.75$ low-risk; $I=100$ per loan; $R_L=130$, $p_L=0.9$ (low risk), $R_H=234$, $p_H=0.5$ (high risk); payoff 0 on failure; $\bar R=p_LR_L=p_HR_H$.
  \item $r$: loan rate; $r^*=\bar R/(p_LI)-1$, $r_{max}=\bar R/(p_HI)-1$: exit rates; $m=\alpha p_L+(1-\alpha)p_H$.
  \item $\rho(r)$: banks' expected gross return per unit lent, paid to depositors; $\rho^*=\rho(r^*)$; $\tilde r$: rate above $r^*$ with $\rho(\tilde r)=\rho^*$.
  \item $L^s(\rho)=\max\{0,\,10^6(\rho-0.96)\}$: deposits supplied; $D(r)$: loan demand; $r_W$: rate clearing the market among risky borrowers.
\end{itemize}

\step{1}{Exit rates and the return function (part (a))} We restate the borrower's payoff and solve for the rates at which each type leaves. The means coincide:
\begin{align}
  \bar R = p_LR_L = (0.9)(130) = 117 = (0.5)(234) = p_HR_H. \nonumber
\end{align}
Type $i$ expects $p_i[R_i-(1+r)I]=\bar R-p_i(1+r)I$ and applies while this is non-negative. Setting it to zero,
\begin{align}
  r^* &= \frac{\bar R}{p_LI}-1 = \frac{117}{90}-1 = 1.3-1 = 0.30, \label{t3e17-rstar}\\
  r_{max} &= \frac{\bar R}{p_HI}-1 = \frac{117}{50}-1 = 2.34-1 = 1.34. \label{t3e17-rmax}
\end{align}
For $r\le r^*$ both types apply and the average success probability is
\begin{align}
  m = \alpha p_L+(1-\alpha)p_H = (0.75)(0.9)+(0.25)(0.5) = 0.675+0.125 = 0.80, \nonumber
\end{align}
so the return per unit lent is
\begin{align}
  \rho(r) = \begin{cases} m(1+r) = 0.8(1+r), & r\le0.30,\\ p_H(1+r) = 0.5(1+r), & 0.30<r\le1.34.\end{cases} \label{t3e17-rho}
\end{align}
At $r^*$, $\rho^*=(0.8)(1.3)=1.04$. The rate $\tilde r$ solves $0.5(1+\tilde r)=1.04$:
\begin{align}
  1+\tilde r = \frac{1.04}{0.5} = 2.08 \quad\Longrightarrow\quad \tilde r = 1.08. \label{t3e17-rtilde}
\end{align}

\answer{(a)} $\bar R=117$, $r^*=30\%$, $r_{max}=134\%$; $\rho(r)=0.8(1+r)$ up to $30\%$ and $0.5(1+r)$ above; $\rho^*=1.04$; $\tilde r=108\%$.

\step{2}{Excess demand at $r^*$ (part (b))} At $r^*$ the safe borrowers are just indifferent and apply, so every borrower asks for $I$:
\begin{align}
  D(r^*) &= NI = (1{,}000)(100) = 100{,}000, \nonumber\\
  L^s(\rho^*) &= 10^6(1.04-0.96) = 10^6(0.08) = 80{,}000. \nonumber
\end{align}
Excess demand is $100{,}000-80{,}000=20{,}000$, that is $20{,}000/100=200$ applicants turned down. With random selection each type is turned down in its population share: $(0.75)(200)=150$ low-risk and $(0.25)(200)=50$ high-risk.

\answer{(b)} Demand $100{,}000$, supply $80{,}000$: 800 loans, 200 applicants turned down (150 low-risk, 50 high-risk).

\step{3}{A higher rate (part (c))} At $r=0.40$, $(1+r)I=140$. The two payoffs are
\begin{align}
  U_L &= 117-(0.9)(140) = 117-126 = -9 < 0, \nonumber\\
  U_H &= 117-(0.5)(140) = 117-70 = 47 \ge 0, \nonumber
\end{align}
so only the 250 high-risk borrowers apply, and from \eqref{t3e17-rho}
\begin{align}
  \rho(0.40) &= (0.5)(1.4) = 0.70, \nonumber\\
  L^s(0.70) &= \max\{0,\,10^6(0.70-0.96)\} = \max\{0,-260{,}000\} = 0, \nonumber\\
  D(0.40) &= (250)(100) = 25{,}000. \nonumber
\end{align}

\answer{(c)} At $40\%$ only high-risk borrowers apply; $\rho=0.70<1.04$, deposits fall to 0 against a demand of $25{,}000$. The higher rate worsens the pool so much that the return per unit, and hence the supply of funds, falls: the market moves further from clearing, and the bank that raises its rate earns less.

\step{4}{No clearing rate survives (part (d))} \emph{Rates $r\le r^*$.} Demand is $100{,}000$ throughout, and $\rho(r)=0.8(1+r)\le1.04$ with $L^s$ increasing in $\rho$, so
\begin{align}
  L^s(\rho(r)) \le L^s(1.04) = 80{,}000 < 100{,}000 = D(r), \qquad r\le r^*. \nonumber
\end{align}
\emph{Rates $r^*<r\le r_{max}$.} Demand is $25{,}000$. Supply equals demand when
\begin{align}
  10^6\bigl[0.5(1+r_W)-0.96\bigr] &= 25{,}000 \nonumber\\
  0.5(1+r_W) &= 0.96+0.025 = 0.985 \nonumber\\
  1+r_W &= \frac{0.985}{0.5} = 1.97 \nonumber\\
  r_W &= 0.97, \label{t3e17-rW}
\end{align}
which lies in $(0.30,1.34]$, with $\rho(r_W)=0.985$. \emph{Deviation.} Suppose all banks charge $r_W$. A bank that offers $r^*=30\%$ instead attracts every borrower: the 750 safe ones are unserved at $97\%$ and accept $30\%$, and the risky ones prefer $30\%$ to $97\%$. Its pool is the population mix, so it earns $\rho^*=1.04>0.985$ per unit and can pay depositors more than the other banks: it attracts funds and makes a profit. So $r_W$ is not an equilibrium. Above $r_{max}$ nobody borrows. Hence no rate clears the market. With excess demand at every rate up to $r^*$, competition for deposits pushes banks to the rate with the highest return among those that keep the safe borrowers, $r^*$ (a lower rate pays depositors less and loses funds; a slightly higher one loses the safe borrowers, Step~3).

\answer{(d)} For $r\le30\%$ supply is at most $80{,}000<100{,}000$; among risky borrowers the clearing rate is $r_W=97\%$ with $\rho=0.985$, but a bank offering $30\%$ earns $1.04>0.985$, so $r_W$ is not an equilibrium. The loan rate is $r^*=30\%$ with $20{,}000$ of unmet demand: credit rationing.

\step{5}{Comparative statics (part (e))} \emph{(i) $\alpha=0.9$.} The pool improves:
\begin{align}
  m &= (0.9)(0.9)+(0.1)(0.5) = 0.81+0.05 = 0.86, \nonumber\\
  \rho^* &= (0.86)(1.3) = 1.118, \nonumber\\
  L^s(\rho^*) &= 10^6(1.118-0.96) = 158{,}000 > 100{,}000. \nonumber
\end{align}
There is excess supply at $r^*$, so the market clears below $r^*$ where $L^s(0.86(1+r))=100{,}000$:
\begin{align}
  10^6\bigl[0.86(1+r)-0.96\bigr] &= 100{,}000 \nonumber\\
  0.86(1+r) &= 1.06 \nonumber\\
  1+r &= \frac{1.06}{0.86} = 1.232558 \nonumber\\
  r &= 0.232558. \nonumber
\end{align}
Every borrower is served at $23.26\%$, so no bank can find unserved borrowers at another rate: this is an equilibrium without rationing.

\emph{(ii) $p_H=0.4$, $R_H=292.5$.} The mean is unchanged, $(0.4)(292.5)=117$, and $r^*=117/90-1=0.30$ does not change (it involves only the safe type). The other values change:
\begin{align}
  r_{max} &= \frac{117}{(0.4)(100)}-1 = 2.925-1 = 1.925, \nonumber\\
  m &= (0.75)(0.9)+(0.25)(0.4) = 0.675+0.1 = 0.775, \nonumber\\
  \rho^* &= (0.775)(1.3) = 1.0075, \nonumber\\
  L^s(\rho^*) &= 10^6(1.0075-0.96) = 47{,}500. \nonumber
\end{align}
Unmet demand is $100{,}000-47{,}500=52{,}500$, i.e.\ 525 of 1{,}000 applicants ($52.5\%$, up from $20\%$). The drop at $r^*$ is
\begin{align}
  (m-p_H)(1+r^*) = (0.775-0.4)(1.3) = 0.4875 = \alpha(p_L-p_H)(1+r^*) = (0.75)(0.5)(1.3), \nonumber
\end{align}
up from $(0.8-0.5)(1.3)=0.39$. The risky-only clearing rate solves $10^6[0.4(1+r)-0.96]=25{,}000$, so $0.4(1+r)=0.985$, $1+r=2.4625$, $r=1.4625\in(0.30,1.925]$, with return $0.985<1.0075$: a bank offering $r^*$ again does better, so it is not an equilibrium either.

\answer{(e)} (i) $\rho^*=1.118$, $L^s=158{,}000>100{,}000$: no rationing; the market clears at $23.26\%$. (ii) $r^*=30\%$, $r_{max}=192.5\%$, $\rho^*=1.0075$, $L^s=47{,}500$: $52.5\%$ of applicants turned down; the drop at $r^*$ rises from $0.39$ to $0.4875$; the risky-only clearing rate $146.25\%$ (return $0.985$) is again not an equilibrium.

\step{6}{A second, higher rate (part (f))} Of the 200 applicants turned down in (b), 150 are low-risk and 50 high-risk. At a rate $r'>r^*$ a low-risk borrower expects $117-90(1+r')<117-90(1.3)=0$ and declines; a high-risk borrower expects $117-50(1+r')\ge0$ as long as $r'\le r_{max}$ and accepts. Only the 50 high-risk borrowers accept, and the bank's return is
\begin{align}
  \rho = p_H(1+r') = 0.5(1+r') > 0.5(2.08) = 1.04 \qquad\text{for } r'>\tilde r. \nonumber
\end{align}
The offer is profitable, so other banks make it too, and competition bids $r'$ down to the rate at which the return equals what depositors earn elsewhere, $\rho^*$: $r'=\tilde r=108\%$. The deposits supplied at $\rho^*=1.04$ are then shared between the two rates, $30\%$ for the pool and $108\%$ for rejected high-risk applicants. The rejected low-risk borrowers get nothing: any rate above $r^*$ gives them a negative payoff, and at $r^*$ banks already lend all the funds depositors supply at that return.

\answer{(f)} Only the 50 rejected high-risk borrowers accept; the return $0.5(1+r')$ exceeds $1.04$ for $r'>108\%$, and competition pushes $r'$ to $\tilde r=108\%$. In this two-type example the market can have two rates; the rejected low-risk borrowers remain rationed, because no rate above $30\%$ is acceptable to them.

\intuition Raising the rate is the textbook way to clear excess demand, but here it changes who borrows. Past $r^*$ the safe borrowers leave, the return to lending falls, depositors supply less, and the gap between demand and supply widens instead of closing. Even the rate that does clear the market among risky borrowers fails, because a bank that goes back to $r^*$ gets the better pool and pays more. Rationing is therefore an equilibrium outcome, and it is worse when the pool is worse (fewer safe types, or a riskier risky project). With enough safe types the supply at $r^*$ covers demand and the market clears at a lower rate. In this two-type example the safe borrowers are exactly indifferent at $r^*$, so their exclusion costs them nothing; what the example isolates is the mechanism, why the rate does not rise to clear the market.

\pitfalls
\begin{itemize}
  \item Calling the risky-only clearing rate $r_W$ an equilibrium: supply equals demand there, but a bank offering $r^*$ earns $1.04>0.985$ and breaks it.
  \item Holding $\rho$ fixed when the rate changes: depositors are paid $\rho(r)$, which depends on who borrows, so the supply of funds moves with the composition of the pool.
  \item Forgetting that $r^*$ depends only on the safe type: a riskier risky project with the same mean changes $r_{max}$, $m$ and $\rho^*$, not $r^*$.
\end{itemize}

\refbib{Lecture Topic~3, block 3.6, slides 47--55 (backward-bending supply, two-type example). Tirole \S3.1 Introduction, p.~113 (why lenders do not raise rates when demand exceeds supply). Original exercise: no book answer.}

\exsec{18. Collateral as a separating signal --- Original exercise}
%% BEGIN-STATEMENT 18
\stmtbegin
\textbf{Origin:} original exercise on costly collateral as a signal (Bester 1985), built on
Prof.~Ricca's Topic 3 slides 56--67 (Aula/Topic\_3.pdf) and Tirole §6.3, Application 5, pp.~251--253
(PDF 264--266) · Answer: in the solutions file

An entrepreneur needs \(I = 60\) for a project that yields \(R = 120\) if it succeeds and 0 if it
fails. She has no cash, but she owns an illiquid asset worth 60 to her that she can pledge as
collateral: collateral worth \(C\) to her is worth only \(\delta C\) to a lender, with
\(\delta = 0.5\). A good type succeeds with probability \(p = 0.9\), a bad type with probability
\(q = 0.6\); the share of good types is \(\alpha = 0.5\). Everybody is risk neutral, the riskless
rate is 0, and lenders are competitive (zero expected profit). A contract \(\{R_u, C\}\) says that
the lender receives \(R_u\) if the project succeeds and seizes the collateral \(C\) if it fails; the
entrepreneur keeps \(R - R_u\) on success and loses \(C\) on failure.

\begin{enumerate}
\def\labelenumi{(\alph{enumi})}
\item
  Symmetric information. Compute each type's zero-profit repayment with no collateral, \(R_u^G\) and
  \(R_u^B\), and each type's utility. Then let the good type pledge \(C = 20\): compute her
  zero-profit repayment and her utility, and show that her utility falls by
  \((1 - p)(1 - \delta)C\). Why is \(C = 0\) optimal under symmetric information?
\item
  Asymmetric information without collateral. Compute the prior \(m\) and the pooling repayment
  \(R_u = I/m\), and each type's utility. Show that if the good type asked for her
  symmetric-information contract \(\{I/p, 0\}\), the bad type would mimic her.
\item
  Separating equilibrium. In it the bad type pledges no collateral and repays \(R_u^B = I/q\). The
  good type chooses \(\{\tilde R, \tilde C\}\) to maximise her utility subject to the no-mimicking
  condition of the bad type and the lender's break-even condition. Write the program, explain why
  both constraints bind, and solve for \(C_G\) and \(R_u^G\) (give the general formulas and the
  numbers). Verify both constraints numerically and check that the good type can afford the
  collateral.
\item
  Compute the good type's utility in the separating equilibrium and the expected deadweight cost of
  collateral, \((1 - p)(1 - \delta)C_G\). Compare it with her utility in (a) and in (b). Which
  regime does the good type prefer, and by how much?
\item
  Check that neither type wants the other type's contract in the separating equilibrium. Explain why
  collateral can separate the types here (which type finds it more costly, and why).
\end{enumerate}

\paragraph{Context needed to solve it}

\emph{From Prof.~Ricca's Topic 3 slides 56--67 (block 3.7, signalling with costly collateral):}

\begin{itemize}
\tightlist
\item
  To be a successful separating mechanism, the cost of the signalling device has to differ
  sufficiently among the types (slide 56).
\item
  Projects require \(I\) and yield \(R\) on success and 0 otherwise; the good type succeeds with
  probability \(p\), the bad type with \(q\); the share of good types is \(\alpha\) and the prior is
  \(m = \alpha p + (1 - \alpha)q\); both types have positive NPV, \(pR > I\), \(qR > I\) (slide 57).
\item
  Borrowers have no cash on hand but can pledge an illiquid asset as collateral. Assets worth \(C\)
  in the hands of borrowers are worth only \(\delta C\) in the lenders' hands, \(0 < \delta < 1\);
  collateral is inefficient: \((1 - \delta)C\) is lost in case of default (slide 58). (Tirole writes
  \(\beta\) for \(\delta\).)
\item
  Symmetric information: with collateral \(C\) the good type's zero-profit repayment satisfies
  \(pR_u + (1 - p)\delta C = I\), i.e.~\(R_u = I/p - (1 - p)\delta C/p\), and her utility is
  \(p(R - R_u) + (1 - p)(-C) = pR - I - (1 - p)(1 - \delta)C\), so the best choice is \(C = 0\);
  same for the bad type (slides 59--60).
\item
  Asymmetric information without collateral: both types repay \(R_u = I/m > I/p = R_u^G\) (slide
  61).
\item
  Separating equilibrium: good type \(\{R_u^G, C_G\}\), bad type \(\{R_u^B, C_B\}\); lenders infer
  the type from the contract. The bad type must not want to mimic, so \(C_B = 0\) and
  \(R_u^B = I/q\). The good type's contract \(\{\tilde R, \tilde C\}\) solves
  \[\max_{\tilde R, \tilde C}\ p(R - \tilde R) + (1 - p)(-\tilde C)\] subject to
  \(q(R - \tilde R) + (1 - q)(-\tilde C) \le qR - I\) (no-mimicking) and
  \(p\tilde R + (1 - p)\delta\tilde C \ge I\) (lender's break-even) (slides 62--65).
\item
  Both constraints bind at the solution (slide 66). Signalling is worth it for the good type when
  the adverse-selection cost is large: \(p(R - R_u^G) + (1 - p)(-C_G) \ge p(R - I/m)\) (slide 67).
\end{itemize}

\emph{From Tirole §6.3, Application 5 ``Costly Collateral Pledging'', pp.~251--253 (PDF 264--266):}
the same program written with the borrower's reward \(R_b = R - R_u\); Assumption 6.1,
\(0 < qR - I < pR - I\); the borrower has ``a sizeable amount of assets'' to pledge; signalling
works because pledging collateral is relatively more costly for a bad borrower, while a higher
reward in success is valued more by a good one (the sorting condition).
\stmtend
%% END-STATEMENT

\asked (a) Under symmetric information, find each type's zero-profit repayment and utility, and show that pledging collateral only lowers the good type's utility. (b) Without collateral and with asymmetric information, find the pooled repayment and the utilities, and show that the bad type would copy the good type's symmetric-information contract. (c) Solve the good type's least-cost separating contract (repayment and collateral), explain why both constraints bind, and verify them. (d) Compute the good type's utility under separation, the deadweight cost of collateral, and compare with (a) and (b). (e) Check that neither type wants the other's contract and explain why collateral sorts the types.

\approach This is the costly-collateral signalling model of lecture block 3.7 (Bester 1985; Tirole \S6.3, Application 5). Collateral is a dissipative signal: the lender values it at $\delta C<C$, so it destroys value in default and is never used under symmetric information. Under asymmetric information it is cheaper for the good type, who defaults less often, so a contract with collateral and a lower repayment can be made unattractive to the bad type. The good type's best separating contract maximises her utility subject to the bad type's no-mimicking constraint and the lender's break-even constraint.

\notation
\begin{itemize}
  \item $I=60$: investment; $R=120$: return on success (0 on failure); $p=0.9$, $q=0.6$: success probabilities of the good and the bad type; $\alpha=0.5$: share of good types; $m=\alpha p+(1-\alpha)q$.
  \item $C$: collateral, worth $C$ to the borrower and $\delta C$ to the lender, $\delta=0.5$; the borrower owns assets worth 60.
  \item $R_u$: repayment to the lender on success; $R_u^G=I/p$, $R_u^B=I/q$: symmetric-information repayments; $\{\tilde R,\tilde C\}$: the good type's contract in the separating program, with solution $\{R_u^G,C_G\}$ (written $R_u^{G,sep}$ when it must be told apart from $I/p$).
  \item $U_G$, $U_B$: expected utilities of the good and the bad type; $D\equiv p(1-q)-\delta q(1-p)$.
\end{itemize}

\step{1}{Symmetric information (part (a))} A lender who knows the type breaks even when expected repayment equals $I$. Without collateral,
\begin{align}
  R_u^G &= \frac Ip = \frac{60}{0.9} = 66.667, & R_u^B &= \frac Iq = \frac{60}{0.6} = 100, \nonumber\\
  U_G &= p(R-R_u^G) = pR-I = 108-60 = 48, & U_B &= qR-I = 72-60 = 12. \nonumber
\end{align}
With collateral $C=20$ for the good type, break-even requires $pR_u+(1-p)\delta C=I$:
\begin{align}
  0.9R_u+(0.1)(0.5)(20) &= 60 \nonumber\\
  0.9R_u &= 60-1 = 59 \nonumber\\
  R_u &= \frac{59}{0.9} = 65.556, \nonumber
\end{align}
and her utility is
\begin{align}
  p(R-R_u)-(1-p)C = 0.9(120-65.556)-(0.1)(20) = 49-2 = 47. \nonumber
\end{align}
In general, substituting $pR_u=I-(1-p)\delta C$ from the break-even condition,
\begin{align}
  p(R-R_u)-(1-p)C &= pR-\bigl[I-(1-p)\delta C\bigr]-(1-p)C \nonumber\\
  &= pR-I-(1-p)(1-\delta)C, \label{t3e18-symC}
\end{align}
which here is $48-(0.1)(0.5)(20)=48-1=47$.

\answer{(a)} $R_u^G=66.667$, $R_u^B=100$, $U_G=48$, $U_B=12$. With $C=20$: $R_u=65.556$, $U_G=47=48-(1-p)(1-\delta)C$. Utility falls one for one with the expected collateral loss $(1-p)(1-\delta)C$, so $C=0$ is optimal; the same argument holds for the bad type.

\step{2}{Pooling without collateral (part (b))} Lenders who cannot tell the types apart price the loan on the prior:
\begin{align}
  m &= \alpha p+(1-\alpha)q = (0.5)(0.9)+(0.5)(0.6) = 0.45+0.30 = 0.75, \nonumber\\
  R_u &= \frac Im = \frac{60}{0.75} = 80, \nonumber\\
  U_G &= p(R-R_u) = 0.9(120-80) = 36, \nonumber\\
  U_B &= q(R-R_u) = 0.6(120-80) = 24. \nonumber
\end{align}
The good type loses $48-36=12$ and the bad type gains $24-12=12$ relative to symmetric information (with equal shares the transfers offset: $(0.5)(12)=(0.5)(12)$). If the good type asked for $\{I/p,0\}$, the bad type would get
\begin{align}
  q(R-I/p) = 0.6(120-66.667) = 0.6(53.333) = 32 > 12, \nonumber
\end{align}
so she would mimic, and the lender, facing both types, would receive $m(66.667)=(0.75)(66.667)=50<60$.

\answer{(b)} $m=0.75$, $R_u=80$, $U_G=36$, $U_B=24$ (a cross-subsidy of 12 from good to bad). The bad type gets $32>12$ by copying $\{66.667,0\}$, so the symmetric-information contract cannot be offered to the good type alone.

\step{3}{The separating program and why both constraints bind (part (c))} The bad type, recognised, pledges nothing and repays $I/q=100$, getting $qR-I=12$. The good type solves
\begin{align}
  \max_{\tilde R,\tilde C}\ & p(R-\tilde R)-(1-p)\tilde C \nonumber\\
  \text{s.t. }\ & q(R-\tilde R)-(1-q)\tilde C \le qR-I \qquad\text{(no-mimicking)}, \label{t3e18-NM}\\
  & p\tilde R+(1-p)\delta\tilde C \ge I \qquad\text{(break-even)}. \label{t3e18-BE}
\end{align}
\emph{No-mimicking binds.} Without \eqref{t3e18-NM} the problem is the symmetric-information one, whose solution is $\{I/p,0\}$ by \eqref{t3e18-symC}; Step~2 showed that this violates \eqref{t3e18-NM} ($32>12$). So \eqref{t3e18-NM} binds. \emph{Break-even binds.} Expanding \eqref{t3e18-NM} at equality, $qR-q\tilde R-(1-q)\tilde C=qR-I$, i.e.\ $q\tilde R+(1-q)\tilde C=I$. Along this line $d\tilde C=-\frac{q}{1-q}\,d\tilde R$, and the good type's utility changes by
\begin{align}
  -p\,d\tilde R-(1-p)\,d\tilde C &= -p\,d\tilde R+\frac{q(1-p)}{1-q}\,d\tilde R \nonumber\\
  &= -\frac{p(1-q)-q(1-p)}{1-q}\,d\tilde R = -\frac{p-q}{1-q}\,d\tilde R, \nonumber
\end{align}
so she gains by lowering $\tilde R$ and raising $\tilde C$ along the line. The lender's receipts change by
\begin{align}
  p\,d\tilde R+(1-p)\delta\,d\tilde C = p\,d\tilde R-\frac{\delta q(1-p)}{1-q}\,d\tilde R = \frac{p(1-q)-\delta q(1-p)}{1-q}\,d\tilde R, \nonumber
\end{align}
which falls as $\tilde R$ falls (the coefficient is positive because $p(1-q)>q(1-p)>\delta q(1-p)$). The good type therefore moves along the line until the lender just breaks even: \eqref{t3e18-BE} binds.

\step{4}{Solving the two equations} The binding constraints are
\begin{align}
  q\tilde R+(1-q)\tilde C &= I, \label{t3e18-NMb}\\
  p\tilde R+(1-p)\delta\tilde C &= I. \label{t3e18-BEb}
\end{align}
From \eqref{t3e18-NMb}, $\tilde R=[I-(1-q)\tilde C]/q$. Substituting into \eqref{t3e18-BEb} and multiplying by $q$,
\begin{align}
  \frac pq\bigl[I-(1-q)\tilde C\bigr]+(1-p)\delta\tilde C &= I \nonumber\\
  pI-p(1-q)\tilde C+\delta q(1-p)\tilde C &= qI \nonumber\\
  (p-q)I &= \bigl[p(1-q)-\delta q(1-p)\bigr]\tilde C \nonumber\\
  C_G &= \frac{(p-q)I}{p(1-q)-\delta q(1-p)} = \frac{(p-q)I}{D}. \label{t3e18-C}
\end{align}
Since $p(1-q)=(p-q)+q(1-p)$, $D=(p-q)+q(1-p)(1-\delta)$ and $C_G=I/[1+q(1-p)(1-\delta)/(p-q)]$, the slides' form. Back in \eqref{t3e18-NMb},
\begin{align}
  R_u^G = \frac{I-(1-q)C_G}{q} = \frac Iq\cdot\frac{D-(1-q)(p-q)}{D}, \nonumber
\end{align}
and $D-(1-q)(p-q)=p(1-q)-\delta q(1-p)-p(1-q)+q(1-q)=q[(1-q)-\delta(1-p)]$, so
\begin{align}
  R_u^G = \frac{(1-q)-\delta(1-p)}{p(1-q)-\delta q(1-p)}\,I. \label{t3e18-R}
\end{align}
Numbers:
\begin{align}
  D &= (0.9)(0.4)-(0.5)(0.6)(0.1) = 0.36-0.03 = 0.33, \nonumber\\
  C_G &= \frac{(0.3)(60)}{0.33} = \frac{18}{0.33} = 54.545, \nonumber\\
  R_u^G &= \frac{0.4-(0.5)(0.1)}{0.33}(60) = \frac{(0.35)(60)}{0.33} = \frac{21}{0.33} = 63.636. \nonumber
\end{align}
Checks: break-even, $(0.9)(63.636)+(0.1)(0.5)(54.545)=57.273+2.727=60$; no-mimicking, $0.6(120-63.636)-(0.4)(54.545)=33.818-21.818=12=qR-I$; the collateral $54.545$ is below the assets of 60, and $R_u^G<R$.

\answer{(c)} Both constraints bind; $C_G=(p-q)I/[p(1-q)-\delta q(1-p)]=54.545$ and $R_u^G=[(1-q)-\delta(1-p)]I/[p(1-q)-\delta q(1-p)]=63.636$ (below $I/p=66.667$, because the collateral pays the lender something in default).

\step{5}{Utility and deadweight cost (part (d))} Directly,
\begin{align}
  U_G^{sep} = p(R-R_u^G)-(1-p)C_G = 0.9(56.364)-(0.1)(54.545) = 50.727-5.455 = 45.273. \nonumber
\end{align}
By \eqref{t3e18-symC}, with the break-even condition binding, this equals $pR-I-(1-p)(1-\delta)C_G$, where
\begin{align}
  (1-p)(1-\delta)C_G = (0.1)(0.5)(54.545) = 2.727, \qquad 48-2.727 = 45.273. \nonumber
\end{align}

\answer{(d)} $U_G^{sep}=45.273$; deadweight cost $2.727$. The good type is $2.727$ below her symmetric-information utility of 48 but $45.273-36=9.273$ above pooling: she prefers to signal.

\step{6}{Incentives and sorting (part (e))} The bad type gets $qR-I=12$ from her own contract and, by the binding \eqref{t3e18-NM}, also 12 from the good type's contract: she is indifferent and takes her own (any tiny extra collateral would make the preference strict). The good type taking the bad contract would get $p(R-I/q)=0.9(120-100)=18<45.273$. Collateral costs type $i$ its failure probability times $C$ ($0.1C$ for the good type, $0.4C$ for the bad), while a lower repayment benefits type $i$ by its success probability ($0.9$ vs.\ $0.6$). Along an indifference curve, the repayment cut needed to compensate one unit of collateral is $0.1/0.9=0.111$ for the good type and $0.4/0.6=0.667$ for the bad type: the curves cross once, so a contract with more collateral and a lower repayment attracts only the good type.

\answer{(e)} Bad type: 12 from either contract (indifferent, takes her own); good type: $18<45.273$ from the bad contract. Collateral separates because it is relatively more costly to the type that defaults more often (single crossing).

\intuition Collateral is a pure cost in this model: the lender values it at half its worth to the borrower, so it is never used when types are known. When they are not, the good type is willing to burn part of its value, because the burning is cheap for her (she rarely fails) and expensive for the bad type (who fails 40\% of the time). The least-cost contract piles on just enough collateral to make the bad type indifferent and lowers the repayment until the lender only breaks even. The good type pays $2.727$ in expected deadweight cost instead of a cross-subsidy of $12$.

\pitfalls
\begin{itemize}
  \item Writing the break-even condition with $C$ instead of $\delta C$: the lender recovers only $\delta C$ in default.
  \item Letting the bad type pledge collateral too: in the separating equilibrium she is recognised, and pledging only destroys value for her, so $C_B=0$, $R_u^B=I/q$.
  \item Comparing the separating repayment with the pooling one without counting the collateral: the right comparison is utilities, $45.273$ vs.\ $36$.
\end{itemize}

\refbib{Lecture Topic~3, block 3.7, slides 56--67 (costly collateral pledging). Tirole \S6.3, Application 5 ``Costly Collateral Pledging'', pp.~251--253 (equations (6.3)--(6.5); Tirole writes $\beta$ for $\delta$ and $R_b=R-R_u$). Original exercise: no book answer.}

\exsec{19. When does signalling with collateral pay? --- Original exercise}
%% BEGIN-STATEMENT 19
\stmtbegin
\textbf{Origin:} original exercise on costly collateral as a signal (Bester 1985): separating versus
pooling and comparative statics, built on Prof.~Ricca's Topic 3 slides 56--67 (Aula/Topic\_3.pdf)
and Tirole §6.3, Application 5, pp.~251--253 (PDF 264--266) · Answer: in the solutions file

Same model as exercise 18 (restated in the context below) with \(I = 50\), \(R = 100\), \(p = 0.8\),
\(q = 0.7\) and, unless stated otherwise, \(\delta = 0.2\). The entrepreneur owns illiquid assets
worth 50 to her. The share of good types \(\alpha\) varies across parts.

\begin{enumerate}
\def\labelenumi{(\alph{enumi})}
\item
  Compute the good type's least-cost separating contract \(\{R_u^G, C_G\}\) (both constraints
  binding), check both constraints, and compute her utility \(U^{sep}\) and the expected deadweight
  cost \((1 - p)(1 - \delta)C_G\). Does \(U^{sep}\) depend on \(\alpha\)?
\item
  For \(\alpha = 0.3\) and for \(\alpha = 0.6\), compute \(m\), the pooling repayment \(I/m\) and
  the good type's pooling utility \(U^{pool} = p(R - I/m)\). In which case does she prefer to
  signal?
\item
  Show, for general parameters, that \(U^{sep} \ge U^{pool}\) if and only if
  \[\alpha\,p\,(1 - p)(1 - \delta) \le (1 - \alpha)(p - q),\] and write the threshold \(\hat\alpha\)
  below which the good type prefers to signal. Compute \(\hat\alpha\) and check that at
  \(\hat\alpha\) the two utilities coincide.
\item
  Comparative statics in \(\delta\). Show that \(C_G\) increases with \(\delta\) while the expected
  deadweight cost \((1 - p)(1 - \delta)C_G\) decreases with \(\delta\); evaluate both derivatives at
  \(\delta = 0.2\). Compute \(C_G\), \(R_u^G\), the deadweight cost and \(U^{sep}\) for
  \(\delta = 0\) and \(\delta = 0.6\), and the limits as \(\delta \to 1\). For \(\alpha = 0.6\),
  find the threshold \(\hat\delta\) above which the good type prefers to signal, and check it with
  the numbers for \(\delta = 0.6\).
\item
  Comparative statics in \(q\) (\(\delta = 0.2\)). Show that \(C_G\) falls as \(q\) rises toward
  \(p\), and compute \(C_G\) and the deadweight cost for \(q = 0.6\) and \(q = 0.75\). Interpret.
\item
  Welfare at \(\alpha = 0.3\). Compute total expected surplus per borrower (both types together)
  under pooling and under separation, the bad type's utility in each, and the surplus lost to
  separation. Which outcome is efficient, and why does the good type still choose to signal?
\end{enumerate}

\paragraph{Context needed to solve it}

\emph{From Prof.~Ricca's Topic 3 slides 56--67 (block 3.7, signalling with costly collateral):} the
model of exercise 18, restated.

\begin{itemize}
\tightlist
\item
  Project: cost \(I\), return \(R\) on success, 0 on failure; good type succeeds with probability
  \(p\), bad type with \(q < p\); share of good types \(\alpha\), prior
  \(m = \alpha p + (1 - \alpha)q\); \(pR > I\) and \(qR > I\); no cash on hand; risk neutrality,
  riskless rate 0, competitive lenders (slide 57).
\item
  Collateral worth \(C\) to the borrower is worth \(\delta C\) to the lender, \(0 < \delta < 1\)
  (Tirole: \(\beta\), \(0 \le \beta < 1\)); \((1 - \delta)C\) is lost in default (slide 58). Under
  symmetric information the good type's utility with collateral is
  \(pR - I - (1 - p)(1 - \delta)C\), so \(C = 0\) (slides 59--60).
\item
  Pooling without collateral: \(R_u = I/m\), good type's utility \(p(R - I/m)\) (slides 61, 67).
\item
  Separating: bad type \(C_B = 0\), \(R_u^B = I/q\), utility \(qR - I\). Good type:
  \(\max_{\tilde R, \tilde C} p(R - \tilde R) - (1 - p)\tilde C\) subject to
  \(q(R - \tilde R) - (1 - q)\tilde C \le qR - I\) and \(p\tilde R + (1 - p)\delta\tilde C \ge I\);
  both bind (slides 63--66), giving
  \[R_u^G = \frac{(1 - q) - \delta(1 - p)}{p(1 - q) - \delta q(1 - p)}\,I, \qquad C_G = \frac{I}{1 + \dfrac{q(1 - p)(1 - \delta)}{p - q}}.\]
\item
  If the adverse-selection cost is very large (large share of bad borrowers \(1 - \alpha\), or large
  distance between \(p\) and \(q\)), the good type is better off with the possibility of signalling:
  \(p(R - R_u^G) + (1 - p)(-C_G) \ge p(R - I/m)\) (slide 67).
\end{itemize}

\emph{From Tirole §6.3, Application 5, p.~253 (PDF 266):} ``The good borrower must pledge more
collateral when collateral pledging becomes cheaper for the borrower
(\(\partial C^*/\partial\beta > 0\))'' and ``the stronger the asymmetry of information
(\(\partial C^*/\partial q < 0\)) {[}\ldots{]} \(C^*\) tends to 0 {[}\ldots{]} when \(q\) tends to
\(p\), as we would expect.''
\stmtend
%% END-STATEMENT

\asked (a) Find the good type's least-cost separating contract, her utility and the deadweight cost, and say whether they depend on the share of good types. (b) Compute her pooling utility for $\alpha=0.3$ and $\alpha=0.6$ and pick the better regime in each case. (c) Derive, in general, the condition under which separation beats pooling, its threshold $\hat\alpha$, and check it numerically. (d) Show how the collateral and the deadweight cost move with $\delta$, compute the contract for $\delta=0$ and $\delta=0.6$ and as $\delta\to1$, and find the threshold $\hat\delta$ for $\alpha=0.6$. (e) Show how the collateral moves with $q$ and compute it for $q=0.6$ and $0.75$. (f) Compare total surplus and the bad type's utility under the two regimes at $\alpha=0.3$.

\approach Same costly-collateral model as in lecture block 3.7 (Bester 1985; Tirole \S6.3, Application 5). The separating contract does not depend on $\alpha$ (lenders know the type once contracts separate), while the pooling utility $p(R-I/m)$ rises with $\alpha$. Comparing the two (slide 67) gives a threshold; comparative statics use derivatives of the closed-form $C_G$. Welfare compares total surplus, where pooling has no deadweight cost because both types have positive NPV and are financed either way.

\notation
\begin{itemize}
  \item $I=50$, $R=100$ (0 on failure), $p=0.8$, $q=0.7$; $\delta=0.2$ unless stated (collateral worth $C$ to the borrower is worth $\delta C$ to the lender); assets worth 50.
  \item $\alpha$: share of good types; $m=\alpha p+(1-\alpha)q$; $U^{pool}=p(R-I/m)$, $U^{sep}$: good type's utility in each regime.
  \item $\{R_u^G,C_G\}$: good type's separating repayment and collateral; the bad type repays $I/q$ with no collateral and gets $qR-I$.
  \item $x\equiv(1-p)(1-\delta)$: expected deadweight loss per unit of collateral; $D\equiv p(1-q)-\delta q(1-p)=(p-q)+qx$; $DW\equiv xC_G$: expected deadweight cost; $\hat\alpha$, $\hat\delta$: thresholds.
\end{itemize}

\step{1}{The separating contract (part (a))} The good type maximises $p(R-\tilde R)-(1-p)\tilde C$ subject to the bad type's no-mimicking constraint $q(R-\tilde R)-(1-q)\tilde C\le qR-I$ and the lender's break-even $p\tilde R+(1-p)\delta\tilde C\ge I$. No-mimicking binds because the solution without it, $\{I/p,0\}$, gives the bad type $q(R-I/p)=0.7(100-62.5)=26.25>20=qR-I$. Break-even binds because along the binding no-mimicking line ($d\tilde C=-\frac q{1-q}d\tilde R$) the good type's utility changes by $-\frac{p-q}{1-q}d\tilde R$, so she lowers $\tilde R$ and raises $\tilde C$ until the lender's receipts, which change by $\frac{D}{1-q}d\tilde R$ with $D>0$, fall to $I$. The two equalities are
\begin{align}
  q\tilde R+(1-q)\tilde C &= I, \label{t3e19-NM}\\
  p\tilde R+(1-p)\delta\tilde C &= I. \label{t3e19-BE}
\end{align}
Substituting $\tilde R=[I-(1-q)\tilde C]/q$ from \eqref{t3e19-NM} into \eqref{t3e19-BE} and multiplying by $q$: $pI-p(1-q)\tilde C+\delta q(1-p)\tilde C=qI$, so $(p-q)I=D\tilde C$ and
\begin{align}
  C_G = \frac{(p-q)I}{D}, \qquad R_u^G = \frac{I-(1-q)C_G}{q} = \frac{(1-q)-\delta(1-p)}{D}\,I \label{t3e19-sol}
\end{align}
(the last step uses $D-(1-q)(p-q)=q[(1-q)-\delta(1-p)]$). Numbers:
\begin{align}
  D &= (0.8)(0.3)-(0.2)(0.7)(0.2) = 0.24-0.028 = 0.212, \nonumber\\
  C_G &= \frac{(0.1)(50)}{0.212} = \frac5{0.212} = 23.585, \nonumber\\
  R_u^G &= \frac{0.3-(0.2)(0.2)}{0.212}(50) = \frac{(0.26)(50)}{0.212} = \frac{13}{0.212} = 61.321. \nonumber
\end{align}
Checks: $(0.8)(61.321)+(0.2)(0.2)(23.585)=49.057+0.943=50$; $0.7(100-61.321)-(0.3)(23.585)=27.075-7.075=20=qR-I$. Using the binding break-even, the good type's utility is
\begin{align}
  U^{sep} &= p(R-R_u^G)-(1-p)C_G = pR-I-(1-p)(1-\delta)C_G, \label{t3e19-Usep}\\
  DW &= (0.2)(0.8)(23.585) = 3.774, \qquad U^{sep} = 80-50-3.774 = 26.226. \nonumber
\end{align}

\answer{(a)} $C_G=23.585$, $R_u^G=61.321$, $DW=3.774$, $U^{sep}=26.226$. None depends on $\alpha$: once contracts separate, lenders know each borrower's type.

\step{2}{Pooling utilities (part (b))} The pooled repayment is $I/m$ and the good type gets $p(R-I/m)$:
\begin{align}
  \alpha=0.3:\quad m &= (0.3)(0.8)+(0.7)(0.7) = 0.24+0.49 = 0.73, \nonumber\\
  I/m &= 50/0.73 = 68.493, \qquad U^{pool} = 0.8(100-68.493) = 0.8(31.507) = 25.205; \nonumber\\
  \alpha=0.6:\quad m &= (0.6)(0.8)+(0.4)(0.7) = 0.48+0.28 = 0.76, \nonumber\\
  I/m &= 50/0.76 = 65.789, \qquad U^{pool} = 0.8(100-65.789) = 0.8(34.211) = 27.368. \nonumber
\end{align}
Both pools are financeable ($mR=73$ and $76$ exceed $I=50$).

\answer{(b)} $\alpha=0.3$: $U^{pool}=25.205<26.226$, she signals. $\alpha=0.6$: $U^{pool}=27.368>26.226$, she prefers pooling.

\step{3}{The general condition (part (c))} Using \eqref{t3e19-Usep}, separation is at least as good as pooling iff
\begin{align}
  pR-I-xC_G &\ge pR-\frac{pI}{m} \nonumber\\
  xC_G &\le \frac{pI}m-I = \frac{(p-m)I}{m}. \label{t3e19-cmp}
\end{align}
Since $p-m=p-\alpha p-(1-\alpha)q=(1-\alpha)(p-q)$ and $C_G=(p-q)I/D$, \eqref{t3e19-cmp} reads
\begin{align}
  \frac{x(p-q)I}{D} \le \frac{(1-\alpha)(p-q)I}{m}. \nonumber
\end{align}
Dividing by $(p-q)I>0$ and multiplying by $mD>0$,
\begin{align}
  xm &\le (1-\alpha)D = (1-\alpha)(p-q)+(1-\alpha)qx \nonumber\\
  x\bigl[m-(1-\alpha)q\bigr] &\le (1-\alpha)(p-q). \nonumber
\end{align}
Because $m-(1-\alpha)q=\alpha p$, this is
\begin{align}
  \alpha p(1-p)(1-\delta) \le (1-\alpha)(p-q). \label{t3e19-cond}
\end{align}
Collecting the terms in $\alpha$, $\alpha[p(1-p)(1-\delta)+(p-q)]\le p-q$, so
\begin{align}
  \hat\alpha = \frac{p-q}{(p-q)+p(1-p)(1-\delta)} = \frac{0.1}{0.1+(0.8)(0.2)(0.8)} = \frac{0.1}{0.228} = 0.438596. \label{t3e19-ahat}
\end{align}
Check: at $\hat\alpha$, $m=0.7+0.1\hat\alpha=0.743860$, $I/m=67.217$, and $U^{pool}=0.8(100-67.217)=0.8(32.783)=26.226=U^{sep}$. Also, for $\alpha=0.3$ the two sides of \eqref{t3e19-cond} are $(0.3)(0.128)=0.0384\le0.07$, and for $\alpha=0.6$ they are $0.0768>0.04$, consistent with (b).

\answer{(c)} $U^{sep}\ge U^{pool}\iff\alpha p(1-p)(1-\delta)\le(1-\alpha)(p-q)\iff\alpha\le\hat\alpha=(p-q)/[(p-q)+p(1-p)(1-\delta)]=0.4386$. The good type signals when bad types are numerous (pooling is expensive) or collateral is cheap to burn.

\step{4}{Comparative statics in $\delta$ (part (d))} With $D(\delta)=(p-q)+q(1-p)(1-\delta)$, $D'(\delta)=-q(1-p)$, so
\begin{align}
  \frac{dC_G}{d\delta} &= -\frac{(p-q)I\,D'(\delta)}{D^2} = \frac{(p-q)I\,q(1-p)}{D^2} > 0, \nonumber\\
  \left.\frac{dC_G}{d\delta}\right|_{\delta=0.2} &= \frac{(0.1)(50)(0.7)(0.2)}{0.212^2} = \frac{0.7}{0.044944} = 15.575. \nonumber
\end{align}
For the deadweight cost write $DW=xC_G=(p-q)Ix/[(p-q)+qx]$. Differentiating in $x$ (quotient rule),
\begin{align}
  \frac{dDW}{dx} = (p-q)I\,\frac{(p-q)+qx-qx}{D^2} = \frac{(p-q)^2I}{D^2} > 0, \nonumber
\end{align}
and $dx/d\delta=-(1-p)$, so
\begin{align}
  \frac{dDW}{d\delta} &= -\frac{(1-p)(p-q)^2I}{D^2} < 0, \nonumber\\
  \left.\frac{dDW}{d\delta}\right|_{\delta=0.2} &= -\frac{(0.2)(0.01)(50)}{0.044944} = -\frac{0.1}{0.044944} = -2.225. \nonumber
\end{align}
Evaluating \eqref{t3e19-sol} and \eqref{t3e19-Usep} (with $D=0.24-0.14\delta$, $x=0.2(1-\delta)$ and $U^{sep}=30-DW$):
\begin{center}\small
\begin{tabular}{@{}lccccc@{}}
\toprule
$\delta$ & $D$ & $C_G=5/D$ & $R_u^G=50[0.3-0.2\delta]/D$ & $DW=xC_G$ & $U^{sep}$\\
\midrule
0 & $0.24$ & $5/0.24=20.833$ & $15/0.24=62.5$ & $(0.2)(20.833)=4.167$ & $25.833$\\
0.2 & $0.212$ & $5/0.212=23.585$ & $13/0.212=61.321$ & $(0.16)(23.585)=3.774$ & $26.226$\\
0.6 & $0.24-0.084=0.156$ & $5/0.156=32.051$ & $9/0.156=57.692$ & $(0.08)(32.051)=2.564$ & $27.436$\\
\bottomrule
\end{tabular}
\end{center}
At $\delta=0$, $R_u^G=I/p=62.5$: collateral worth nothing to the lender cannot lower the repayment, it only deters the bad type.
As $\delta\to1$: $D\to p-q$, so $C_G\to I=50$, $R_u^G\to I(p-q)/(p-q)=I=50$, $DW\to0$ and $U^{sep}\to pR-I=30$ (the loan is repaid in full in both states). For $\alpha=0.6$, solve \eqref{t3e19-cond} for $\delta$:
\begin{align}
  1-\delta \le \frac{(1-\alpha)(p-q)}{\alpha p(1-p)} = \frac{(0.4)(0.1)}{(0.6)(0.8)(0.2)} = \frac{0.04}{0.096} = 0.416667 \iff \delta \ge \hat\delta = 0.583333. \nonumber
\end{align}
At $\delta=0.6>\hat\delta$, $U^{sep}=27.436>27.368=U^{pool}$, as predicted.

\answer{(d)} $dC_G/d\delta=15.575>0$ and $dDW/d\delta=-2.225<0$ at $\delta=0.2$. $\delta=0$: $C_G=20.833$, $R_u^G=62.5$, $DW=4.167$, $U^{sep}=25.833$. $\delta=0.6$: $C_G=32.051$, $R_u^G=57.692$, $DW=2.564$, $U^{sep}=27.436$. $\delta\to1$: $C_G\to50$, $R_u^G\to50$, $DW\to0$, $U^{sep}\to30$. For $\alpha=0.6$: $\hat\delta=0.5833$.

\step{5}{Comparative statics in $q$ (part (e))} Here $x$ does not depend on $q$ and $dD/dq=-1+x$. By the quotient rule,
\begin{align}
  \frac{dC_G}{dq} &= \frac{-I\cdot D-(p-q)I(-1+x)}{D^2} = \frac{I\bigl[-(p-q)-qx+(p-q)-(p-q)x\bigr]}{D^2} \nonumber\\
  &= -\frac{Ipx}{D^2} < 0, \nonumber\\
  \left.\frac{dC_G}{dq}\right|_{q=0.7} &= -\frac{(50)(0.8)(0.16)}{0.044944} = -\frac{6.4}{0.044944} = -142.40. \nonumber
\end{align}
At $q=0.6$: $D=0.2+(0.6)(0.16)=0.296$, $C_G=(0.2)(50)/0.296=33.784$, $DW=(0.16)(33.784)=5.405$ (and $qR-I=10>0$, so the bad type is still creditworthy). At $q=0.75$: $D=0.05+(0.75)(0.16)=0.17$, $C_G=(0.05)(50)/0.17=14.706$, $DW=(0.16)(14.706)=2.353$. As $q\to p$, $C_G\to0$.

\answer{(e)} $dC_G/dq=-Ipx/D^2<0$; $C_G=33.784$ ($DW=5.405$) at $q=0.6$ and $14.706$ ($DW=2.353$) at $q=0.75$. A wider gap between the types gives the bad type more to gain from mimicking, so more collateral is needed to deter her; the result assumes both types stay creditworthy ($qR>I$).

\step{6}{Welfare at $\alpha=0.3$ (part (f))} Lenders break even in both regimes, so total surplus per borrower is the average of the two types' utilities. Under pooling nothing is destroyed:
\begin{align}
  \alpha(pR-I)+(1-\alpha)(qR-I) &= (0.3)(30)+(0.7)(20) = 9+14 = 23, \nonumber\\
  \text{check: } \alpha U^{pool}+(1-\alpha)q(R-I/m) &= (0.3)(25.205)+(0.7)(0.7)(31.507) = 7.562+15.438 = 23. \nonumber
\end{align}
Under separation the good types burn $DW$ each:
\begin{align}
  23-\alpha DW &= 23-(0.3)(3.774) = 23-1.132 = 21.868, \nonumber\\
  \text{check: } \alpha U^{sep}+(1-\alpha)(qR-I) &= (0.3)(26.226)+(0.7)(20) = 7.868+14 = 21.868. \nonumber
\end{align}
The bad type gets $q(R-I/m)=22.055$ under pooling and $qR-I=20$ under separation.

\answer{(f)} Surplus $23$ (pooling) vs.\ $21.868$ (separation): $1.132$ is lost. The bad type falls from $22.055$ to $20$, the good type rises from $25.205$ to $26.226$. Pooling is efficient (both projects have positive NPV and are financed, nothing is destroyed); the good type still signals because she gains $1.021$ by ending her cross-subsidy, while the cost falls on the bad type and on the collateral destroyed.

\intuition Signalling with collateral is a private fix for adverse selection, not a social one. The good type compares two ways of paying for the bad type's presence: the cross-subsidy in the pool, which grows with the share of bad types and with $p-q$, and the expected collateral loss $(1-p)(1-\delta)C_G$, which does not depend on $\alpha$. When bad types are many, the cross-subsidy is the bigger bill and she signals. Cheaper collateral (higher $\delta$) requires more of it but wastes less, so signalling becomes more attractive; a smaller gap between types requires less of it. When both types are worth financing, every unit of collateral destroyed is a social loss: pooling would be efficient.

\pitfalls
\begin{itemize}
  \item Letting $C_G$ depend on $\alpha$: in a separating equilibrium lenders know the type, so the prior plays no role in the contract; $\alpha$ enters only through the pooling alternative.
  \item Reading ``more collateral'' as ``more waste'': as $\delta$ rises, $C_G$ rises but $DW=(1-p)(1-\delta)C_G$ falls.
  \item Concluding that separation is efficient because the good type prefers it: total surplus is lower by $\alpha(1-p)(1-\delta)C_G$, paid for by the bad type's lost cross-subsidy.
\end{itemize}

\refbib{Lecture Topic~3, block 3.7, slides 56--67 (costly collateral pledging; slide 67 compares signalling with pooling). Tirole \S6.3, Application 5 ``Costly Collateral Pledging'', pp.~251--253 (equation (6.4) and the determinants of collateralization, $\partial C^*/\partial\beta>0$, $\partial C^*/\partial q<0$, p.~253). Original exercise: no book answer.}

\end{document}
